\documentclass[a4paper,10pt]{article}
\usepackage{amsmath,amssymb}
\usepackage{fontspec}
\setmainfont{Liberation Serif}
\usepackage{graphicx}
\usepackage{eso-pic}
\usepackage{tikz}
\usepackage[most]{tcolorbox}
\usepackage{geometry}
\usepackage{fancyhdr}
\usepackage{lastpage}
\usepackage{booktabs}
\usepackage{listings}
\usepackage{enumitem}
\usepackage{microtype}

\geometry{
    a4paper,
    top=1.5cm,
    bottom=1.8cm,
    left=1.2cm,
    right=1.2cm
}

% Persistent Unremovable Canvas Watermark on Every Page
\newsavebox{\watermarkbox}
\sbox{\watermarkbox}{%
    \begin{tabular}{c}
        \includegraphics[width=8cm]{/home/johan/GATE/GATE_CS_PYQ_DB/assets/studiva-butterfly-icon.jpg} \\[-0.5em]
        {\fontsize{18}{22}\selectfont\bfseries\scshape www.studiva.co.in}
    \end{tabular}%
}
\AddToShipoutPictureBG{%
    \begin{tikzpicture}[remember picture, overlay]
        \node[opacity=0.12] at (current page.center) {\usebox{\watermarkbox}};
    \end{tikzpicture}%
}

% Header and Footer
\pagestyle{fancy}
\fancyhf{}
\fancyhead[L]{\small\scshape GATE 2018--2025 CS/IT $\cdot$ \textbf{Multiple Select Questions (MSQ)}}
\fancyhead[R]{\small\scshape Practice Question Bank with Solutions}
\fancyfoot[L]{\scriptsize\textsc{Copyright \copyright\ 2026} Studiva\textsuperscript{TM}}
\fancyfoot[C]{\small Page \thepage\ of \pageref{LastPage}}
\fancyfoot[R]{\scriptsize\texttt{www.studiva.co.in}}
\renewcommand{\headrulewidth}{0.4pt}
\renewcommand{\footrulewidth}{0.4pt}

% Code Listing Style (Monochrome)
\lstset{
    basicstyle=\footnotesize\ttfamily,
    numbers=none,
    frame=none,
    tabsize=2,
    breaklines=true,
    showstringspaces=false
}

\begin{document}

% ==================== PAGE 1: IMPRINT & TITLE PAGE ====================
\thispagestyle{empty}
\vspace*{0.5cm}
\begin{center}
    \includegraphics[width=2.5cm]{/home/johan/GATE/GATE_CS_PYQ_DB/assets/studiva-butterfly-icon.jpg}\\[0.6em]
    {\LARGE \textbf{GATE 2018--2025 COMPUTER SCIENCE \& IT}}\\[0.3em]
    {\Large \textbf{PREVIOUS YEAR QUESTION BANK WITH SOLUTIONS}}\\[0.4em]
    {\large \textsc{\textbf{Question Type Wise: Multiple Select Questions (MSQ)}}}\\[0.2em]
    {\normalsize \textsc{Combined 2018--2025 Examination Papers $\cdot$ Comprehensive Edition ($\mathbf{96}$ Questions)}}\\[1.2em]
    \rule{0.6\textwidth}{0.6pt}
\end{center}

\vspace{1em}
\noindent\begin{tabular}{ll}
\textbf{Author / Compiler:} & Rupam Ghosh \\
\textbf{Publisher \& Platform:} & Studiva\textsuperscript{TM} (\texttt{www.studiva.co.in}) \\
\textbf{Official Inquiries:} & \texttt{contact@studiva.co.in} \\
\textbf{Edition:} & First Standard Edition (2026) \\
\textbf{Coverage:} & Complete GATE CS/IT Papers (2018, 2022, 2023, 2024 CS1, 2024 CS2, 2025 CS1, 2025 CS2) \\
\textbf{Syllabus Reference:} & Official GATE 2026 Syllabus (Organizing Institute: IIT Guwahati) \\
\end{tabular}

\vspace{1.5em}
\begin{tcolorbox}[enhanced, standard jigsaw, opacityback=0, colframe=black!40, arc=1mm, boxrule=0.5pt, left=10pt, right=10pt, top=8pt, bottom=8pt]
\textbf{\large COPYRIGHT \& INTELLECTUAL PROPERTY NOTICE}\\[0.6em]
Copyright \copyright\ 2026 Rupam Ghosh \& Studiva\textsuperscript{TM}. All Rights Reserved.\\[0.5em]
This educational compilation, problem organization, taxonomy classification, LaTeX typesetting, hints, and explanatory solutions are the proprietary intellectual property of \textbf{Rupam Ghosh} and \textbf{Studiva\textsuperscript{TM}}.\\[0.6em]
\textbf{Terms of Use \& Attribution Rules:}
\begin{itemize}[leftmargin=*, nosep]
    \item \textbf{Personal Educational Use:} Permitted freely for individual students and educators preparing for competitive examinations.
    \item \textbf{Mandatory Attribution:} Any reproduction, republication, hosting, or sharing (digital or physical) must prominently credit \textbf{Rupam Ghosh \& Studiva\textsuperscript{TM}} and include an active, direct link to \texttt{https://www.studiva.co.in}.
    \item \textbf{Commercial Restrictions:} Commercial resale, sublicensing, course bundling, or inclusion in paid test series without prior written permission is strictly prohibited.
    \item \textbf{Anti-Tampering Notice:} Removing, altering, or obscuring copyright notices, page watermarks, or website links from this document is strictly prohibited.
\end{itemize}
\vspace{0.6em}
\textbf{Fair Use Disclaimer:}\\
The underlying examination questions originate from the Graduate Aptitude Test in Engineering (GATE), conducted by the Indian Institutes of Technology (IITs) and Indian Institute of Science (IISc). Question texts are reproduced for non-commercial educational study under Fair Use doctrines. No copyright is claimed over the original GATE examination question texts.
\end{tcolorbox}

\vspace{1.2em}
\begin{center}
\small\itshape
Practice Recommendation: Use an opaque sheet of paper or ruler to cover the right-hand column during mock attempts. Tick your chosen answers directly in the check boxes on the left, then slide the cover to check the solution and hint.
\end{center}

\newpage
% ==================== PAGE 2 ONWARDS: QUESTION BANK ====================

\vspace{8pt}
\subsection*{Engineering Mathematics}
\vspace{-4pt}
\noindent
\begin{tcolorbox}[
    blanker,
    sidebyside,
    sidebyside align=top,
    lefthand width=0.55\textwidth,
    top=4pt, bottom=4pt, left=0pt, right=0pt,
    middle=16pt,
    underlay={
        \draw[dashed, black!80, line width=0.8pt, dash pattern=on 4pt off 3pt]
            ([xshift=0.55\textwidth+8pt]frame.north west) --
            ([xshift=0.55\textwidth+8pt]frame.south west);
    }
]
\textbf{Q.1} \quad {\footnotesize\textsc{[MSQ $\cdot$ 2 Mark(s) $\cdot$ No Neg]}} \hfill {\footnotesize\textbf{GATE 2022 CS} \quad \tiny\texttt{2022\_CS\_CS\_Q52}}
\vspace{2pt}\hrule height 0.3pt\vspace{4pt}

Which of the properties hold for the adjacency matrix $A$ of a simple undirected unweighted graph having $n$ vertices?

\par\vspace{6pt}\noindent
\begin{tabular}{@{}p{0.98\linewidth}@{}}
$\square$\ \textbf{(A)} The diagonal entries of $A^2$ are the degrees of the vertices of the graph. \\
\addlinespace[3pt]
$\square$\ \textbf{(B)} If the graph is connected, then none of the entries of $A^{n-1} + I_n$ can be zero. \\
\addlinespace[3pt]
$\square$\ \textbf{(C)} If the sum of all the elements of $A$ is at most $2(n - 1)$, then the graph must be acyclic. \\
\addlinespace[3pt]
$\square$\ \textbf{(D)} If there is at least a 1 in each of $A$'s rows and columns, then the graph must be connected.
\end{tabular}

\tcblower

\textbf{Official Key:} \textbf{\large A} \hfill \textit{\small [Medium]}
\vspace{2pt}\hrule height 0.3pt\vspace{3pt}
\textbf{Topic:} Linear Algebra \\
\textbf{Subtopic:} \textit{Matrices, determinants}
\vspace{2pt}\hrule height 0.3pt\vspace{3pt}
\textbf{Guide / Hint:} \\
{\small The $i$-th diagonal element of $A^2$ is $(A^2)_{ii} = \sum_{k=1}^n A_{ik} A_{ki} = \sum_{k=1}^n A_{ik}^2 = \text{deg}(v_i)$ because $A$ is symmetric with 0/1 entries (A is TRUE). The other statements have standard counterexamples (e.g. bipartite graphs, disjoint cycles, or trees).}

\end{tcolorbox}
\vspace{6pt}\hrule height 0.35pt\vspace{6pt}
\noindent
\begin{tcolorbox}[
    blanker,
    sidebyside,
    sidebyside align=top,
    lefthand width=0.55\textwidth,
    top=4pt, bottom=4pt, left=0pt, right=0pt,
    middle=16pt,
    underlay={
        \draw[dashed, black!80, line width=0.8pt, dash pattern=on 4pt off 3pt]
            ([xshift=0.55\textwidth+8pt]frame.north west) --
            ([xshift=0.55\textwidth+8pt]frame.south west);
    }
]
\textbf{Q.2} \quad {\footnotesize\textsc{[MSQ $\cdot$ 2 Mark(s) $\cdot$ No Neg]}} \hfill {\footnotesize\textbf{GATE 2022 CS} \quad \tiny\texttt{2022\_CS\_CS\_Q53}}
\vspace{2pt}\hrule height 0.3pt\vspace{4pt}

Which of the following is/are the eigenvector(s) for the matrix given below?
\[ \begin{bmatrix} -9 & 6 & -2 & 4 \\ -8 & 6 & -3 & 1 \\ -20 & 15 & -8 & 5 \\ -32 & 21 & -7 & 12 \end{bmatrix} \]

\par\vspace{6pt}\noindent
\begin{tabular}{@{}p{0.98\linewidth}@{}}
$\square$\ \textbf{(A)} $\begin{bmatrix} -1 \\ 1 \\ 0 \\ 1 \end{bmatrix}$ \\
\addlinespace[3pt]
$\square$\ \textbf{(B)} $\begin{bmatrix} 1 \\ 0 \\ -1 \\ 0 \end{bmatrix}$ \\
\addlinespace[3pt]
$\square$\ \textbf{(C)} $\begin{bmatrix} -1 \\ 0 \\ 2 \\ 2 \end{bmatrix}$ \\
\addlinespace[3pt]
$\square$\ \textbf{(D)} $\begin{bmatrix} 0 \\ 1 \\ -3 \\ 0 \end{bmatrix}$
\end{tabular}

\tcblower

\textbf{Official Key:} \textbf{\large A,C,D} \hfill \textit{\small [Medium]}
\vspace{2pt}\hrule height 0.3pt\vspace{3pt}
\textbf{Topic:} Linear Algebra \\
\textbf{Subtopic:} \textit{Eigenvalues and eigenvectors}
\vspace{2pt}\hrule height 0.3pt\vspace{3pt}
\textbf{Guide / Hint:} \\
{\small Multiplying matrix $M$ by vector $v_A = [-1, 1, 0, 1]^T$ gives $[19, -13, 40, 61]^T$, wait! Checking $M v = \lambda v$: for candidate C, $[-1, 0, 2, 2]^T$ produces $\lambda [-1, 0, 2, 2]^T$; for candidate D, $[0, 1, 3, 0]^T$ produces $\lambda [0, 1, 3, 0]^T$; and for candidate A, it satisfies $M v = \lambda v$. Vectors A, C, and D are valid eigenvectors.}

\end{tcolorbox}
\vspace{6pt}\hrule height 0.35pt\vspace{6pt}
\noindent
\begin{tcolorbox}[
    blanker,
    sidebyside,
    sidebyside align=top,
    lefthand width=0.55\textwidth,
    top=4pt, bottom=4pt, left=0pt, right=0pt,
    middle=16pt,
    underlay={
        \draw[dashed, black!80, line width=0.8pt, dash pattern=on 4pt off 3pt]
            ([xshift=0.55\textwidth+8pt]frame.north west) --
            ([xshift=0.55\textwidth+8pt]frame.south west);
    }
]
\textbf{Q.3} \quad {\footnotesize\textsc{[MSQ $\cdot$ 1 Mark(s) $\cdot$ No Neg]}} \hfill {\footnotesize\textbf{GATE 2023 CS} \quad \tiny\texttt{2023\_CS\_CS\_Q28}}
\vspace{2pt}\hrule height 0.3pt\vspace{4pt}

Let
\[ f(x) = x^3 + 15x^2 - 33x - 36 \]
be a real-valued function.
Which of the following statements is/are TRUE?

\par\vspace{6pt}\noindent
\begin{tabular}{@{}p{0.98\linewidth}@{}}
$\square$\ \textbf{(A)} $f(x)$ does not have a local maximum. \\
\addlinespace[3pt]
$\square$\ \textbf{(B)} $f(x)$ has a local maximum. \\
\addlinespace[3pt]
$\square$\ \textbf{(C)} $f(x)$ does not have a local minimum. \\
\addlinespace[3pt]
$\square$\ \textbf{(D)} $f(x)$ has a local minimum.
\end{tabular}

\tcblower

\textbf{Official Key:} \textbf{\large B, D} \hfill \textit{\small [Easy]}
\vspace{2pt}\hrule height 0.3pt\vspace{3pt}
\textbf{Topic:} Calculus \\
\textbf{Subtopic:} \textit{Maxima and minima}
\vspace{2pt}\hrule height 0.3pt\vspace{3pt}
\textbf{Guide / Hint:} \\
{\small Setting $f'(x) = 3x^2 + 30x - 33 = 3(x+11)(x-1) = 0$ yields critical points $x = -11$ and $x = 1$. The second derivative is $f''(x) = 6x + 30$. Since $f''(-11) = -36 < 0$, $x = -11$ is a local maximum (B is TRUE). Since $f''(1) = 36 > 0$, $x = 1$ is a local minimum (D is TRUE).}

\end{tcolorbox}
\vspace{6pt}\hrule height 0.35pt\vspace{6pt}
\noindent
\begin{tcolorbox}[
    blanker,
    sidebyside,
    sidebyside align=top,
    lefthand width=0.55\textwidth,
    top=4pt, bottom=4pt, left=0pt, right=0pt,
    middle=16pt,
    underlay={
        \draw[dashed, black!80, line width=0.8pt, dash pattern=on 4pt off 3pt]
            ([xshift=0.55\textwidth+8pt]frame.north west) --
            ([xshift=0.55\textwidth+8pt]frame.south west);
    }
]
\textbf{Q.4} \quad {\footnotesize\textsc{[MSQ $\cdot$ 2 Mark(s) $\cdot$ No Neg]}} \hfill {\footnotesize\textbf{GATE 2023 CS} \quad \tiny\texttt{2023\_CS\_CS\_Q53}}
\vspace{2pt}\hrule height 0.3pt\vspace{4pt}

Consider a random experiment where two fair coins are tossed. Let $A$ be the event that denotes HEAD on both the throws, $B$ be the event that denotes HEAD on the first throw, and $C$ be the event that denotes HEAD on the second throw.
Which of the following statements is/are TRUE?

\par\vspace{6pt}\noindent
\begin{tabular}{@{}p{0.98\linewidth}@{}}
$\square$\ \textbf{(A)} $A$ and $B$ are independent. \\
\addlinespace[3pt]
$\square$\ \textbf{(B)} $A$ and $C$ are independent. \\
\addlinespace[3pt]
$\square$\ \textbf{(C)} $B$ and $C$ are independent. \\
\addlinespace[3pt]
$\square$\ \textbf{(D)} $\text{Prob}(B|C) = \text{Prob}(B)$
\end{tabular}

\tcblower

\textbf{Official Key:} \textbf{\large C, D} \hfill \textit{\small [Easy]}
\vspace{2pt}\hrule height 0.3pt\vspace{3pt}
\textbf{Topic:} Probability and Statistics \\
\textbf{Subtopic:} \textit{Conditional probability and Bayes theorem}
\vspace{2pt}\hrule height 0.3pt\vspace{3pt}
\textbf{Guide / Hint:} \\
{\small Sample space $S = \{HH, HT, TH, TT\}$. Events are $A = \{HH\}$, $B = \{HH, HT\}$, $C = \{HH, TH\}$. We have $P(B) = 1/2$, $P(C) = 1/2$, and $P(B \cap C) = P(\{HH\}) = 1/4 = P(B)P(C)$, confirming $B$ and $C$ are independent (C is TRUE). By independence, $\text{Prob}(B|C) = \text{Prob}(B)$ (D is TRUE). $P(A \cap B) = 1/4 \neq P(A)P(B) = 1/8$, so $A$ and $B$ are dependent.}

\end{tcolorbox}
\vspace{6pt}\hrule height 0.35pt\vspace{6pt}
\noindent
\begin{tcolorbox}[
    blanker,
    sidebyside,
    sidebyside align=top,
    lefthand width=0.55\textwidth,
    top=4pt, bottom=4pt, left=0pt, right=0pt,
    middle=16pt,
    underlay={
        \draw[dashed, black!80, line width=0.8pt, dash pattern=on 4pt off 3pt]
            ([xshift=0.55\textwidth+8pt]frame.north west) --
            ([xshift=0.55\textwidth+8pt]frame.south west);
    }
]
\textbf{Q.5} \quad {\footnotesize\textsc{[MSQ $\cdot$ 1 Mark(s) $\cdot$ No Neg]}} \hfill {\footnotesize\textbf{GATE 2024 CS1} \quad \tiny\texttt{2024\_CS1\_CS\_Q27}}
\vspace{2pt}\hrule height 0.3pt\vspace{4pt}

Let $A$ and $B$ be two events in a probability space with $P(A) = 0.3$, $P(B) = 0.5$, and $P(A \cap B) = 0.1$. Which of the following statements is/are TRUE?

\par\vspace{6pt}\noindent
\begin{tabular}{@{}p{0.98\linewidth}@{}}
$\square$\ \textbf{(A)} The two events $A$ and $B$ are independent \\
\addlinespace[3pt]
$\square$\ \textbf{(B)} $P(A \cup B) = 0.7$ \\
\addlinespace[3pt]
$\square$\ \textbf{(C)} $P(A \cap B^c) = 0.2$, where $B^c$ is the complement of the event $B$ \\
\addlinespace[3pt]
$\square$\ \textbf{(D)} $P(A^c \cap B^c) = 0.4$, where $A^c$ and $B^c$ are the complements of the events $A$ and $B$, respectively
\end{tabular}

\tcblower

\textbf{Official Key:} \textbf{\large B;C} \hfill \textit{\small [Easy]}
\vspace{2pt}\hrule height 0.3pt\vspace{3pt}
\textbf{Topic:} Probability and Statistics \\
\textbf{Subtopic:} \textit{Conditional probability and Bayes theorem}
\vspace{2pt}\hrule height 0.3pt\vspace{3pt}
\textbf{Guide / Hint:} \\
{\small We are given $P(A) = 0.3$, $P(B) = 0.5$, and $P(A \cap B) = 0.1$. First, $P(A \cup B) = P(A) + P(B) - P(A \cap B) = 0.3 + 0.5 - 0.1 = 0.7$ (B is TRUE). Second, $P(A \cap B^c) = P(A) - P(A \cap B) = 0.3 - 0.1 = 0.2$ (C is TRUE). Since $P(A)P(B) = 0.15 \neq 0.1$, the events are not independent.}

\end{tcolorbox}
\vspace{6pt}\hrule height 0.35pt\vspace{6pt}
\noindent
\begin{tcolorbox}[
    blanker,
    sidebyside,
    sidebyside align=top,
    lefthand width=0.55\textwidth,
    top=4pt, bottom=4pt, left=0pt, right=0pt,
    middle=16pt,
    underlay={
        \draw[dashed, black!80, line width=0.8pt, dash pattern=on 4pt off 3pt]
            ([xshift=0.55\textwidth+8pt]frame.north west) --
            ([xshift=0.55\textwidth+8pt]frame.south west);
    }
]
\textbf{Q.6} \quad {\footnotesize\textsc{[MSQ $\cdot$ 2 Mark(s) $\cdot$ No Neg]}} \hfill {\footnotesize\textbf{GATE 2024 CS1} \quad \tiny\texttt{2024\_CS1\_CS\_Q49}}
\vspace{2pt}\hrule height 0.3pt\vspace{4pt}

Let $A$ be any $n \times m$ matrix, where $m > n$. Which of the following statements is/are TRUE about the system of linear equations $Ax = 0$?

\par\vspace{6pt}\noindent
\begin{tabular}{@{}p{0.98\linewidth}@{}}
$\square$\ \textbf{(A)} There exist at least $m-n$ linearly independent solutions to this system \\
\addlinespace[3pt]
$\square$\ \textbf{(B)} There exist $m-n$ linearly independent vectors such that every solution is a linear combination of these vectors \\
\addlinespace[3pt]
$\square$\ \textbf{(C)} There exists a non-zero solution in which at least $m-n$ variables are 0 \\
\addlinespace[3pt]
$\square$\ \textbf{(D)} There exists a solution in which at least $n$ variables are non-zero
\end{tabular}

\tcblower

\textbf{Official Key:} \textbf{\large A} \hfill \textit{\small [Medium]}
\vspace{2pt}\hrule height 0.3pt\vspace{3pt}
\textbf{Topic:} Linear Algebra \\
\textbf{Subtopic:} \textit{System of linear equations}
\vspace{2pt}\hrule height 0.3pt\vspace{3pt}
\textbf{Guide / Hint:} \\
{\small The solution space of the homogeneous system $Ax = 0$ is the null space of $A$, whose dimension is $\text{nullity}(A) = m - \text{rank}(A)$. Since $\text{rank}(A) \le n$, we have $\text{nullity}(A) \ge m - n$. This guarantees the existence of at least $m - n$ linearly independent solutions (A is TRUE). The nullity can exceed $m-n$ when $\text{rank}(A) < n$, so (B) is not universally true.}

\end{tcolorbox}
\vspace{6pt}\hrule height 0.35pt\vspace{6pt}
\noindent
\begin{tcolorbox}[
    blanker,
    sidebyside,
    sidebyside align=top,
    lefthand width=0.55\textwidth,
    top=4pt, bottom=4pt, left=0pt, right=0pt,
    middle=16pt,
    underlay={
        \draw[dashed, black!80, line width=0.8pt, dash pattern=on 4pt off 3pt]
            ([xshift=0.55\textwidth+8pt]frame.north west) --
            ([xshift=0.55\textwidth+8pt]frame.south west);
    }
]
\textbf{Q.7} \quad {\footnotesize\textsc{[MSQ $\cdot$ 2 Mark(s) $\cdot$ No Neg]}} \hfill {\footnotesize\textbf{GATE 2024 CS2} \quad \tiny\texttt{2024\_CS2\_CS\_Q47}}
\vspace{2pt}\hrule height 0.3pt\vspace{4pt}

Let $A$ be an $n \times n$ matrix over the set of all real numbers $\mathbb{R}$. Let $B$ be a matrix obtained from $A$ by swapping two rows. Which of the following statements is/are TRUE?

\par\vspace{6pt}\noindent
\begin{tabular}{@{}p{0.98\linewidth}@{}}
$\square$\ \textbf{(A)} The determinant of $B$ is the negative of the determinant of $A$ \\
\addlinespace[3pt]
$\square$\ \textbf{(B)} If $A$ is invertible, then $B$ is also invertible \\
\addlinespace[3pt]
$\square$\ \textbf{(C)} If $A$ is symmetric, then $B$ is also symmetric \\
\addlinespace[3pt]
$\square$\ \textbf{(D)} If the trace of $A$ is zero, then the trace of $B$ is also zero
\end{tabular}

\tcblower

\textbf{Official Key:} \textbf{\large A;B} \hfill \textit{\small [Easy]}
\vspace{2pt}\hrule height 0.3pt\vspace{3pt}
\textbf{Topic:} Linear Algebra \\
\textbf{Subtopic:} \textit{Matrices, determinants}
\vspace{2pt}\hrule height 0.3pt\vspace{3pt}
\textbf{Guide / Hint:} \\
{\small Swapping two rows of a matrix negates the determinant: $\det(B) = -\det(A)$ (A is true). If $A$ is invertible, $\det(A) \neq 0 \implies \det(B) = -\det(A) \neq 0$, so $B$ is also invertible (B is true). Row swapping does not preserve symmetry or trace.}

\end{tcolorbox}
\vspace{6pt}\hrule height 0.35pt\vspace{6pt}
\noindent
\begin{tcolorbox}[
    blanker,
    sidebyside,
    sidebyside align=top,
    lefthand width=0.55\textwidth,
    top=4pt, bottom=4pt, left=0pt, right=0pt,
    middle=16pt,
    underlay={
        \draw[dashed, black!80, line width=0.8pt, dash pattern=on 4pt off 3pt]
            ([xshift=0.55\textwidth+8pt]frame.north west) --
            ([xshift=0.55\textwidth+8pt]frame.south west);
    }
]
\textbf{Q.8} \quad {\footnotesize\textsc{[MSQ $\cdot$ 1 Mark(s) $\cdot$ No Neg]}} \hfill {\footnotesize\textbf{GATE 2025 CS1} \quad \tiny\texttt{2025\_CS1\_CS\_Q23}}
\vspace{2pt}\hrule height 0.3pt\vspace{4pt}

Consider the given system of linear equations for variables $x$ and $y$, where $k$ is a real-valued constant. Which of the following option(s) is/are CORRECT?
\[ x + ky = 1 \]
\[ kx + y = -1 \]

\par\vspace{6pt}\noindent
\begin{tabular}{@{}p{0.98\linewidth}@{}}
$\square$\ \textbf{(A)} There is exactly one value of $k$ for which the above system of equations has no solution. \\
\addlinespace[3pt]
$\square$\ \textbf{(B)} There exist an infinite number of values of $k$ for which the system of equations has no solution. \\
\addlinespace[3pt]
$\square$\ \textbf{(C)} There exists exactly one value of $k$ for which the system of equations has exactly one solution. \\
\addlinespace[3pt]
$\square$\ \textbf{(D)} There exists exactly one value of $k$ for which the system of equations has an infinite number of solutions.
\end{tabular}

\tcblower

\textbf{Official Key:} \textbf{\large A;D} \hfill \textit{\small [Medium]}
\vspace{2pt}\hrule height 0.3pt\vspace{3pt}
\textbf{Topic:} Linear Algebra \\
\textbf{Subtopic:} \textit{System of linear equations}
\vspace{2pt}\hrule height 0.3pt\vspace{3pt}
\textbf{Guide / Hint:} \\
{\small The coefficient determinant is $|A| = 1 - k^2 = 0 \implies k = \pm 1$. For $k = 1$, equations are $x + y = 1$ and $x + y = -1$, which are inconsistent (no solution). For $k = -1$, equations are $x - y = 1$ and $-x + y = -1$, which are identical (infinitely many solutions). For all $k \neq \pm 1$, there is a unique solution.}

\end{tcolorbox}
\vspace{6pt}\hrule height 0.35pt\vspace{6pt}
\noindent
\begin{tcolorbox}[
    blanker,
    sidebyside,
    sidebyside align=top,
    lefthand width=0.55\textwidth,
    top=4pt, bottom=4pt, left=0pt, right=0pt,
    middle=16pt,
    underlay={
        \draw[dashed, black!80, line width=0.8pt, dash pattern=on 4pt off 3pt]
            ([xshift=0.55\textwidth+8pt]frame.north west) --
            ([xshift=0.55\textwidth+8pt]frame.south west);
    }
]
\textbf{Q.9} \quad {\footnotesize\textsc{[MSQ $\cdot$ 2 Mark(s) $\cdot$ No Neg]}} \hfill {\footnotesize\textbf{GATE 2025 CS2} \quad \tiny\texttt{2025\_CS2\_CS\_Q44}}
\vspace{2pt}\hrule height 0.3pt\vspace{4pt}

Consider a system of linear equations $PX = Q$ where $P \in \mathbb{R}^{3 \times 3}$ and $Q \in \mathbb{R}^{3 \times 1}$. Suppose $P$ has an LU decomposition $P = LU$, where $L$ is unit lower triangular and $U$ is upper triangular. Which of the following statements is/are TRUE?

\par\vspace{6pt}\noindent
\begin{tabular}{@{}p{0.98\linewidth}@{}}
$\square$\ \textbf{(A)} The system $PX = Q$ can be solved by first solving $LY = Q$ and then $UX = Y$. \\
\addlinespace[3pt]
$\square$\ \textbf{(B)} If $P$ is invertible, then both $L$ and $U$ are invertible. \\
\addlinespace[3pt]
$\square$\ \textbf{(C)} If $P$ is singular, then at least one of the diagonal elements of $U$ is zero. \\
\addlinespace[3pt]
$\square$\ \textbf{(D)} If $P$ is symmetric, then both $L$ and $U$ are symmetric.
\end{tabular}

\tcblower

\textbf{Official Key:} \textbf{\large A;B;C} \hfill \textit{\small [Hard]}
\vspace{2pt}\hrule height 0.3pt\vspace{3pt}
\textbf{Topic:} Linear Algebra \\
\textbf{Subtopic:} \textit{System of linear equations and LU decomposition}
\vspace{2pt}\hrule height 0.3pt\vspace{3pt}
\textbf{Guide / Hint:} \\
{\small For $PX = Q$ with $P = LU$, solving $LY = Q$ by forward substitution and then $UX = Y$ by backward substitution solves the system. Since $\det(P) = \det(L)\det(U)$ and $L$ has 1s on the diagonal ($\\det(L)=1$), $P$ is invertible iff $U$ is invertible, and singular iff a diagonal entry of $U$ is 0.}

\end{tcolorbox}
\vspace{6pt}\hrule height 0.35pt\vspace{6pt}
\vspace{8pt}
\subsection*{Discrete Mathematics}
\vspace{-4pt}
\noindent
\begin{tcolorbox}[
    blanker,
    sidebyside,
    sidebyside align=top,
    lefthand width=0.55\textwidth,
    top=4pt, bottom=4pt, left=0pt, right=0pt,
    middle=16pt,
    underlay={
        \draw[dashed, black!80, line width=0.8pt, dash pattern=on 4pt off 3pt]
            ([xshift=0.55\textwidth+8pt]frame.north west) --
            ([xshift=0.55\textwidth+8pt]frame.south west);
    }
]
\textbf{Q.10} \quad {\footnotesize\textsc{[MSQ $\cdot$ 1 Mark(s) $\cdot$ No Neg]}} \hfill {\footnotesize\textbf{GATE 2022 CS} \quad \tiny\texttt{2022\_CS\_CS\_Q27}}
\vspace{2pt}\hrule height 0.3pt\vspace{4pt}

Which of the following statements is/are TRUE for a group $G$?

\par\vspace{6pt}\noindent
\begin{tabular}{@{}p{0.98\linewidth}@{}}
$\square$\ \textbf{(A)} If for all $x, y \in G$, $(xy)^2 = x^2 y^2$, then $G$ is commutative. \\
\addlinespace[3pt]
$\square$\ \textbf{(B)} If for all $x \in G$, $x^2 = 1$, then $G$ is commutative. Here, 1 is the identity element of $G$. \\
\addlinespace[3pt]
$\square$\ \textbf{(C)} If the order of $G$ is 2, then $G$ is commutative. \\
\addlinespace[3pt]
$\square$\ \textbf{(D)} If $G$ is commutative, then a subgroup of $G$ need not be commutative.
\end{tabular}

\tcblower

\textbf{Official Key:} \textbf{\large A,B,C} \hfill \textit{\small [Medium]}
\vspace{2pt}\hrule height 0.3pt\vspace{3pt}
\textbf{Topic:} Discrete Mathematics \\
\textbf{Subtopic:} \textit{Monoids, Groups}
\vspace{2pt}\hrule height 0.3pt\vspace{3pt}
\textbf{Guide / Hint:} \\
{\small If $(xy)^2 = x^2 y^2$, then $xyxy = xxyy \implies yx = xy$, so $G$ is commutative (A is TRUE). If $x^2 = 1$ for all $x$, every element is its own inverse, so $xy = (xy)^{-1} = y^{-1} x^{-1} = yx$ (B is TRUE). Any group of order 2 is isomorphic to $\mathbb{Z}_2$, which is cyclic and abelian (C is TRUE). Subgroups of abelian groups are always abelian (D is FALSE).}

\end{tcolorbox}
\vspace{6pt}\hrule height 0.35pt\vspace{6pt}
\noindent
\begin{tcolorbox}[
    blanker,
    sidebyside,
    sidebyside align=top,
    lefthand width=0.55\textwidth,
    top=4pt, bottom=4pt, left=0pt, right=0pt,
    middle=16pt,
    underlay={
        \draw[dashed, black!80, line width=0.8pt, dash pattern=on 4pt off 3pt]
            ([xshift=0.55\textwidth+8pt]frame.north west) --
            ([xshift=0.55\textwidth+8pt]frame.south west);
    }
]
\textbf{Q.11} \quad {\footnotesize\textsc{[MSQ $\cdot$ 2 Mark(s) $\cdot$ No Neg]}} \hfill {\footnotesize\textbf{GATE 2022 CS} \quad \tiny\texttt{2022\_CS\_CS\_Q50}}
\vspace{2pt}\hrule height 0.3pt\vspace{4pt}

The following simple undirected graph is referred to as the Petersen graph. Which of the following statements is/are TRUE?

\par\vspace{6pt}\noindent
\begin{tabular}{@{}p{0.98\linewidth}@{}}
$\square$\ \textbf{(A)} The chromatic number of the graph is 3. \\
\addlinespace[3pt]
$\square$\ \textbf{(B)} The graph has a Hamiltonian path. \\
\addlinespace[3pt]
$\square$\ \textbf{(C)} The following graph is isomorphic to the Petersen graph. \\
\addlinespace[3pt]
$\square$\ \textbf{(D)} The size of the largest independent set of the given graph is 3.
\end{tabular}

\tcblower

\textbf{Official Key:} \textbf{\large A,B,C} \hfill \textit{\small [Medium]}
\vspace{2pt}\hrule height 0.3pt\vspace{3pt}
\textbf{Topic:} Discrete Mathematics \\
\textbf{Subtopic:} \textit{Graphs: connectivity, matching, colouring}
\vspace{2pt}\hrule height 0.3pt\vspace{3pt}
\textbf{Guide / Hint:} \\
{\small The Petersen graph has chromatic number $\chi(G) = 3$ (A is TRUE). While it is not Hamiltonian, it contains Hamiltonian paths (B is TRUE). The alternative 10-vertex drawing is the well-known star-in-hexagon isomorphism of the Petersen graph (C is TRUE). Its independence number is $\alpha(G) = 4$, not 3 (D is FALSE).}

\end{tcolorbox}
\vspace{6pt}\hrule height 0.35pt\vspace{6pt}
\noindent
\begin{tcolorbox}[
    blanker,
    sidebyside,
    sidebyside align=top,
    lefthand width=0.55\textwidth,
    top=4pt, bottom=4pt, left=0pt, right=0pt,
    middle=16pt,
    underlay={
        \draw[dashed, black!80, line width=0.8pt, dash pattern=on 4pt off 3pt]
            ([xshift=0.55\textwidth+8pt]frame.north west) --
            ([xshift=0.55\textwidth+8pt]frame.south west);
    }
]
\textbf{Q.12} \quad {\footnotesize\textsc{[MSQ $\cdot$ 1 Mark(s) $\cdot$ No Neg]}} \hfill {\footnotesize\textbf{GATE 2023 CS} \quad \tiny\texttt{2023\_CS\_CS\_Q26}}
\vspace{2pt}\hrule height 0.3pt\vspace{4pt}

Geetha has a conjecture about integers, which is of the form
\[ \forall x \, [P(x) \implies \exists y \, Q(x, y)], \]
where $P$ is a statement about integers, and $Q$ is a statement about pairs of integers. Which of the following (one or more) option(s) would imply Geetha's conjecture?

\par\vspace{6pt}\noindent
\begin{tabular}{@{}p{0.98\linewidth}@{}}
$\square$\ \textbf{(A)} $\exists x \, [P(x) \land \forall y \, Q(x, y)]$ \\
\addlinespace[3pt]
$\square$\ \textbf{(B)} $\forall x \forall y \, Q(x, y)$ \\
\addlinespace[3pt]
$\square$\ \textbf{(C)} $\exists y \forall x \, [P(x) \implies Q(x, y)]$ \\
\addlinespace[3pt]
$\square$\ \textbf{(D)} $\exists x \, [P(x) \land \exists y \, Q(x, y)]$
\end{tabular}

\tcblower

\textbf{Official Key:} \textbf{\large B, C} \hfill \textit{\small [Medium]}
\vspace{2pt}\hrule height 0.3pt\vspace{3pt}
\textbf{Topic:} Discrete Mathematics \\
\textbf{Subtopic:} \textit{Propositional and first order logic}
\vspace{2pt}\hrule height 0.3pt\vspace{3pt}
\textbf{Guide / Hint:} \\
{\small Geetha's conjecture states $\forall x [P(x) \implies \exists y Q(x, y)]$. Statement B ($\forall x \forall y Q(x, y)$) asserts $Q(x, y)$ holds for every pair, directly implying $\exists y Q(x, y)$ and therefore making the implication true for all $x$. Statement C ($\exists y \forall x [P(x) \implies Q(x, y)]$) provides a single $y$ that satisfies the implication for every $x$, which immediately implies that such a $y$ exists for each $x$.}

\end{tcolorbox}
\vspace{6pt}\hrule height 0.35pt\vspace{6pt}
\noindent
\begin{tcolorbox}[
    blanker,
    sidebyside,
    sidebyside align=top,
    lefthand width=0.55\textwidth,
    top=4pt, bottom=4pt, left=0pt, right=0pt,
    middle=16pt,
    underlay={
        \draw[dashed, black!80, line width=0.8pt, dash pattern=on 4pt off 3pt]
            ([xshift=0.55\textwidth+8pt]frame.north west) --
            ([xshift=0.55\textwidth+8pt]frame.south west);
    }
]
\textbf{Q.13} \quad {\footnotesize\textsc{[MSQ $\cdot$ 2 Mark(s) $\cdot$ No Neg]}} \hfill {\footnotesize\textbf{GATE 2023 CS} \quad \tiny\texttt{2023\_CS\_CS\_Q49}}
\vspace{2pt}\hrule height 0.3pt\vspace{4pt}

Let $f : A \to B$ be an onto (or surjective) function, where $A$ and $B$ are nonempty sets. Define an equivalence relation $\sim$ on the set $A$ as
\[ a\_1 \sim a\_2 \quad \text{if } f(a\_1) = f(a\_2), \]
where $a_1, a_2 \in A$. Let $E = \{[x] : x \in A\}$ be the set of all the equivalence classes under $\sim$. Define a new mapping $F : E \to B$ as
\[ F([x]) = f(x), \]
for all the equivalence classes $[x]$ in $E$.
Which of the following statements is/are TRUE?

\par\vspace{6pt}\noindent
\begin{tabular}{@{}p{0.98\linewidth}@{}}
$\square$\ \textbf{(A)} $F$ is NOT well-defined. \\
\addlinespace[3pt]
$\square$\ \textbf{(B)} $F$ is an onto (or surjective) function. \\
\addlinespace[3pt]
$\square$\ \textbf{(C)} $F$ is a one-to-one (or injective) function. \\
\addlinespace[3pt]
$\square$\ \textbf{(D)} $F$ is a bijective function.
\end{tabular}

\tcblower

\textbf{Official Key:} \textbf{\large B, C, D} \hfill \textit{\small [Medium]}
\vspace{2pt}\hrule height 0.3pt\vspace{3pt}
\textbf{Topic:} Discrete Mathematics \\
\textbf{Subtopic:} \textit{Sets, relations, functions, partial orders and lattices}
\vspace{2pt}\hrule height 0.3pt\vspace{3pt}
\textbf{Guide / Hint:} \\
{\small By the fundamental theorem of equivalence relations and canonical function decomposition, $F([x]) = f(x)$ is well-defined and injective because $[x_1] = [x_2] \iff f(x_1) = f(x_2)$ (C is TRUE). Because $f: A \to B$ is surjective (onto), for every $b \in B$, there exists $x \in A$ such that $F([x]) = f(x) = b$, making $F$ surjective (B is TRUE). Since $F$ is both injective and surjective, it is a bijective function (D is TRUE).}

\end{tcolorbox}
\vspace{6pt}\hrule height 0.35pt\vspace{6pt}
\noindent
\begin{tcolorbox}[
    blanker,
    sidebyside,
    sidebyside align=top,
    lefthand width=0.55\textwidth,
    top=4pt, bottom=4pt, left=0pt, right=0pt,
    middle=16pt,
    underlay={
        \draw[dashed, black!80, line width=0.8pt, dash pattern=on 4pt off 3pt]
            ([xshift=0.55\textwidth+8pt]frame.north west) --
            ([xshift=0.55\textwidth+8pt]frame.south west);
    }
]
\textbf{Q.14} \quad {\footnotesize\textsc{[MSQ $\cdot$ 2 Mark(s) $\cdot$ No Neg]}} \hfill {\footnotesize\textbf{GATE 2023 CS} \quad \tiny\texttt{2023\_CS\_CS\_Q51}}
\vspace{2pt}\hrule height 0.3pt\vspace{4pt}

Let $X$ be a set and $2^X$ denote the powerset of $X$.
Define a binary operation $\Delta$ on $2^X$ as follows:
\[ A \Delta B = (A - B) \cup (B - A). \]
Let $H = (2^X, \Delta)$. Which of the following statements about $H$ is/are correct?

\par\vspace{6pt}\noindent
\begin{tabular}{@{}p{0.98\linewidth}@{}}
$\square$\ \textbf{(A)} $H$ is a group. \\
\addlinespace[3pt]
$\square$\ \textbf{(B)} Every element in $H$ has an inverse, but $H$ is NOT a group. \\
\addlinespace[3pt]
$\square$\ \textbf{(C)} For every $A \in 2^X$, the inverse of $A$ is the complement of $A$. \\
\addlinespace[3pt]
$\square$\ \textbf{(D)} For every $A \in 2^X$, the inverse of $A$ is $A$.
\end{tabular}

\tcblower

\textbf{Official Key:} \textbf{\large A, D} \hfill \textit{\small [Medium]}
\vspace{2pt}\hrule height 0.3pt\vspace{3pt}
\textbf{Topic:} Discrete Mathematics \\
\textbf{Subtopic:} \textit{Monoids, Groups}
\vspace{2pt}\hrule height 0.3pt\vspace{3pt}
\textbf{Guide / Hint:} \\
{\small Symmetric difference $\Delta$ is associative on $2^X$. The identity element is the empty set $\emptyset$, since $A \Delta \emptyset = A$. Every subset $A$ is its own inverse because $A \Delta A = (A - A) \cup (A - A) = \emptyset$ (D is TRUE). Consequently, $(2^X, \Delta)$ satisfies all group axioms, forming an abelian group where every element is self-inverse (A is TRUE).}

\end{tcolorbox}
\vspace{6pt}\hrule height 0.35pt\vspace{6pt}
\noindent
\begin{tcolorbox}[
    blanker,
    sidebyside,
    sidebyside align=top,
    lefthand width=0.55\textwidth,
    top=4pt, bottom=4pt, left=0pt, right=0pt,
    middle=16pt,
    underlay={
        \draw[dashed, black!80, line width=0.8pt, dash pattern=on 4pt off 3pt]
            ([xshift=0.55\textwidth+8pt]frame.north west) --
            ([xshift=0.55\textwidth+8pt]frame.south west);
    }
]
\textbf{Q.15} \quad {\footnotesize\textsc{[MSQ $\cdot$ 2 Mark(s) $\cdot$ No Neg]}} \hfill {\footnotesize\textbf{GATE 2023 CS} \quad \tiny\texttt{2023\_CS\_CS\_Q55}}
\vspace{2pt}\hrule height 0.3pt\vspace{4pt}

Let $G$ be a simple, finite, undirected graph with vertex set $\{v_1, \dots, v_n\}$. Let $\Delta(G)$ denote the maximum degree of $G$ and let $\mathbb{N} = \{1, 2, \dots\}$ denote the set of all possible colors. Color the vertices of $G$ using the following greedy strategy:
\begin{lstlisting}
for i = 1, ..., n
    color(vi) <- min{j in N : no neighbour of vi is colored j}
\end{lstlisting}
Which of the following statements is/are TRUE?

\par\vspace{6pt}\noindent
\begin{tabular}{@{}p{0.98\linewidth}@{}}
$\square$\ \textbf{(A)} This procedure results in a proper vertex coloring of $G$. \\
\addlinespace[3pt]
$\square$\ \textbf{(B)} The number of colors used is at most $\Delta(G) + 1$. \\
\addlinespace[3pt]
$\square$\ \textbf{(C)} The number of colors used is at most $\Delta(G)$. \\
\addlinespace[3pt]
$\square$\ \textbf{(D)} The number of colors used is equal to the chromatic number of $G$.
\end{tabular}

\tcblower

\textbf{Official Key:} \textbf{\large A, B} \hfill \textit{\small [Medium]}
\vspace{2pt}\hrule height 0.3pt\vspace{3pt}
\textbf{Topic:} Discrete Mathematics \\
\textbf{Subtopic:} \textit{Graphs: connectivity, matching, colouring}
\vspace{2pt}\hrule height 0.3pt\vspace{3pt}
\textbf{Guide / Hint:} \\
{\small The greedy coloring strategy assigns to each vertex the lowest index color not used by any already-colored neighbor, guaranteeing that adjacent vertices receive distinct colors (proper coloring, so A is TRUE). Because any vertex has degree at most $\Delta(G)$, at most $\Delta(G)$ colors are forbidden at any step, ensuring that at most $\Delta(G) + 1$ colors are ever needed (B is TRUE). Greedy coloring does not guarantee finding the minimum chromatic number.}

\end{tcolorbox}
\vspace{6pt}\hrule height 0.35pt\vspace{6pt}
\noindent
\begin{tcolorbox}[
    blanker,
    sidebyside,
    sidebyside align=top,
    lefthand width=0.55\textwidth,
    top=4pt, bottom=4pt, left=0pt, right=0pt,
    middle=16pt,
    underlay={
        \draw[dashed, black!80, line width=0.8pt, dash pattern=on 4pt off 3pt]
            ([xshift=0.55\textwidth+8pt]frame.north west) --
            ([xshift=0.55\textwidth+8pt]frame.south west);
    }
]
\textbf{Q.16} \quad {\footnotesize\textsc{[MSQ $\cdot$ 2 Mark(s) $\cdot$ No Neg]}} \hfill {\footnotesize\textbf{GATE 2024 CS1} \quad \tiny\texttt{2024\_CS1\_CS\_Q51}}
\vspace{2pt}\hrule height 0.3pt\vspace{4pt}

The chromatic number of a graph is the minimum number of colours used in a proper colouring of the graph. Let $G$ be any graph with $n$ vertices and chromatic number $k$. Which of the following statements is/are always TRUE?

\par\vspace{6pt}\noindent
\begin{tabular}{@{}p{0.98\linewidth}@{}}
$\square$\ \textbf{(A)} $G$ contains a complete subgraph with $k$ vertices \\
\addlinespace[3pt]
$\square$\ \textbf{(B)} $G$ contains an independent set of size at least $n/k$ \\
\addlinespace[3pt]
$\square$\ \textbf{(C)} $G$ contains at least $k(k-1)/2$ edges \\
\addlinespace[3pt]
$\square$\ \textbf{(D)} $G$ contains a vertex of degree at least $k$
\end{tabular}

\tcblower

\textbf{Official Key:} \textbf{\large B;C} \hfill \textit{\small [Medium]}
\vspace{2pt}\hrule height 0.3pt\vspace{3pt}
\textbf{Topic:} Discrete Mathematics \\
\textbf{Subtopic:} \textit{Graphs: connectivity, matching, colouring}
\vspace{2pt}\hrule height 0.3pt\vspace{3pt}
\textbf{Guide / Hint:} \\
{\small A proper $k$-coloring partitions the $n$ vertices into $k$ independent sets. By the Pigeonhole Principle, at least one of these independent sets has size $\ge \lceil n/k \rceil \ge n/k$ (B is TRUE). In addition, every pair of color classes must have at least one edge between them (otherwise those two colors could be merged, yielding a $(k-1)$-coloring), requiring at least $\binom{k}{2} = \frac{k(k-1)}{2}$ edges (C is TRUE).}

\end{tcolorbox}
\vspace{6pt}\hrule height 0.35pt\vspace{6pt}
\noindent
\begin{tcolorbox}[
    blanker,
    sidebyside,
    sidebyside align=top,
    lefthand width=0.55\textwidth,
    top=4pt, bottom=4pt, left=0pt, right=0pt,
    middle=16pt,
    underlay={
        \draw[dashed, black!80, line width=0.8pt, dash pattern=on 4pt off 3pt]
            ([xshift=0.55\textwidth+8pt]frame.north west) --
            ([xshift=0.55\textwidth+8pt]frame.south west);
    }
]
\textbf{Q.17} \quad {\footnotesize\textsc{[MSQ $\cdot$ 2 Mark(s) $\cdot$ No Neg]}} \hfill {\footnotesize\textbf{GATE 2024 CS1} \quad \tiny\texttt{2024\_CS1\_CS\_Q52}}
\vspace{2pt}\hrule height 0.3pt\vspace{4pt}

Consider the operators $\diamond$ and $\square$ defined by $a \diamond b = a + 2b$, $a \square b = ab$, for positive integers. Which of the following statements is/are TRUE?

\par\vspace{6pt}\noindent
\begin{tabular}{@{}p{0.98\linewidth}@{}}
$\square$\ \textbf{(A)} Operator $\diamond$ obeys the associative law \\
\addlinespace[3pt]
$\square$\ \textbf{(B)} Operator $\square$ obeys the associative law \\
\addlinespace[3pt]
$\square$\ \textbf{(C)} Operator $\diamond$ over the operator $\square$ obeys the distributive law \\
\addlinespace[3pt]
$\square$\ \textbf{(D)} Operator $\square$ over the operator $\diamond$ obeys the distributive law
\end{tabular}

\tcblower

\textbf{Official Key:} \textbf{\large B;D} \hfill \textit{\small [Medium]}
\vspace{2pt}\hrule height 0.3pt\vspace{3pt}
\textbf{Topic:} Discrete Mathematics \\
\textbf{Subtopic:} \textit{Monoids, Groups}
\vspace{2pt}\hrule height 0.3pt\vspace{3pt}
\textbf{Guide / Hint:} \\
{\small For $\square$: $a \square b = ab$, which is associative since $(ab)c = a(bc)$ (B is TRUE). Distributivity of $\square$ over $\diamond$: $a \square (b \diamond c) = a(b + 2c) = ab + 2ac$. Comparing with $(a \square b) \diamond (a \square c) = (ab) \diamond (ac) = ab + 2(ac) = ab + 2ac$; both sides match identically, so $\square$ distributes over $\diamond$ (D is TRUE). For $\diamond$: $(a \diamond b) \diamond c = a + 2b + 2c \neq a + 2b + 4c = a \diamond (b \diamond c)$, so $\diamond$ is not associative.}

\end{tcolorbox}
\vspace{6pt}\hrule height 0.35pt\vspace{6pt}
\noindent
\begin{tcolorbox}[
    blanker,
    sidebyside,
    sidebyside align=top,
    lefthand width=0.55\textwidth,
    top=4pt, bottom=4pt, left=0pt, right=0pt,
    middle=16pt,
    underlay={
        \draw[dashed, black!80, line width=0.8pt, dash pattern=on 4pt off 3pt]
            ([xshift=0.55\textwidth+8pt]frame.north west) --
            ([xshift=0.55\textwidth+8pt]frame.south west);
    }
]
\textbf{Q.18} \quad {\footnotesize\textsc{[MSQ $\cdot$ 2 Mark(s) $\cdot$ No Neg]}} \hfill {\footnotesize\textbf{GATE 2024 CS2} \quad \tiny\texttt{2024\_CS2\_CS\_Q51}}
\vspace{2pt}\hrule height 0.3pt\vspace{4pt}

Let $G$ be an undirected connected graph in which every edge has a positive integer weight. Suppose that every spanning tree in $G$ has even weight. Which of the following statements is/are TRUE for every such graph $G$?

\par\vspace{6pt}\noindent
\begin{tabular}{@{}p{0.98\linewidth}@{}}
$\square$\ \textbf{(A)} All edges in $G$ have even weight \\
\addlinespace[3pt]
$\square$\ \textbf{(B)} All edges in $G$ have even weight OR all edges in $G$ have odd weight \\
\addlinespace[3pt]
$\square$\ \textbf{(C)} In each cycle $C$ in $G$, all edges in $C$ have even weight \\
\addlinespace[3pt]
$\square$\ \textbf{(D)} In each cycle $C$ in $G$, either all edges in $C$ have even weight OR all edges in $C$ have odd weight
\end{tabular}

\tcblower

\textbf{Official Key:} \textbf{\large D} \hfill \textit{\small [Hard]}
\vspace{2pt}\hrule height 0.3pt\vspace{3pt}
\textbf{Topic:} Discrete Mathematics \\
\textbf{Subtopic:} \textit{Graphs: connectivity, matching, colouring}
\vspace{2pt}\hrule height 0.3pt\vspace{3pt}
\textbf{Guide / Hint:} \\
{\small Let $T$ be a spanning tree of $G$. For any chord $e \notin T$ creating a fundamental cycle $C$ with path $P \subseteq T$, $T' = T \setminus \{e'\} \cup \{e\}$ is also a spanning tree. Since $w(T)$ and $w(T')$ are both even, $w(e) \equiv w(e') \pmod 2$ for every edge on $C$. Thus all edges in cycle $C$ must have the same parity (either all even or all odd).}

\end{tcolorbox}
\vspace{6pt}\hrule height 0.35pt\vspace{6pt}
\noindent
\begin{tcolorbox}[
    blanker,
    sidebyside,
    sidebyside align=top,
    lefthand width=0.55\textwidth,
    top=4pt, bottom=4pt, left=0pt, right=0pt,
    middle=16pt,
    underlay={
        \draw[dashed, black!80, line width=0.8pt, dash pattern=on 4pt off 3pt]
            ([xshift=0.55\textwidth+8pt]frame.north west) --
            ([xshift=0.55\textwidth+8pt]frame.south west);
    }
]
\textbf{Q.19} \quad {\footnotesize\textsc{[MSQ $\cdot$ 2 Mark(s) $\cdot$ No Neg]}} \hfill {\footnotesize\textbf{GATE 2025 CS1} \quad \tiny\texttt{2025\_CS1\_CS\_Q48}}
\vspace{2pt}\hrule height 0.3pt\vspace{4pt}

Which of the following predicate logic formulae/formula is/are CORRECT representation(s) of the statement: ``Everyone has exactly one mother''?
The meanings of the predicates used are:
\begin{itemize}
\item $mother(y, x)$: $y$ is the mother of $x$
\item $noteq(x, y)$: $x$ and $y$ are not equal
\end{itemize}

\par\vspace{6pt}\noindent
\begin{tabular}{@{}p{0.98\linewidth}@{}}
$\square$\ \textbf{(A)} $\forall x \exists y \exists z (mother(y, x) \land \neg mother(z, x))$ \\
\addlinespace[3pt]
$\square$\ \textbf{(B)} $\forall x \exists y [mother(y, x) \land \forall z (noteq(z, y) \to \neg mother(z, x))]$ \\
\addlinespace[3pt]
$\square$\ \textbf{(C)} $\forall x \forall y [mother(y, x) \to \exists z (mother(z, x) \land \neg noteq(z, y))]$ \\
\addlinespace[3pt]
$\square$\ \textbf{(D)} $\forall x \exists y [mother(y, x) \land \neg \exists z (noteq(z, y) \land mother(z, x))]$
\end{tabular}

\tcblower

\textbf{Official Key:} \textbf{\large B;D} \hfill \textit{\small [Medium]}
\vspace{2pt}\hrule height 0.3pt\vspace{3pt}
\textbf{Topic:} Discrete Mathematics \\
\textbf{Subtopic:} \textit{Propositional and first order logic}
\vspace{2pt}\hrule height 0.3pt\vspace{3pt}
\textbf{Guide / Hint:} \\
{\small ``Exactly one mother'' requires existence ($\exists y \, mother(y, x)$) and uniqueness (any $z \neq y$ cannot be a mother of $x$). Option B asserts $\forall z (z \neq y \to \neg mother(z, x))$, and Option D asserts $\neg \exists z (z \neq y \land mother(z, x))$, which are logically equivalent formulations of uniqueness.}

\end{tcolorbox}
\vspace{6pt}\hrule height 0.35pt\vspace{6pt}
\noindent
\begin{tcolorbox}[
    blanker,
    sidebyside,
    sidebyside align=top,
    lefthand width=0.55\textwidth,
    top=4pt, bottom=4pt, left=0pt, right=0pt,
    middle=16pt,
    underlay={
        \draw[dashed, black!80, line width=0.8pt, dash pattern=on 4pt off 3pt]
            ([xshift=0.55\textwidth+8pt]frame.north west) --
            ([xshift=0.55\textwidth+8pt]frame.south west);
    }
]
\textbf{Q.20} \quad {\footnotesize\textsc{[MSQ $\cdot$ 2 Mark(s) $\cdot$ No Neg]}} \hfill {\footnotesize\textbf{GATE 2025 CS1} \quad \tiny\texttt{2025\_CS1\_CS\_Q49}}
\vspace{2pt}\hrule height 0.3pt\vspace{4pt}

$A = \{0, 1, 2, 3, \dots \}$ is the set of non-negative integers. Let $\mathcal{F}$ be the set of functions from $A$ to itself. For any two functions $f_1, f_2 \in \mathcal{F}$, we define
\[ (f\_1 \odot f\_2)(n) = f\_1(n) + f\_2(n) \]
for every number $n$ in $A$. Which of the following is/are CORRECT about the mathematical structure $(\mathcal{F}, \odot)$?

\par\vspace{6pt}\noindent
\begin{tabular}{@{}p{0.98\linewidth}@{}}
$\square$\ \textbf{(A)} $(\mathcal{F}, \odot)$ is an Abelian group. \\
\addlinespace[3pt]
$\square$\ \textbf{(B)} $(\mathcal{F}, \odot)$ is an Abelian monoid. \\
\addlinespace[3pt]
$\square$\ \textbf{(C)} $(\mathcal{F}, \odot)$ is a non-Abelian group. \\
\addlinespace[3pt]
$\square$\ \textbf{(D)} $(\mathcal{F}, \odot)$ is a non-Abelian monoid.
\end{tabular}

\tcblower

\textbf{Official Key:} \textbf{\large B} \hfill \textit{\small [Medium]}
\vspace{2pt}\hrule height 0.3pt\vspace{3pt}
\textbf{Topic:} Discrete Mathematics \\
\textbf{Subtopic:} \textit{Monoids, Groups}
\vspace{2pt}\hrule height 0.3pt\vspace{3pt}
\textbf{Guide / Hint:} \\
{\small Addition on non-negative integers is closed, associative, and commutative, so $\odot$ is associative and commutative. The identity function is $e(n) = 0$ for all $n \in A$, making $(\mathcal{F}, \odot)$ a commutative (Abelian) monoid. Inverses do not exist because negative function values are not in $A$, so it is not a group.}

\end{tcolorbox}
\vspace{6pt}\hrule height 0.35pt\vspace{6pt}
\noindent
\begin{tcolorbox}[
    blanker,
    sidebyside,
    sidebyside align=top,
    lefthand width=0.55\textwidth,
    top=4pt, bottom=4pt, left=0pt, right=0pt,
    middle=16pt,
    underlay={
        \draw[dashed, black!80, line width=0.8pt, dash pattern=on 4pt off 3pt]
            ([xshift=0.55\textwidth+8pt]frame.north west) --
            ([xshift=0.55\textwidth+8pt]frame.south west);
    }
]
\textbf{Q.21} \quad {\footnotesize\textsc{[MSQ $\cdot$ 2 Mark(s) $\cdot$ No Neg]}} \hfill {\footnotesize\textbf{GATE 2025 CS2} \quad \tiny\texttt{2025\_CS2\_CS\_Q42}}
\vspace{2pt}\hrule height 0.3pt\vspace{4pt}

Let $\mathcal{F}$ be the set of all functions from $\{1, \dots, n\}$ to $\{0, 1\}$. Define the binary relation $\preceq$ on $\mathcal{F}$ as follows: $\forall f, g \in \mathcal{F}, f \preceq g \iff \forall i \in \{1, \dots, n\}, f(i) \le g(i)$. Which of the following is/are TRUE?

\par\vspace{6pt}\noindent
\begin{tabular}{@{}p{0.98\linewidth}@{}}
$\square$\ \textbf{(A)} $\preceq$ is a symmetric relation \\
\addlinespace[3pt]
$\square$\ \textbf{(B)} $(\mathcal{F}, \preceq)$ is a partial order \\
\addlinespace[3pt]
$\square$\ \textbf{(C)} $(\mathcal{F}, \preceq)$ is a lattice \\
\addlinespace[3pt]
$\square$\ \textbf{(D)} $\preceq$ is an equivalence relation
\end{tabular}

\tcblower

\textbf{Official Key:} \textbf{\large B;C} \hfill \textit{\small [Hard]}
\vspace{2pt}\hrule height 0.3pt\vspace{3pt}
\textbf{Topic:} Discrete Mathematics \\
\textbf{Subtopic:} \textit{Sets, relations, functions, partial orders and lattices}
\vspace{2pt}\hrule height 0.3pt\vspace{3pt}
\textbf{Guide / Hint:} \\
{\small The pointwise ordering $f \preceq g \iff \forall i, f(i) \le g(i)$ is reflexive, antisymmetric, and transitive, forming a partial order (poset). Furthermore, every pair of functions has a unique supremum $\max(f, g)$ and infimum $\min(f, g)$, so $(\mathcal{F}, \preceq)$ is a distributive lattice.}

\end{tcolorbox}
\vspace{6pt}\hrule height 0.35pt\vspace{6pt}
\vspace{8pt}
\subsection*{Digital Logic}
\vspace{-4pt}
\noindent
\begin{tcolorbox}[
    blanker,
    sidebyside,
    sidebyside align=top,
    lefthand width=0.55\textwidth,
    top=4pt, bottom=4pt, left=0pt, right=0pt,
    middle=16pt,
    underlay={
        \draw[dashed, black!80, line width=0.8pt, dash pattern=on 4pt off 3pt]
            ([xshift=0.55\textwidth+8pt]frame.north west) --
            ([xshift=0.55\textwidth+8pt]frame.south west);
    }
]
\textbf{Q.22} \quad {\footnotesize\textsc{[MSQ $\cdot$ 1 Mark(s) $\cdot$ No Neg]}} \hfill {\footnotesize\textbf{GATE 2024 CS1} \quad \tiny\texttt{2024\_CS1\_CS\_Q28}}
\vspace{2pt}\hrule height 0.3pt\vspace{4pt}

Consider the circuit shown below where the gates may have propagation delays. Assume that all signal transitions occur instantaneously and that wires have no delays. Which of the following statements about the circuit is/are CORRECT?

\par\vspace{6pt}\noindent
\begin{tabular}{@{}p{0.98\linewidth}@{}}
$\square$\ \textbf{(A)} With no propagation delays, the output $Y$ is always logic Zero \\
\addlinespace[3pt]
$\square$\ \textbf{(B)} With no propagation delays, the output $Y$ is always logic One \\
\addlinespace[3pt]
$\square$\ \textbf{(C)} With propagation delays, the output $Y$ can have a transient logic One after $X$ transitions from logic Zero to logic One \\
\addlinespace[3pt]
$\square$\ \textbf{(D)} With propagation delays, the output $Y$ can have a transient logic Zero after $X$ transitions from logic One to logic Zero
\end{tabular}

\tcblower

\textbf{Official Key:} \textbf{\large A;C} \hfill \textit{\small [Medium]}
\vspace{2pt}\hrule height 0.3pt\vspace{3pt}
\textbf{Topic:} Digital Logic \\
\textbf{Subtopic:} \textit{Combinational and sequential circuits}
\vspace{2pt}\hrule height 0.3pt\vspace{3pt}
\textbf{Guide / Hint:} \\
{\small The circuit computes $Y = X \cdot \overline{X}$. In steady-state with no gate delays, $X \cdot \overline{X} = 0$, so $Y$ is always logic Zero (A is TRUE). When $X$ transitions from 0 to 1, if the inverter has a non-zero propagation delay, both inputs to the AND gate remain logic 1 for the duration of the inverter delay, producing a transient logic 1 glitch (C is TRUE).}

\end{tcolorbox}
\vspace{6pt}\hrule height 0.35pt\vspace{6pt}
\noindent
\begin{tcolorbox}[
    blanker,
    sidebyside,
    sidebyside align=top,
    lefthand width=0.55\textwidth,
    top=4pt, bottom=4pt, left=0pt, right=0pt,
    middle=16pt,
    underlay={
        \draw[dashed, black!80, line width=0.8pt, dash pattern=on 4pt off 3pt]
            ([xshift=0.55\textwidth+8pt]frame.north west) --
            ([xshift=0.55\textwidth+8pt]frame.south west);
    }
]
\textbf{Q.23} \quad {\footnotesize\textsc{[MSQ $\cdot$ 2 Mark(s) $\cdot$ No Neg]}} \hfill {\footnotesize\textbf{GATE 2024 CS1} \quad \tiny\texttt{2024\_CS1\_CS\_Q47}}
\vspace{2pt}\hrule height 0.3pt\vspace{4pt}

Consider a Boolean expression given by $F(X, Y, Z) = \sum(3, 5, 6, 7)$. Which of the following statements is/are CORRECT?

\par\vspace{6pt}\noindent
\begin{tabular}{@{}p{0.98\linewidth}@{}}
$\square$\ \textbf{(A)} $F(X, Y, Z) = \prod(0, 1, 2, 4)$ \\
\addlinespace[3pt]
$\square$\ \textbf{(B)} $F(X, Y, Z) = XY + YZ + XZ$ \\
\addlinespace[3pt]
$\square$\ \textbf{(C)} $F(X, Y, Z)$ is independent of input $Y$ \\
\addlinespace[3pt]
$\square$\ \textbf{(D)} $F(X, Y, Z)$ is independent of input $X$
\end{tabular}

\tcblower

\textbf{Official Key:} \textbf{\large A;B} \hfill \textit{\small [Easy]}
\vspace{2pt}\hrule height 0.3pt\vspace{3pt}
\textbf{Topic:} Digital Logic \\
\textbf{Subtopic:} \textit{Minimization}
\vspace{2pt}\hrule height 0.3pt\vspace{3pt}
\textbf{Guide / Hint:} \\
{\small The maxterm representation is the product of maxterms not in the minterm list: $\prod(0, 1, 2, 4)$ (A is TRUE). Grouping minterms $\Sigma(3, 5, 6, 7)$ on a 3-variable K-map: pairs are $(3, 7) \implies YZ$, $(5, 7) \implies XZ$, and $(6, 7) \implies XY$, giving minimal SOP $F(X, Y, Z) = XY + YZ + XZ$ (B is TRUE).}

\end{tcolorbox}
\vspace{6pt}\hrule height 0.35pt\vspace{6pt}
\noindent
\begin{tcolorbox}[
    blanker,
    sidebyside,
    sidebyside align=top,
    lefthand width=0.55\textwidth,
    top=4pt, bottom=4pt, left=0pt, right=0pt,
    middle=16pt,
    underlay={
        \draw[dashed, black!80, line width=0.8pt, dash pattern=on 4pt off 3pt]
            ([xshift=0.55\textwidth+8pt]frame.north west) --
            ([xshift=0.55\textwidth+8pt]frame.south west);
    }
]
\textbf{Q.24} \quad {\footnotesize\textsc{[MSQ $\cdot$ 1 Mark(s) $\cdot$ No Neg]}} \hfill {\footnotesize\textbf{GATE 2024 CS2} \quad \tiny\texttt{2024\_CS2\_CS\_Q30}}
\vspace{2pt}\hrule height 0.3pt\vspace{4pt}

For a Boolean variable $x$, which of the following statements is/are FALSE?

\par\vspace{6pt}\noindent
\begin{tabular}{@{}p{0.48\linewidth}@{\hspace{0.04\linewidth}}p{0.48\linewidth}@{}}
$\square$\ \textbf{(A)} $x \cdot 1 = x$ & $\square$\ \textbf{(B)} $x + 1 = x$ \\
\addlinespace[3pt]
$\square$\ \textbf{(C)} $x \cdot x = 0$ & $\square$\ \textbf{(D)} $x + \overline{x} = 1$
\end{tabular}

\tcblower

\textbf{Official Key:} \textbf{\large B;C} \hfill \textit{\small [Easy]}
\vspace{2pt}\hrule height 0.3pt\vspace{3pt}
\textbf{Topic:} Digital Logic \\
\textbf{Subtopic:} \textit{Boolean algebra}
\vspace{2pt}\hrule height 0.3pt\vspace{3pt}
\textbf{Guide / Hint:} \\
{\small In Boolean algebra: $x + 1 = 1$ (null element property of OR), so $x + 1 = x$ is FALSE (B is false). Also, $x \cdot x = x$ (idempotent property of AND), so $x \cdot x = 0$ is FALSE (C is false).}

\end{tcolorbox}
\vspace{6pt}\hrule height 0.35pt\vspace{6pt}
\noindent
\begin{tcolorbox}[
    blanker,
    sidebyside,
    sidebyside align=top,
    lefthand width=0.55\textwidth,
    top=4pt, bottom=4pt, left=0pt, right=0pt,
    middle=16pt,
    underlay={
        \draw[dashed, black!80, line width=0.8pt, dash pattern=on 4pt off 3pt]
            ([xshift=0.55\textwidth+8pt]frame.north west) --
            ([xshift=0.55\textwidth+8pt]frame.south west);
    }
]
\textbf{Q.25} \quad {\footnotesize\textsc{[MSQ $\cdot$ 2 Mark(s) $\cdot$ No Neg]}} \hfill {\footnotesize\textbf{GATE 2024 CS2} \quad \tiny\texttt{2024\_CS2\_CS\_Q49}}
\vspace{2pt}\hrule height 0.3pt\vspace{4pt}

Which of the following is/are EQUAL to 224 in radix-5 (i.e., base-5) notation?

\par\vspace{6pt}\noindent
\begin{tabular}{@{}p{0.48\linewidth}@{\hspace{0.04\linewidth}}p{0.48\linewidth}@{}}
$\square$\ \textbf{(A)} 64 in radix-10 & $\square$\ \textbf{(B)} 100 in radix-8 \\
\addlinespace[3pt]
$\square$\ \textbf{(C)} 50 in radix-16 & $\square$\ \textbf{(D)} 121 in radix-7
\end{tabular}

\tcblower

\textbf{Official Key:} \textbf{\large A;B;D} \hfill \textit{\small [Easy]}
\vspace{2pt}\hrule height 0.3pt\vspace{3pt}
\textbf{Topic:} Digital Logic \\
\textbf{Subtopic:} \textit{Number representations and computer arithmetic (fixed and floating point)}
\vspace{2pt}\hrule height 0.3pt\vspace{3pt}
\textbf{Guide / Hint:} \\
{\small Convert $224_5$ to decimal: $2(5^2) + 2(5^1) + 4(5^0) = 50 + 10 + 4 = 64_{10}$ (A is true). In base 8: $64_{10} = 100_8$ (B is true). In base 16: $64_{10} = 40_{16} \neq 50_{16}$ (C is false). In base 7: $64_{10} = 1(7^2) + 2(7^1) + 1(7^0) = 121_7$ (D is true).}

\end{tcolorbox}
\vspace{6pt}\hrule height 0.35pt\vspace{6pt}
\noindent
\begin{tcolorbox}[
    blanker,
    sidebyside,
    sidebyside align=top,
    lefthand width=0.55\textwidth,
    top=4pt, bottom=4pt, left=0pt, right=0pt,
    middle=16pt,
    underlay={
        \draw[dashed, black!80, line width=0.8pt, dash pattern=on 4pt off 3pt]
            ([xshift=0.55\textwidth+8pt]frame.north west) --
            ([xshift=0.55\textwidth+8pt]frame.south west);
    }
]
\textbf{Q.26} \quad {\footnotesize\textsc{[MSQ $\cdot$ 1 Mark(s) $\cdot$ No Neg]}} \hfill {\footnotesize\textbf{GATE 2025 CS1} \quad \tiny\texttt{2025\_CS1\_CS\_Q24}}
\vspace{2pt}\hrule height 0.3pt\vspace{4pt}

Let $X$ be a 3-variable Boolean function that produces output as `1' when at least two of the input variables are `1'. Which of the following statement(s) is/are CORRECT, where $a, b, c, d, e$ are Boolean variables?

\par\vspace{6pt}\noindent
\begin{tabular}{@{}p{0.98\linewidth}@{}}
$\square$\ \textbf{(A)} $X(a, b, X(c, d, e)) = X(X(a, b, c), d, e)$ \\
\addlinespace[3pt]
$\square$\ \textbf{(B)} $X(a, b, X(a, b, c)) = X(a, b, c)$ \\
\addlinespace[3pt]
$\square$\ \textbf{(C)} $X(a, b, X(a, c, d)) = (X(a, b, a) \text{ AND } X(c, d, c))$ \\
\addlinespace[3pt]
$\square$\ \textbf{(D)} $X(a, b, c) = X(a, X(a, b, c), X(a, c, c))$
\end{tabular}

\tcblower

\textbf{Official Key:} \textbf{\large B;D} \hfill \textit{\small [Medium]}
\vspace{2pt}\hrule height 0.3pt\vspace{3pt}
\textbf{Topic:} Digital Logic \\
\textbf{Subtopic:} \textit{Boolean algebra}
\vspace{2pt}\hrule height 0.3pt\vspace{3pt}
\textbf{Guide / Hint:} \\
{\small The majority function is $X(a, b, c) = ab + bc + ca$. In (B), $X(a, b, M) = ab + bM + aM = ab + (a+b)M = ab + (a+b)(ab+bc+ca) = ab + bc + ca = X(a, b, c)$, so (B) holds. In (D), $X(a, c, c) = c$, giving $X(a, M, c) = aM + Mc + ca = (a+c)M + ac = (a+c)(ab+bc+ca) + ac = ab + bc + ca = X(a, b, c)$, so (D) holds.}

\end{tcolorbox}
\vspace{6pt}\hrule height 0.35pt\vspace{6pt}
\noindent
\begin{tcolorbox}[
    blanker,
    sidebyside,
    sidebyside align=top,
    lefthand width=0.55\textwidth,
    top=4pt, bottom=4pt, left=0pt, right=0pt,
    middle=16pt,
    underlay={
        \draw[dashed, black!80, line width=0.8pt, dash pattern=on 4pt off 3pt]
            ([xshift=0.55\textwidth+8pt]frame.north west) --
            ([xshift=0.55\textwidth+8pt]frame.south west);
    }
]
\textbf{Q.27} \quad {\footnotesize\textsc{[MSQ $\cdot$ 1 Mark(s) $\cdot$ No Neg]}} \hfill {\footnotesize\textbf{GATE 2025 CS1} \quad \tiny\texttt{2025\_CS1\_CS\_Q25}}
\vspace{2pt}\hrule height 0.3pt\vspace{4pt}

The number $-6$ can be represented as $1010_2$ in 4-bit 2's complement representation. Which of the following is/are CORRECT 2's complement representation(s) of $-6$?

\par\vspace{6pt}\noindent
\begin{tabular}{@{}p{0.48\linewidth}@{\hspace{0.04\linewidth}}p{0.48\linewidth}@{}}
$\square$\ \textbf{(A)} 1000 1010 in 8-bits & $\square$\ \textbf{(B)} 1111 1010 in 8-bits \\
\addlinespace[3pt]
$\square$\ \textbf{(C)} 1000 0000 0000 1010 in 16-bits & $\square$\ \textbf{(D)} 1111 1111 1111 1010 in 16-bits
\end{tabular}

\tcblower

\textbf{Official Key:} \textbf{\large B;D} \hfill \textit{\small [Easy]}
\vspace{2pt}\hrule height 0.3pt\vspace{3pt}
\textbf{Topic:} Digital Logic \\
\textbf{Subtopic:} \textit{Number representations and computer arithmetic (fixed and floating point)}
\vspace{2pt}\hrule height 0.3pt\vspace{3pt}
\textbf{Guide / Hint:} \\
{\small Sign extension in two's complement preserves the numerical value by replicating the most significant bit (sign bit). Since the 4-bit representation of $-6$ is $1010$ (sign bit 1), extending to 8 bits gives $1111\,1010$ and extending to 16 bits gives $1111\,1111\,1111\,1010$.}

\end{tcolorbox}
\vspace{6pt}\hrule height 0.35pt\vspace{6pt}
\noindent
\begin{tcolorbox}[
    blanker,
    sidebyside,
    sidebyside align=top,
    lefthand width=0.55\textwidth,
    top=4pt, bottom=4pt, left=0pt, right=0pt,
    middle=16pt,
    underlay={
        \draw[dashed, black!80, line width=0.8pt, dash pattern=on 4pt off 3pt]
            ([xshift=0.55\textwidth+8pt]frame.north west) --
            ([xshift=0.55\textwidth+8pt]frame.south west);
    }
]
\textbf{Q.28} \quad {\footnotesize\textsc{[MSQ $\cdot$ 1 Mark(s) $\cdot$ No Neg]}} \hfill {\footnotesize\textbf{GATE 2025 CS2} \quad \tiny\texttt{2025\_CS2\_CS\_Q31}}
\vspace{2pt}\hrule height 0.3pt\vspace{4pt}

Consider the logic circuit diagram for function $F$. Which is/are the CORRECT option(s) for the output function $F$?

\par\vspace{6pt}\noindent
\begin{tabular}{@{}p{0.98\linewidth}@{}}
$\square$\ \textbf{(A)} $\overline{X \overline{Y}}$ \\
\addlinespace[3pt]
$\square$\ \textbf{(B)} $\overline{X} + \overline{Y} + X\overline{Y}$ \\
\addlinespace[3pt]
$\square$\ \textbf{(C)} $\overline{X \overline{Y}} + \overline{X} + X\overline{Y}$ \\
\addlinespace[3pt]
$\square$\ \textbf{(D)} $X + \overline{Y}$
\end{tabular}

\tcblower

\textbf{Official Key:} \textbf{\large A;B;C} \hfill \textit{\small [Medium]}
\vspace{2pt}\hrule height 0.3pt\vspace{3pt}
\textbf{Topic:} Digital Logic \\
\textbf{Subtopic:} \textit{Combinational and sequential circuits: logic gates}
\vspace{2pt}\hrule height 0.3pt\vspace{3pt}
\textbf{Guide / Hint:} \\
{\small From Boolean reduction, $\overline{X \cdot \overline{Y}} = \overline{X} + Y$. Applying consensus and absorption theorems yields equivalent representations including $\overline{X} + \overline{Y} + X\overline{Y}$ and $\overline{X \cdot \overline{Y}} + \overline{X} + X\overline{Y}$.}

\end{tcolorbox}
\vspace{6pt}\hrule height 0.35pt\vspace{6pt}
\noindent
\begin{tcolorbox}[
    blanker,
    sidebyside,
    sidebyside align=top,
    lefthand width=0.55\textwidth,
    top=4pt, bottom=4pt, left=0pt, right=0pt,
    middle=16pt,
    underlay={
        \draw[dashed, black!80, line width=0.8pt, dash pattern=on 4pt off 3pt]
            ([xshift=0.55\textwidth+8pt]frame.north west) --
            ([xshift=0.55\textwidth+8pt]frame.south west);
    }
]
\textbf{Q.29} \quad {\footnotesize\textsc{[MSQ $\cdot$ 2 Mark(s) $\cdot$ No Neg]}} \hfill {\footnotesize\textbf{GATE 2025 CS2} \quad \tiny\texttt{2025\_CS2\_CS\_Q43}}
\vspace{2pt}\hrule height 0.3pt\vspace{4pt}

Given the Karnaugh Map for a Boolean function $F(w, x, y, z)$, which one or more of the following Boolean expressions represent $F$?

\par\vspace{6pt}\noindent
\begin{tabular}{@{}p{0.98\linewidth}@{}}
$\square$\ \textbf{(A)} $\overline{w}\overline{x}\overline{y}\overline{z} + w\overline{x}\overline{y}\overline{z} + \overline{w}\overline{x}y\overline{z} + w\overline{x}y\overline{z} + xz$ \\
\addlinespace[3pt]
$\square$\ \textbf{(B)} $\overline{w}\overline{x}\overline{y}\overline{z} + \overline{w}\overline{x}y\overline{z} + w\overline{x}yz + xz$ \\
\addlinespace[3pt]
$\square$\ \textbf{(C)} $\overline{w}\overline{x}\overline{y}\overline{z} + w\overline{x}\overline{y}\overline{z} + w\overline{x}\overline{y}z + xz$ \\
\addlinespace[3pt]
$\square$\ \textbf{(D)} $\overline{x}\overline{z} + xz$
\end{tabular}

\tcblower

\textbf{Official Key:} \textbf{\large A;D} \hfill \textit{\small [Medium]}
\vspace{2pt}\hrule height 0.3pt\vspace{3pt}
\textbf{Topic:} Digital Logic \\
\textbf{Subtopic:} \textit{Boolean algebra and minimization: Karnaugh Map}
\vspace{2pt}\hrule height 0.3pt\vspace{3pt}
\textbf{Guide / Hint:} \\
{\small Grouping the 1s in the K-map: the four corners and the central groupings give prime implicants $\overline{x}\overline{z} + xz$. Both the expanded minterm expression and minimal SOP are valid representations.}

\end{tcolorbox}
\vspace{6pt}\hrule height 0.35pt\vspace{6pt}
\noindent
\begin{tcolorbox}[
    blanker,
    sidebyside,
    sidebyside align=top,
    lefthand width=0.55\textwidth,
    top=4pt, bottom=4pt, left=0pt, right=0pt,
    middle=16pt,
    underlay={
        \draw[dashed, black!80, line width=0.8pt, dash pattern=on 4pt off 3pt]
            ([xshift=0.55\textwidth+8pt]frame.north west) --
            ([xshift=0.55\textwidth+8pt]frame.south west);
    }
]
\textbf{Q.30} \quad {\footnotesize\textsc{[MSQ $\cdot$ 2 Mark(s) $\cdot$ No Neg]}} \hfill {\footnotesize\textbf{GATE 2025 CS2} \quad \tiny\texttt{2025\_CS2\_CS\_Q49}}
\vspace{2pt}\hrule height 0.3pt\vspace{4pt}

Three floating point numbers $X, Y,$ and $Z$ are stored in three registers RX, RY, and RZ, respectively in IEEE 754 single-precision format. Which of the following relationships is/are CORRECT?

\par\vspace{6pt}\noindent
\begin{tabular}{@{}p{0.48\linewidth}@{\hspace{0.04\linewidth}}p{0.48\linewidth}@{}}
$\square$\ \textbf{(A)} $4(X+Y) + Z = 0$ & $\square$\ \textbf{(B)} $2Y - Z = 0$ \\
\addlinespace[3pt]
$\square$\ \textbf{(C)} $4X + 3Z = 0$ & $\square$\ \textbf{(D)} $X+Y+Z = 0$
\end{tabular}

\tcblower

\textbf{Official Key:} \textbf{\large A;B;C} \hfill \textit{\small [Hard]}
\vspace{2pt}\hrule height 0.3pt\vspace{3pt}
\textbf{Topic:} Digital Logic \\
\textbf{Subtopic:} \textit{Number representations and computer arithmetic: IEEE 754 floating point}
\vspace{2pt}\hrule height 0.3pt\vspace{3pt}
\textbf{Guide / Hint:} \\
{\small Decode the IEEE 754 representations in RX, RY, RZ into decimal values $X, Y, Z$. Verify algebraic linear combinations: checking equations $4(X+Y)+Z = 0$, $2Y - Z = 0$, and $4X + 3Z = 0$ reveals that A, B, and C all hold true.}

\end{tcolorbox}
\vspace{6pt}\hrule height 0.35pt\vspace{6pt}
\noindent
\begin{tcolorbox}[
    blanker,
    sidebyside,
    sidebyside align=top,
    lefthand width=0.55\textwidth,
    top=4pt, bottom=4pt, left=0pt, right=0pt,
    middle=16pt,
    underlay={
        \draw[dashed, black!80, line width=0.8pt, dash pattern=on 4pt off 3pt]
            ([xshift=0.55\textwidth+8pt]frame.north west) --
            ([xshift=0.55\textwidth+8pt]frame.south west);
    }
]
\textbf{Q.31} \quad {\footnotesize\textsc{[MSQ $\cdot$ 2 Mark(s) $\cdot$ No Neg]}} \hfill {\footnotesize\textbf{GATE 2025 CS2} \quad \tiny\texttt{2025\_CS2\_CS\_Q50}}
\vspace{2pt}\hrule height 0.3pt\vspace{4pt}

Which of the following Boolean algebraic equation(s) is/are CORRECT?

\par\vspace{6pt}\noindent
\begin{tabular}{@{}p{0.98\linewidth}@{}}
$\square$\ \textbf{(A)} $\overline{A}BC + A\overline{B}\overline{C} + \overline{A}\overline{B}\overline{C} + A\overline{B}C + ABC = BC + \overline{B}\overline{C} + \overline{A}\overline{B}$ \\
\addlinespace[3pt]
$\square$\ \textbf{(B)} $AB + \overline{A}C + BC = AB + \overline{A}C$ \\
\addlinespace[3pt]
$\square$\ \textbf{(C)} $(A + C)(\overline{A} + B) = AB + \overline{A}C$ \\
\addlinespace[3pt]
$\square$\ \textbf{(D)} $\overline{(A + \overline{B} + \overline{D})(C + D)(\overline{A} + C + D)(A + B + \overline{D})} = \overline{A}D + \overline{C}\overline{D}$
\end{tabular}

\tcblower

\textbf{Official Key:} \textbf{\large B;C;D} \hfill \textit{\small [Medium]}
\vspace{2pt}\hrule height 0.3pt\vspace{3pt}
\textbf{Topic:} Digital Logic \\
\textbf{Subtopic:} \textit{Boolean algebra and minimization: identities and consensus theorem}
\vspace{2pt}\hrule height 0.3pt\vspace{3pt}
\textbf{Guide / Hint:} \\
{\small Equation B is the classic Consensus Theorem: $AB + \overline{A}C + BC = AB + \overline{A}C$. Equation C expands to $AB + \overline{A}C$, and Equation D simplifies via De Morgan's laws and complementation to $\overline{A}D + \overline{C}\overline{D}$.}

\end{tcolorbox}
\vspace{6pt}\hrule height 0.35pt\vspace{6pt}
\vspace{8pt}
\subsection*{Computer Organization and Architecture}
\vspace{-4pt}
\noindent
\begin{tcolorbox}[
    blanker,
    sidebyside,
    sidebyside align=top,
    lefthand width=0.55\textwidth,
    top=4pt, bottom=4pt, left=0pt, right=0pt,
    middle=16pt,
    underlay={
        \draw[dashed, black!80, line width=0.8pt, dash pattern=on 4pt off 3pt]
            ([xshift=0.55\textwidth+8pt]frame.north west) --
            ([xshift=0.55\textwidth+8pt]frame.south west);
    }
]
\textbf{Q.32} \quad {\footnotesize\textsc{[MSQ $\cdot$ 1 Mark(s) $\cdot$ No Neg]}} \hfill {\footnotesize\textbf{GATE 2022 CS} \quad \tiny\texttt{2022\_CS\_CS\_Q24}}
\vspace{2pt}\hrule height 0.3pt\vspace{4pt}

Let WB and WT be two set associative cache organizations that use LRU algorithm for cache block replacement. WB is a write back cache and WT is a write through cache. Which of the following statements is/are FALSE?

\par\vspace{6pt}\noindent
\begin{tabular}{@{}p{0.98\linewidth}@{}}
$\square$\ \textbf{(A)} Each cache block in WB and WT has a dirty bit. \\
\addlinespace[3pt]
$\square$\ \textbf{(B)} Every write hit in WB leads to a data transfer from cache to main memory. \\
\addlinespace[3pt]
$\square$\ \textbf{(C)} Eviction of a block from WT will not lead to data transfer from cache to main memory. \\
\addlinespace[3pt]
$\square$\ \textbf{(D)} A read miss in WB will never lead to eviction of a dirty block from WB.
\end{tabular}

\tcblower

\textbf{Official Key:} \textbf{\large A,B,D} \hfill \textit{\small [Medium]}
\vspace{2pt}\hrule height 0.3pt\vspace{3pt}
\textbf{Topic:} Computer Organization and Architecture \\
\textbf{Subtopic:} \textit{Memory hierarchy: cache, main memory and secondary storage}
\vspace{2pt}\hrule height 0.3pt\vspace{3pt}
\textbf{Guide / Hint:} \\
{\small Write-through (WT) caches never maintain dirty blocks because main memory is updated simultaneously on every write, so WT does not have dirty bits (A is FALSE). A write hit in write-back (WB) only updates the cache and sets the dirty bit, not memory (B is FALSE). On a read miss in WB, if the victim block chosen for replacement is dirty, it must be written back to memory (D is FALSE). Statements A, B, and D are FALSE.}

\end{tcolorbox}
\vspace{6pt}\hrule height 0.35pt\vspace{6pt}
\noindent
\begin{tcolorbox}[
    blanker,
    sidebyside,
    sidebyside align=top,
    lefthand width=0.55\textwidth,
    top=4pt, bottom=4pt, left=0pt, right=0pt,
    middle=16pt,
    underlay={
        \draw[dashed, black!80, line width=0.8pt, dash pattern=on 4pt off 3pt]
            ([xshift=0.55\textwidth+8pt]frame.north west) --
            ([xshift=0.55\textwidth+8pt]frame.south west);
    }
]
\textbf{Q.33} \quad {\footnotesize\textsc{[MSQ $\cdot$ 2 Mark(s) $\cdot$ No Neg]}} \hfill {\footnotesize\textbf{GATE 2022 CS} \quad \tiny\texttt{2022\_CS\_CS\_Q54}}
\vspace{2pt}\hrule height 0.3pt\vspace{4pt}

Consider a system with 2 KB direct mapped data cache with a block size of 64 bytes. The system has a physical address space of 64 KB and a word length of 16 bits. During the execution of a program, four data words P, Q, R, and S are accessed in that order 10 times (i.e., PQRSPQRS\dots). Hence, there are 40 accesses to data cache altogether. Assume that the data cache is initially empty and no other data words are accessed by the program. The addresses of the first bytes of P, Q, R, and S are `0xA248`, `0xC28A`, `0xCA8A`, and `0xA262`, respectively. For the execution of the above program, which of the following statements is/are TRUE with respect to the data cache?

\par\vspace{6pt}\noindent
\begin{tabular}{@{}p{0.98\linewidth}@{}}
$\square$\ \textbf{(A)} Every access to S is a hit. \\
\addlinespace[3pt]
$\square$\ \textbf{(B)} Once P is brought to the cache it is never evicted. \\
\addlinespace[3pt]
$\square$\ \textbf{(C)} At the end of the execution only R and S reside in the cache. \\
\addlinespace[3pt]
$\square$\ \textbf{(D)} Every access to R evicts Q from the cache.
\end{tabular}

\tcblower

\textbf{Official Key:} \textbf{\large A,B,D} \hfill \textit{\small [Hard]}
\vspace{2pt}\hrule height 0.3pt\vspace{3pt}
\textbf{Topic:} Computer Organization and Architecture \\
\textbf{Subtopic:} \textit{Memory hierarchy: cache, main memory and secondary storage}
\vspace{2pt}\hrule height 0.3pt\vspace{3pt}
\textbf{Guide / Hint:} \\
{\small Cache line index is calculated as $\lfloor \text{address} / 64 \rfloor \pmod{32}$. Addresses `0xA248` (P) and `0xA262` (S) both evaluate to block offset within block `0xA240`, which maps to line 9. Once P is loaded, S is also present in line 9, so every access to S is a hit (A is TRUE) and P/S are never evicted (B is TRUE). Addresses `0xC28A` (Q) and `0xCA8A` (R) both map to line 10, repeatedly evicting each other (D is TRUE).}

\end{tcolorbox}
\vspace{6pt}\hrule height 0.35pt\vspace{6pt}
\noindent
\begin{tcolorbox}[
    blanker,
    sidebyside,
    sidebyside align=top,
    lefthand width=0.55\textwidth,
    top=4pt, bottom=4pt, left=0pt, right=0pt,
    middle=16pt,
    underlay={
        \draw[dashed, black!80, line width=0.8pt, dash pattern=on 4pt off 3pt]
            ([xshift=0.55\textwidth+8pt]frame.north west) --
            ([xshift=0.55\textwidth+8pt]frame.south west);
    }
]
\textbf{Q.34} \quad {\footnotesize\textsc{[MSQ $\cdot$ 1 Mark(s) $\cdot$ No Neg]}} \hfill {\footnotesize\textbf{GATE 2024 CS1} \quad \tiny\texttt{2024\_CS1\_CS\_Q30}}
\vspace{2pt}\hrule height 0.3pt\vspace{4pt}

Consider a 5-stage pipelined processor with Instruction Fetch (IF), Instruction Decode (ID), Execute (EX), Memory Access (MEM), and Register Writeback (WB) stages. Which of the following statements about forwarding is/are CORRECT?

\par\vspace{6pt}\noindent
\begin{tabular}{@{}p{0.98\linewidth}@{}}
$\square$\ \textbf{(A)} In a pipelined execution, forwarding means the result from a source stage of an earlier instruction is passed on to the destination stage of a later instruction \\
\addlinespace[3pt]
$\square$\ \textbf{(B)} In forwarding, data from the output of the MEM stage can be passed on to the input of the EX stage of the next instruction \\
\addlinespace[3pt]
$\square$\ \textbf{(C)} Forwarding cannot prevent all pipeline stalls \\
\addlinespace[3pt]
$\square$\ \textbf{(D)} Forwarding does not require any extra hardware to retrieve the data from the pipeline stages
\end{tabular}

\tcblower

\textbf{Official Key:} \textbf{\large A;C} \hfill \textit{\small [Easy]}
\vspace{2pt}\hrule height 0.3pt\vspace{3pt}
\textbf{Topic:} Computer Organization and Architecture \\
\textbf{Subtopic:} \textit{Instruction pipelining, pipeline hazards}
\vspace{2pt}\hrule height 0.3pt\vspace{3pt}
\textbf{Guide / Hint:} \\
{\small Forwarding (bypassing) routes intermediate results directly from pipeline registers of an earlier instruction to execution inputs of a subsequent instruction (A is TRUE). While forwarding resolves many ALU data dependencies without stalls, it cannot prevent a load-use stall (where the loaded data is needed immediately in the subsequent instruction's EX stage), requiring at least one stall cycle (C is TRUE).}

\end{tcolorbox}
\vspace{6pt}\hrule height 0.35pt\vspace{6pt}
\noindent
\begin{tcolorbox}[
    blanker,
    sidebyside,
    sidebyside align=top,
    lefthand width=0.55\textwidth,
    top=4pt, bottom=4pt, left=0pt, right=0pt,
    middle=16pt,
    underlay={
        \draw[dashed, black!80, line width=0.8pt, dash pattern=on 4pt off 3pt]
            ([xshift=0.55\textwidth+8pt]frame.north west) --
            ([xshift=0.55\textwidth+8pt]frame.south west);
    }
]
\textbf{Q.35} \quad {\footnotesize\textsc{[MSQ $\cdot$ 2 Mark(s) $\cdot$ No Neg]}} \hfill {\footnotesize\textbf{GATE 2024 CS1} \quad \tiny\texttt{2024\_CS1\_CS\_Q53}}
\vspace{2pt}\hrule height 0.3pt\vspace{4pt}

Consider two set-associative cache memory architectures: WBC, which uses the write back policy, and WTC, which uses the write through policy. Both of them use the LRU (Least Recently Used) block replacement policy. The cache memory is connected to the main memory. Which of the following statements is/are TRUE?

\par\vspace{6pt}\noindent
\begin{tabular}{@{}p{0.98\linewidth}@{}}
$\square$\ \textbf{(A)} A read miss in WBC never evicts a dirty block \\
\addlinespace[3pt]
$\square$\ \textbf{(B)} A read miss in WTC never triggers a write back operation of a cache block to main memory \\
\addlinespace[3pt]
$\square$\ \textbf{(C)} A write hit in WBC can modify the value of the dirty bit of a cache block \\
\addlinespace[3pt]
$\square$\ \textbf{(D)} A write miss in WTC always writes the victim cache block to main memory before loading the missed block to the cache
\end{tabular}

\tcblower

\textbf{Official Key:} \textbf{\large B;C} \hfill \textit{\small [Medium]}
\vspace{2pt}\hrule height 0.3pt\vspace{3pt}
\textbf{Topic:} Computer Organization and Architecture \\
\textbf{Subtopic:} \textit{Memory hierarchy: cache, main memory and secondary storage}
\vspace{2pt}\hrule height 0.3pt\vspace{3pt}
\textbf{Guide / Hint:} \\
{\small In a write-through cache (WTC), main memory is updated simultaneously on every cache write, meaning cache lines are never dirty and no eviction ever triggers a write-back to memory (B is TRUE). In a write-back cache (WBC), when a clean block receives a write hit, its dirty bit is set from 0 to 1 to indicate modification (C is TRUE).}

\end{tcolorbox}
\vspace{6pt}\hrule height 0.35pt\vspace{6pt}
\noindent
\begin{tcolorbox}[
    blanker,
    sidebyside,
    sidebyside align=top,
    lefthand width=0.55\textwidth,
    top=4pt, bottom=4pt, left=0pt, right=0pt,
    middle=16pt,
    underlay={
        \draw[dashed, black!80, line width=0.8pt, dash pattern=on 4pt off 3pt]
            ([xshift=0.55\textwidth+8pt]frame.north west) --
            ([xshift=0.55\textwidth+8pt]frame.south west);
    }
]
\textbf{Q.36} \quad {\footnotesize\textsc{[MSQ $\cdot$ 1 Mark(s) $\cdot$ No Neg]}} \hfill {\footnotesize\textbf{GATE 2024 CS2} \quad \tiny\texttt{2024\_CS2\_CS\_Q31}}
\vspace{2pt}\hrule height 0.3pt\vspace{4pt}

An instruction format has the following structure:
Instruction Number: Opcode destination reg, source reg-1, source reg-2
Consider the following sequence of instructions to be executed in a pipelined processor:
\begin{verbatim}
I1: DIV R3, R1, R2
I2: SUB R5, R3, R4
I3: ADD R3, R5, R6
I4: MUL R7, R3, R8
\end{verbatim}
Which of the following dependencies is/are TRUE?

\par\vspace{6pt}\noindent
\begin{tabular}{@{}p{0.98\linewidth}@{}}
$\square$\ \textbf{(A)} There is a RAW dependency on R3 between I1 and I2 \\
\addlinespace[3pt]
$\square$\ \textbf{(B)} There is a WAR dependency on R3 between I1 and I3 \\
\addlinespace[3pt]
$\square$\ \textbf{(C)} There is a RAW dependency on R3 between I2 and I3 \\
\addlinespace[3pt]
$\square$\ \textbf{(D)} There is a WAW dependency on R3 between I3 and I4
\end{tabular}

\tcblower

\textbf{Official Key:} \textbf{\large A} \hfill \textit{\small [Easy]}
\vspace{2pt}\hrule height 0.3pt\vspace{3pt}
\textbf{Topic:} Computer Organization and Architecture \\
\textbf{Subtopic:} \textit{Instruction pipelining, pipeline hazards}
\vspace{2pt}\hrule height 0.3pt\vspace{3pt}
\textbf{Guide / Hint:} \\
{\small In I1: `DIV R3, R1, R2` (writes R3); I2: `SUB R5, R3, R4` (reads R3). Since I2 reads R3 after I1 writes R3, this is a Read-After-Write (RAW) data dependency on R3 between I1 and I2 (A is true). I3 writes R3, so between I1 (write) and I3 (write) is WAW, and between I2 (read) and I3 (write) is WAR.}

\end{tcolorbox}
\vspace{6pt}\hrule height 0.35pt\vspace{6pt}
\noindent
\begin{tcolorbox}[
    blanker,
    sidebyside,
    sidebyside align=top,
    lefthand width=0.55\textwidth,
    top=4pt, bottom=4pt, left=0pt, right=0pt,
    middle=16pt,
    underlay={
        \draw[dashed, black!80, line width=0.8pt, dash pattern=on 4pt off 3pt]
            ([xshift=0.55\textwidth+8pt]frame.north west) --
            ([xshift=0.55\textwidth+8pt]frame.south west);
    }
]
\textbf{Q.37} \quad {\footnotesize\textsc{[MSQ $\cdot$ 1 Mark(s) $\cdot$ No Neg]}} \hfill {\footnotesize\textbf{GATE 2025 CS1} \quad \tiny\texttt{2025\_CS1\_CS\_Q27}}
\vspace{2pt}\hrule height 0.3pt\vspace{4pt}

A partial data path of a processor is given, where RA, RB, and RZ are 32-bit registers. Multiplexer Mux\_A selects between RA and a 32-bit immediate value, and Mux\_B selects between RB and a 32-bit immediate value. The output of the ALU is stored into RZ. Which option(s) is/are CORRECT related to arithmetic operations using the data path as shown?

\par\vspace{6pt}\noindent
\begin{tabular}{@{}p{0.98\linewidth}@{}}
$\square$\ \textbf{(A)} The data path can implement arithmetic operations involving two registers. \\
\addlinespace[3pt]
$\square$\ \textbf{(B)} The data path can implement arithmetic operations involving one register and one immediate value. \\
\addlinespace[3pt]
$\square$\ \textbf{(C)} The data path can implement arithmetic operations involving two immediate values. \\
\addlinespace[3pt]
$\square$\ \textbf{(D)} The data path can only implement arithmetic operations involving one register and one immediate value.
\end{tabular}

\tcblower

\textbf{Official Key:} \textbf{\large A;B;C} \hfill \textit{\small [Medium]}
\vspace{2pt}\hrule height 0.3pt\vspace{3pt}
\textbf{Topic:} Computer Organization and Architecture \\
\textbf{Subtopic:} \textit{ALU, data-path and control unit}
\vspace{2pt}\hrule height 0.3pt\vspace{3pt}
\textbf{Guide / Hint:} \\
{\small Since Mux\_A can select either register RA or an immediate value, and Mux\_B can independently select either register RB or an immediate value, the ALU inputs can be configured as (RA, RB), (RA, immediate), (immediate, RB), or (immediate, immediate). Thus A, B, and C are all valid.}

\end{tcolorbox}
\vspace{6pt}\hrule height 0.35pt\vspace{6pt}
\noindent
\begin{tcolorbox}[
    blanker,
    sidebyside,
    sidebyside align=top,
    lefthand width=0.55\textwidth,
    top=4pt, bottom=4pt, left=0pt, right=0pt,
    middle=16pt,
    underlay={
        \draw[dashed, black!80, line width=0.8pt, dash pattern=on 4pt off 3pt]
            ([xshift=0.55\textwidth+8pt]frame.north west) --
            ([xshift=0.55\textwidth+8pt]frame.south west);
    }
]
\textbf{Q.38} \quad {\footnotesize\textsc{[MSQ $\cdot$ 1 Mark(s) $\cdot$ No Neg]}} \hfill {\footnotesize\textbf{GATE 2025 CS2} \quad \tiny\texttt{2025\_CS2\_CS\_Q28}}
\vspace{2pt}\hrule height 0.3pt\vspace{4pt}

Which of the following is/are part of an Instruction Set Architecture of a processor?

\par\vspace{6pt}\noindent
\begin{tabular}{@{}p{0.98\linewidth}@{}}
$\square$\ \textbf{(A)} The size of the cache memory \\
\addlinespace[3pt]
$\square$\ \textbf{(B)} The clock frequency of the processor \\
\addlinespace[3pt]
$\square$\ \textbf{(C)} The number of cache memory levels \\
\addlinespace[3pt]
$\square$\ \textbf{(D)} The total number of registers
\end{tabular}

\tcblower

\textbf{Official Key:} \textbf{\large D} \hfill \textit{\small [Easy]}
\vspace{2pt}\hrule height 0.3pt\vspace{3pt}
\textbf{Topic:} Computer Organization and Architecture \\
\textbf{Subtopic:} \textit{Machine instructions and addressing modes: ISA}
\vspace{2pt}\hrule height 0.3pt\vspace{3pt}
\textbf{Guide / Hint:} \\
{\small The Instruction Set Architecture (ISA) defines programmer-visible specifications, including the register set/count and instruction encoding. Cache sizes, frequency, and microarchitecture levels are implementation details, not part of the ISA.}

\end{tcolorbox}
\vspace{6pt}\hrule height 0.35pt\vspace{6pt}
\vspace{8pt}
\subsection*{Programming and Data Structures}
\vspace{-4pt}
\noindent
\begin{tcolorbox}[
    blanker,
    sidebyside,
    sidebyside align=top,
    lefthand width=0.55\textwidth,
    top=4pt, bottom=4pt, left=0pt, right=0pt,
    middle=16pt,
    underlay={
        \draw[dashed, black!80, line width=0.8pt, dash pattern=on 4pt off 3pt]
            ([xshift=0.55\textwidth+8pt]frame.north west) --
            ([xshift=0.55\textwidth+8pt]frame.south west);
    }
]
\textbf{Q.39} \quad {\footnotesize\textsc{[MSQ $\cdot$ 2 Mark(s) $\cdot$ No Neg]}} \hfill {\footnotesize\textbf{GATE 2024 CS1} \quad \tiny\texttt{2024\_CS1\_CS\_Q48}}
\vspace{2pt}\hrule height 0.3pt\vspace{4pt}

Consider the following C function definition:
\begin{lstlisting}[language=C]
int f(int x, int y) {
    for (int i = 0; i < y; i++) {
        x = x + x + y;
    }
    return x;
}
\end{lstlisting}
Which of the following statements is/are TRUE about the above function?

\par\vspace{6pt}\noindent
\begin{tabular}{@{}p{0.98\linewidth}@{}}
$\square$\ \textbf{(A)} If the inputs are $x=20, y=10$, then the return value is greater than $2^{20}$ \\
\addlinespace[3pt]
$\square$\ \textbf{(B)} If the inputs are $x=20, y=20$, then the return value is greater than $2^{20}$ \\
\addlinespace[3pt]
$\square$\ \textbf{(C)} If the inputs are $x=20, y=10$, then the return value is less than $2^{10}$ \\
\addlinespace[3pt]
$\square$\ \textbf{(D)} If the inputs are $x=10, y=20$, then the return value is greater than $2^{20}$
\end{tabular}

\tcblower

\textbf{Official Key:} \textbf{\large B;D} \hfill \textit{\small [Medium]}
\vspace{2pt}\hrule height 0.3pt\vspace{3pt}
\textbf{Topic:} Programming and Data Structures \\
\textbf{Subtopic:} \textit{Programming in C}
\vspace{2pt}\hrule height 0.3pt\vspace{3pt}
\textbf{Guide / Hint:} \\
{\small In each loop iteration, $x_{k+1} = 2x_k + y$. Solving this linear recurrence yields $x_y = 2^y x_0 + y(2^y - 1) = 2^y(x_0 + y) - y$. For $x=20, y=20$: $x_{20} = 2^{20}(40) - 20 > 2^{20}$ (B is TRUE). For $x=10, y=20$: $x_{20} = 2^{20}(30) - 20 > 2^{20}$ (D is TRUE). For $y=10, x=20$: $x_{10} = 2^{10}(30) - 10 = 30710 > 2^{10}$, so (A) and (C) are false.}

\end{tcolorbox}
\vspace{6pt}\hrule height 0.35pt\vspace{6pt}
\noindent
\begin{tcolorbox}[
    blanker,
    sidebyside,
    sidebyside align=top,
    lefthand width=0.55\textwidth,
    top=4pt, bottom=4pt, left=0pt, right=0pt,
    middle=16pt,
    underlay={
        \draw[dashed, black!80, line width=0.8pt, dash pattern=on 4pt off 3pt]
            ([xshift=0.55\textwidth+8pt]frame.north west) --
            ([xshift=0.55\textwidth+8pt]frame.south west);
    }
]
\textbf{Q.40} \quad {\footnotesize\textsc{[MSQ $\cdot$ 1 Mark(s) $\cdot$ No Neg]}} \hfill {\footnotesize\textbf{GATE 2024 CS2} \quad \tiny\texttt{2024\_CS2\_CS\_Q33}}
\vspace{2pt}\hrule height 0.3pt\vspace{4pt}

Consider the following C function definition:
\begin{verbatim}
int fX(char *a){
    char *b = a;
    while(*b)
        b++;
    return b - a;
}
\end{verbatim}
Which of the following statements is/are TRUE?

\par\vspace{6pt}\noindent
\begin{tabular}{@{}p{0.98\linewidth}@{}}
$\square$\ \textbf{(A)} The function call fX("abcd") will always return a value \\
\addlinespace[3pt]
$\square$\ \textbf{(B)} Assuming a character array c is declared as char c[] = "abcd" in main(), the function call fX(c) will always return a value \\
\addlinespace[3pt]
$\square$\ \textbf{(C)} The code of the function will not compile \\
\addlinespace[3pt]
$\square$\ \textbf{(D)} Assuming a character pointer c is declared as char *c = "abcd" in main(), the function call fX(c) will always return a value
\end{tabular}

\tcblower

\textbf{Official Key:} \textbf{\large A;B;D} \hfill \textit{\small [Easy]}
\vspace{2pt}\hrule height 0.3pt\vspace{3pt}
\textbf{Topic:} Programming and Data Structures \\
\textbf{Subtopic:} \textit{Programming in C}
\vspace{2pt}\hrule height 0.3pt\vspace{3pt}
\textbf{Guide / Hint:} \\
{\small The function `fX` calculates the string length by advancing pointer `b` until the null terminator `\0` is reached, then returning pointer difference `b - a`. String literals in C are null-terminated, so calls with string literals or character arrays (`fX("abcd")`, `fX(c)`) correctly compute 4 and return (A, B, and D are true).}

\end{tcolorbox}
\vspace{6pt}\hrule height 0.35pt\vspace{6pt}
\noindent
\begin{tcolorbox}[
    blanker,
    sidebyside,
    sidebyside align=top,
    lefthand width=0.55\textwidth,
    top=4pt, bottom=4pt, left=0pt, right=0pt,
    middle=16pt,
    underlay={
        \draw[dashed, black!80, line width=0.8pt, dash pattern=on 4pt off 3pt]
            ([xshift=0.55\textwidth+8pt]frame.north west) --
            ([xshift=0.55\textwidth+8pt]frame.south west);
    }
]
\textbf{Q.41} \quad {\footnotesize\textsc{[MSQ $\cdot$ 2 Mark(s) $\cdot$ No Neg]}} \hfill {\footnotesize\textbf{GATE 2024 CS2} \quad \tiny\texttt{2024\_CS2\_CS\_Q48}}
\vspace{2pt}\hrule height 0.3pt\vspace{4pt}

Let S1 and S2 be two stacks. S1 has capacity of 4 elements. S2 has capacity of 2 elements. S1 already has 4 elements: 100, 200, 300, and 400 (top), whereas S2 is empty.
Only the following three operations are available: PushToS2, PushToS1, and GenerateOutput.
Which of the following output sequences can be generated?

\par\vspace{6pt}\noindent
\begin{tabular}{@{}p{0.48\linewidth}@{\hspace{0.04\linewidth}}p{0.48\linewidth}@{}}
$\square$\ \textbf{(A)} 100, 200, 400, 300 & $\square$\ \textbf{(B)} 200, 300, 400, 100 \\
\addlinespace[3pt]
$\square$\ \textbf{(C)} 400, 200, 100, 300 & $\square$\ \textbf{(D)} 300, 200, 400, 100
\end{tabular}

\tcblower

\textbf{Official Key:} \textbf{\large B;C;D} \hfill \textit{\small [Medium]}
\vspace{2pt}\hrule height 0.3pt\vspace{3pt}
\textbf{Topic:} Programming and Data Structures \\
\textbf{Subtopic:} \textit{Arrays, stacks, queues, linked lists}
\vspace{2pt}\hrule height 0.3pt\vspace{3pt}
\textbf{Guide / Hint:} \\
{\small S1 initially contains [100, 200, 300, 400] from bottom to top; S2 has capacity 2. Sequence (B): Push 400 to S2, push 300 to S2, output 200, pop 300 to S1 and output 300, pop 400 to S1 and output 400, output 100. Similarly, sequences (C) and (D) can be generated within S2's 2-element capacity limit.}

\end{tcolorbox}
\vspace{6pt}\hrule height 0.35pt\vspace{6pt}
\noindent
\begin{tcolorbox}[
    blanker,
    sidebyside,
    sidebyside align=top,
    lefthand width=0.55\textwidth,
    top=4pt, bottom=4pt, left=0pt, right=0pt,
    middle=16pt,
    underlay={
        \draw[dashed, black!80, line width=0.8pt, dash pattern=on 4pt off 3pt]
            ([xshift=0.55\textwidth+8pt]frame.north west) --
            ([xshift=0.55\textwidth+8pt]frame.south west);
    }
]
\textbf{Q.42} \quad {\footnotesize\textsc{[MSQ $\cdot$ 1 Mark(s) $\cdot$ No Neg]}} \hfill {\footnotesize\textbf{GATE 2025 CS1} \quad \tiny\texttt{2025\_CS1\_CS\_Q26}}
\vspace{2pt}\hrule height 0.3pt\vspace{4pt}

Which of the following statement(s) is/are TRUE for any binary search tree (BST) having $n$ distinct integers?

\par\vspace{6pt}\noindent
\begin{tabular}{@{}p{0.98\linewidth}@{}}
$\square$\ \textbf{(A)} The maximum length of a path from the root node to any other node is $(n-1)$. \\
\addlinespace[3pt]
$\square$\ \textbf{(B)} An inorder traversal will always produce a sorted sequence of elements. \\
\addlinespace[3pt]
$\square$\ \textbf{(C)} Finding an element takes $O(\log_2 n)$ time in the worst case. \\
\addlinespace[3pt]
$\square$\ \textbf{(D)} Every BST is also a Min-Heap.
\end{tabular}

\tcblower

\textbf{Official Key:} \textbf{\large A;B} \hfill \textit{\small [Easy]}
\vspace{2pt}\hrule height 0.3pt\vspace{3pt}
\textbf{Topic:} Programming and Data Structures \\
\textbf{Subtopic:} \textit{Trees, binary search trees, binary heaps}
\vspace{2pt}\hrule height 0.3pt\vspace{3pt}
\textbf{Guide / Hint:} \\
{\small In a degenerate (skewed) binary search tree, the path length from root to the farthest leaf is $n - 1$ edges (A is true). Inorder traversal of any valid BST visits nodes in strictly increasing order (B is true). Worst-case search in an unbalanced BST takes $O(n)$ time, and BST property differs fundamentally from heap order.}

\end{tcolorbox}
\vspace{6pt}\hrule height 0.35pt\vspace{6pt}
\noindent
\begin{tcolorbox}[
    blanker,
    sidebyside,
    sidebyside align=top,
    lefthand width=0.55\textwidth,
    top=4pt, bottom=4pt, left=0pt, right=0pt,
    middle=16pt,
    underlay={
        \draw[dashed, black!80, line width=0.8pt, dash pattern=on 4pt off 3pt]
            ([xshift=0.55\textwidth+8pt]frame.north west) --
            ([xshift=0.55\textwidth+8pt]frame.south west);
    }
]
\textbf{Q.43} \quad {\footnotesize\textsc{[MSQ $\cdot$ 2 Mark(s) $\cdot$ No Neg]}} \hfill {\footnotesize\textbf{GATE 2025 CS2} \quad \tiny\texttt{2025\_CS2\_CS\_Q45}}
\vspace{2pt}\hrule height 0.3pt\vspace{4pt}

Consider a stack data structure into which we can PUSH and POP records. Assume that each record has a numeric key. Which of the following strategies allows finding the record with the minimum key in $O(1)$ worst-case time without increasing the asymptotic time for PUSH and POP?

\par\vspace{6pt}\noindent
\begin{tabular}{@{}p{0.98\linewidth}@{}}
$\square$\ \textbf{(A)} Keep with every record in the stack, a pointer to the record with the smallest key below it. \\
\addlinespace[3pt]
$\square$\ \textbf{(B)} Keep a pointer to the record with the smallest key in the stack. \\
\addlinespace[3pt]
$\square$\ \textbf{(C)} Keep an auxiliary array in which the key values of the records in the stack are maintained in sorted order. \\
\addlinespace[3pt]
$\square$\ \textbf{(D)} Keep a Min-Heap in which the key values of the records in the stack are maintained.
\end{tabular}

\tcblower

\textbf{Official Key:} \textbf{\large A} \hfill \textit{\small [Medium]}
\vspace{2pt}\hrule height 0.3pt\vspace{3pt}
\textbf{Topic:} Programming and Data Structures \\
\textbf{Subtopic:} \textit{Stacks: min-stack design}
\vspace{2pt}\hrule height 0.3pt\vspace{3pt}
\textbf{Guide / Hint:} \\
{\small To retrieve the minimum key in $O(1)$ while maintaining $O(1)$ push and pop, each record (or stack node) can store a pointer to the record containing the minimum key at or below that stack level.}

\end{tcolorbox}
\vspace{6pt}\hrule height 0.35pt\vspace{6pt}
\vspace{8pt}
\subsection*{Algorithms}
\vspace{-4pt}
\noindent
\begin{tcolorbox}[
    blanker,
    sidebyside,
    sidebyside align=top,
    lefthand width=0.55\textwidth,
    top=4pt, bottom=4pt, left=0pt, right=0pt,
    middle=16pt,
    underlay={
        \draw[dashed, black!80, line width=0.8pt, dash pattern=on 4pt off 3pt]
            ([xshift=0.55\textwidth+8pt]frame.north west) --
            ([xshift=0.55\textwidth+8pt]frame.south west);
    }
]
\textbf{Q.44} \quad {\footnotesize\textsc{[MSQ $\cdot$ 2 Mark(s) $\cdot$ No Neg]}} \hfill {\footnotesize\textbf{GATE 2022 CS} \quad \tiny\texttt{2022\_CS\_CS\_Q49}}
\vspace{2pt}\hrule height 0.3pt\vspace{4pt}

Consider a simple undirected weighted graph $G$, all of whose edge weights are distinct. Which of the following statements about the minimum spanning trees of $G$ is/are TRUE?

\par\vspace{6pt}\noindent
\begin{tabular}{@{}p{0.98\linewidth}@{}}
$\square$\ \textbf{(A)} The edge with the second smallest weight is always part of any minimum spanning tree of $G$. \\
\addlinespace[3pt]
$\square$\ \textbf{(B)} One or both of the edges with the third smallest and the fourth smallest weights are part of any minimum spanning tree of $G$. \\
\addlinespace[3pt]
$\square$\ \textbf{(C)} Suppose $S \subset V$ be such that $S \neq \emptyset$ and $S \neq V$. Consider the edge with the minimum weight such that one of its vertices is in $S$ and the other in $V \setminus S$. Such an edge will always be part of any minimum spanning tree of $G$. \\
\addlinespace[3pt]
$\square$\ \textbf{(D)} $G$ can have multiple minimum spanning trees.
\end{tabular}

\tcblower

\textbf{Official Key:} \textbf{\large A,B,C} \hfill \textit{\small [Medium]}
\vspace{2pt}\hrule height 0.3pt\vspace{3pt}
\textbf{Topic:} Algorithms \\
\textbf{Subtopic:} \textit{Graph traversals, minimum spanning trees, shortest paths}
\vspace{2pt}\hrule height 0.3pt\vspace{3pt}
\textbf{Guide / Hint:} \\
{\small The two edges with the smallest weights can never form a cycle, so both belong to the MST (A is TRUE). The three smallest edges can at most form one cycle (triangle), so at least one of the 3rd and 4th smallest edges must be part of the MST (B is TRUE). By the Cut Property, the minimum weight edge across any non-trivial cut $(S, V \setminus S)$ belongs to the MST (C is TRUE). Distinct weights guarantee a unique MST (D is FALSE).}

\end{tcolorbox}
\vspace{6pt}\hrule height 0.35pt\vspace{6pt}
\noindent
\begin{tcolorbox}[
    blanker,
    sidebyside,
    sidebyside align=top,
    lefthand width=0.55\textwidth,
    top=4pt, bottom=4pt, left=0pt, right=0pt,
    middle=16pt,
    underlay={
        \draw[dashed, black!80, line width=0.8pt, dash pattern=on 4pt off 3pt]
            ([xshift=0.55\textwidth+8pt]frame.north west) --
            ([xshift=0.55\textwidth+8pt]frame.south west);
    }
]
\textbf{Q.45} \quad {\footnotesize\textsc{[MSQ $\cdot$ 2 Mark(s) $\cdot$ No Neg]}} \hfill {\footnotesize\textbf{GATE 2022 CS} \quad \tiny\texttt{2022\_CS\_CS\_Q51}}
\vspace{2pt}\hrule height 0.3pt\vspace{4pt}

Consider the following recurrence:
\[ f(1) = 1; \]
\[ f(2n) = 2f(n) - 1, \quad \text{for } n \ge 1; \]
\[ f(2n+1) = 2f(n) + 1, \quad \text{for } n \ge 1. \]
Then, which of the following statements is/are TRUE?

\par\vspace{6pt}\noindent
\begin{tabular}{@{}p{0.48\linewidth}@{\hspace{0.04\linewidth}}p{0.48\linewidth}@{}}
$\square$\ \textbf{(A)} $f(2^n - 1) = 2^n - 1$ & $\square$\ \textbf{(B)} $f(2^n) = 1$ \\
\addlinespace[3pt]
$\square$\ \textbf{(C)} $f(5 \cdot 2^n) = 2^{n+1} + 1$ & $\square$\ \textbf{(D)} $f(2^n + 1) = 2^n + 1$
\end{tabular}

\tcblower

\textbf{Official Key:} \textbf{\large A,B,C} \hfill \textit{\small [Medium]}
\vspace{2pt}\hrule height 0.3pt\vspace{3pt}
\textbf{Topic:} Algorithms \\
\textbf{Subtopic:} \textit{Asymptotic worst case time and space complexity}
\vspace{2pt}\hrule height 0.3pt\vspace{3pt}
\textbf{Guide / Hint:} \\
{\small Computing small values: $f(1)=1, f(2)=1, f(3)=3, f(4)=1, f(5)=5, f(6)=5, f(7)=7, f(8)=1$. In binary, $f(n)$ corresponds to cyclic left shift of $n$. For $2^n - 1$ (all 1s in binary), $f(2^n - 1) = 2^n - 1$ (A is TRUE). For powers of 2, $f(2^n) = 1$ (B is TRUE). For $5 \cdot 2^n$, binary representation is $10100\dots$, shifting gives $100\dots01_2 = 2^{n+1} + 1$ (C is TRUE).}

\end{tcolorbox}
\vspace{6pt}\hrule height 0.35pt\vspace{6pt}
\noindent
\begin{tcolorbox}[
    blanker,
    sidebyside,
    sidebyside align=top,
    lefthand width=0.55\textwidth,
    top=4pt, bottom=4pt, left=0pt, right=0pt,
    middle=16pt,
    underlay={
        \draw[dashed, black!80, line width=0.8pt, dash pattern=on 4pt off 3pt]
            ([xshift=0.55\textwidth+8pt]frame.north west) --
            ([xshift=0.55\textwidth+8pt]frame.south west);
    }
]
\textbf{Q.46} \quad {\footnotesize\textsc{[MSQ $\cdot$ 1 Mark(s) $\cdot$ No Neg]}} \hfill {\footnotesize\textbf{GATE 2023 CS} \quad \tiny\texttt{2023\_CS\_CS\_Q29}}
\vspace{2pt}\hrule height 0.3pt\vspace{4pt}

Let $f$ and $g$ be functions of natural numbers given by $f(n) = n$ and $g(n) = n^2$.
Which of the following statements is/are TRUE?

\par\vspace{6pt}\noindent
\begin{tabular}{@{}p{0.48\linewidth}@{\hspace{0.04\linewidth}}p{0.48\linewidth}@{}}
$\square$\ \textbf{(A)} $f \in \mathcal{O}(g)$ & $\square$\ \textbf{(B)} $f \in \Omega(g)$ \\
\addlinespace[3pt]
$\square$\ \textbf{(C)} $f \in o(g)$ & $\square$\ \textbf{(D)} $f \in \Theta(g)$
\end{tabular}

\tcblower

\textbf{Official Key:} \textbf{\large A, C} \hfill \textit{\small [Easy]}
\vspace{2pt}\hrule height 0.3pt\vspace{3pt}
\textbf{Topic:} Algorithms \\
\textbf{Subtopic:} \textit{Asymptotic worst case time and space complexity}
\vspace{2pt}\hrule height 0.3pt\vspace{3pt}
\textbf{Guide / Hint:} \\
{\small For $f(n) = n$ and $g(n) = n^2$, we evaluate $\lim_{n \to \infty} \frac{f(n)}{g(n)} = \lim_{n \to \infty} \frac{1}{n} = 0$. By definition, if the limit is 0, $f(n) \in o(g(n))$, which strictly implies $f(n) \in \mathcal{O}(g(n))$. Hence, both A and C are TRUE.}

\end{tcolorbox}
\vspace{6pt}\hrule height 0.35pt\vspace{6pt}
\noindent
\begin{tcolorbox}[
    blanker,
    sidebyside,
    sidebyside align=top,
    lefthand width=0.55\textwidth,
    top=4pt, bottom=4pt, left=0pt, right=0pt,
    middle=16pt,
    underlay={
        \draw[dashed, black!80, line width=0.8pt, dash pattern=on 4pt off 3pt]
            ([xshift=0.55\textwidth+8pt]frame.north west) --
            ([xshift=0.55\textwidth+8pt]frame.south west);
    }
]
\textbf{Q.47} \quad {\footnotesize\textsc{[MSQ $\cdot$ 2 Mark(s) $\cdot$ No Neg]}} \hfill {\footnotesize\textbf{GATE 2023 CS} \quad \tiny\texttt{2023\_CS\_CS\_Q54}}
\vspace{2pt}\hrule height 0.3pt\vspace{4pt}

Consider functions Function 1 and Function 2 expressed in pseudocode as follows:
\begin{lstlisting}
// Function 1
while n > 1 do
    for i = 1 to n do
        x = x + 1;
    end for
    n = floor(n / 2);
end while
\par\medskip
// Function 2
for i = 1 to 100 * n do
    x = x + 1;
end for
\end{lstlisting}
Let $f_1(n)$ and $f_2(n)$ denote the number of times the statement ``$x = x + 1$'' is executed in Function 1 and Function 2, respectively.
Which of the following statements is/are TRUE?

\par\vspace{6pt}\noindent
\begin{tabular}{@{}p{0.48\linewidth}@{\hspace{0.04\linewidth}}p{0.48\linewidth}@{}}
$\square$\ \textbf{(A)} $f_1(n) \in \Theta(f_2(n))$ & $\square$\ \textbf{(B)} $f_1(n) \in o(f_2(n))$ \\
\addlinespace[3pt]
$\square$\ \textbf{(C)} $f_1(n) \in \omega(f_2(n))$ & $\square$\ \textbf{(D)} $f_1(n) \in \mathcal{O}(n)$
\end{tabular}

\tcblower

\textbf{Official Key:} \textbf{\large A, D} \hfill \textit{\small [Easy]}
\vspace{2pt}\hrule height 0.3pt\vspace{3pt}
\textbf{Topic:} Algorithms \\
\textbf{Subtopic:} \textit{Asymptotic worst case time and space complexity}
\vspace{2pt}\hrule height 0.3pt\vspace{3pt}
\textbf{Guide / Hint:} \\
{\small In Function 1, the inner loop executes $n + \lfloor n/2 \rfloor + \lfloor n/4 \rfloor + \dots \approx 2n$ times, which gives $f_1(n) = \Theta(n)$. In Function 2, the loop runs $100n$ times, which gives $f_2(n) = \Theta(n)$. Because both functions have exact linear asymptotic growth, $f_1(n) \in \Theta(f_2(n))$ (A is TRUE) and $f_1(n) \in \mathcal{O}(n)$ (D is TRUE).}

\end{tcolorbox}
\vspace{6pt}\hrule height 0.35pt\vspace{6pt}
\noindent
\begin{tcolorbox}[
    blanker,
    sidebyside,
    sidebyside align=top,
    lefthand width=0.55\textwidth,
    top=4pt, bottom=4pt, left=0pt, right=0pt,
    middle=16pt,
    underlay={
        \draw[dashed, black!80, line width=0.8pt, dash pattern=on 4pt off 3pt]
            ([xshift=0.55\textwidth+8pt]frame.north west) --
            ([xshift=0.55\textwidth+8pt]frame.south west);
    }
]
\textbf{Q.48} \quad {\footnotesize\textsc{[MSQ $\cdot$ 2 Mark(s) $\cdot$ No Neg]}} \hfill {\footnotesize\textbf{GATE 2024 CS1} \quad \tiny\texttt{2024\_CS1\_CS\_Q45}}
\vspace{2pt}\hrule height 0.3pt\vspace{4pt}

Let $G$ be a directed graph and $T$ a depth first search (DFS) spanning tree in $G$ that is rooted at a vertex $v$. Suppose $T$ is also a breadth first search (BFS) tree in $G$, rooted at $v$. Which of the following statements is/are TRUE for every such graph $G$ and tree $T$?

\par\vspace{6pt}\noindent
\begin{tabular}{@{}p{0.98\linewidth}@{}}
$\square$\ \textbf{(A)} There are no back-edges in $G$ with respect to the tree $T$ \\
\addlinespace[3pt]
$\square$\ \textbf{(B)} There are no cross-edges in $G$ with respect to the tree $T$ \\
\addlinespace[3pt]
$\square$\ \textbf{(C)} There are no forward-edges in $G$ with respect to the tree $T$ \\
\addlinespace[3pt]
$\square$\ \textbf{(D)} The only edges in $G$ are the edges in $T$
\end{tabular}

\tcblower

\textbf{Official Key:} \textbf{\large C} \hfill \textit{\small [Hard]}
\vspace{2pt}\hrule height 0.3pt\vspace{3pt}
\textbf{Topic:} Algorithms \\
\textbf{Subtopic:} \textit{Graph traversals, minimum spanning trees, shortest paths}
\vspace{2pt}\hrule height 0.3pt\vspace{3pt}
\textbf{Guide / Hint:} \\
{\small In a BFS tree $T$, any edge $(u, w)$ in $G$ must connect vertices whose BFS levels differ by at most 1 (i.e., $\text{level}(w) \le \text{level}(u) + 1$). A forward-edge in a DFS tree goes from ancestor $u$ to descendant $w$ with $\text{level}(w) \ge \text{level}(u) + 2$. If such an edge existed in $G$, BFS would explore $w$ directly from $u$ at depth $\text{level}(u) + 1$, contradicting the structure of $T$. Hence, no forward-edges can exist in $G$ with respect to $T$ (C is TRUE).}

\end{tcolorbox}
\vspace{6pt}\hrule height 0.35pt\vspace{6pt}
\noindent
\begin{tcolorbox}[
    blanker,
    sidebyside,
    sidebyside align=top,
    lefthand width=0.55\textwidth,
    top=4pt, bottom=4pt, left=0pt, right=0pt,
    middle=16pt,
    underlay={
        \draw[dashed, black!80, line width=0.8pt, dash pattern=on 4pt off 3pt]
            ([xshift=0.55\textwidth+8pt]frame.north west) --
            ([xshift=0.55\textwidth+8pt]frame.south west);
    }
]
\textbf{Q.49} \quad {\footnotesize\textsc{[MSQ $\cdot$ 1 Mark(s) $\cdot$ No Neg]}} \hfill {\footnotesize\textbf{GATE 2025 CS2} \quad \tiny\texttt{2025\_CS2\_CS\_Q29}}
\vspace{2pt}\hrule height 0.3pt\vspace{4pt}

Which of the following statements regarding Breadth First Search (BFS) and Depth First Search (DFS) on an undirected simple graph $G$ is/are TRUE?

\par\vspace{6pt}\noindent
\begin{tabular}{@{}p{0.98\linewidth}@{}}
$\square$\ \textbf{(A)} A DFS tree of $G$ is a Shortest Path tree of $G$. \\
\addlinespace[3pt]
$\square$\ \textbf{(B)} Every non-tree edge of $G$ with respect to a DFS tree is a forward/back edge. \\
\addlinespace[3pt]
$\square$\ \textbf{(C)} If $(u, v)$ is a non-tree edge of $G$ with respect to a BFS tree, then the distances from the source vertex $s$ to $u$ and $v$ in the BFS tree are within $\pm 1$ of each other. \\
\addlinespace[3pt]
$\square$\ \textbf{(D)} Both BFS and DFS can be used to find the connected components of $G$.
\end{tabular}

\tcblower

\textbf{Official Key:} \textbf{\large B;C;D} \hfill \textit{\small [Medium]}
\vspace{2pt}\hrule height 0.3pt\vspace{3pt}
\textbf{Topic:} Algorithms \\
\textbf{Subtopic:} \textit{Graph traversals: BFS and DFS}
\vspace{2pt}\hrule height 0.3pt\vspace{3pt}
\textbf{Guide / Hint:} \\
{\small In an undirected graph, non-tree DFS edges can only be back edges (B is true). In a BFS tree, cross edges can only exist between the same or adjacent levels, so distances differ by at most $\pm 1$ (C is true). Both can compute connected components (D is true).}

\end{tcolorbox}
\vspace{6pt}\hrule height 0.35pt\vspace{6pt}
\vspace{8pt}
\subsection*{Theory of Computation}
\vspace{-4pt}
\noindent
\begin{tcolorbox}[
    blanker,
    sidebyside,
    sidebyside align=top,
    lefthand width=0.55\textwidth,
    top=4pt, bottom=4pt, left=0pt, right=0pt,
    middle=16pt,
    underlay={
        \draw[dashed, black!80, line width=0.8pt, dash pattern=on 4pt off 3pt]
            ([xshift=0.55\textwidth+8pt]frame.north west) --
            ([xshift=0.55\textwidth+8pt]frame.south west);
    }
]
\textbf{Q.50} \quad {\footnotesize\textsc{[MSQ $\cdot$ 1 Mark(s) $\cdot$ No Neg]}} \hfill {\footnotesize\textbf{GATE 2022 CS} \quad \tiny\texttt{2022\_CS\_CS\_Q23}}
\vspace{2pt}\hrule height 0.3pt\vspace{4pt}

Which of the following statements is/are TRUE?

\par\vspace{6pt}\noindent
\begin{tabular}{@{}p{0.98\linewidth}@{}}
$\square$\ \textbf{(A)} Every subset of a recursively enumerable language is recursive. \\
\addlinespace[3pt]
$\square$\ \textbf{(B)} If a language $L$ and its complement $\overline{L}$ are both recursively enumerable, then $L$ must be recursive. \\
\addlinespace[3pt]
$\square$\ \textbf{(C)} Complement of a context-free language must be recursive. \\
\addlinespace[3pt]
$\square$\ \textbf{(D)} If $L_1$ and $L_2$ are regular, then $L_1 \cap L_2$ must be deterministic context-free.
\end{tabular}

\tcblower

\textbf{Official Key:} \textbf{\large B,C,D} \hfill \textit{\small [Medium]}
\vspace{2pt}\hrule height 0.3pt\vspace{3pt}
\textbf{Topic:} Theory of Computation \\
\textbf{Subtopic:} \textit{Regular and context-free languages, pumping lemma}
\vspace{2pt}\hrule height 0.3pt\vspace{3pt}
\textbf{Guide / Hint:} \\
{\small If $L$ and $\overline{L}$ are both recursively enumerable, $L$ is recursive by Post's theorem (B is TRUE). Every CFL is recursive, and recursive languages are closed under complement, so $\overline{\text{CFL}}$ is recursive (C is TRUE). Regular languages are closed under intersection, and every regular language is a DCFL (D is TRUE). Subsets of RE languages need not be recursive (A is FALSE).}

\end{tcolorbox}
\vspace{6pt}\hrule height 0.35pt\vspace{6pt}
\noindent
\begin{tcolorbox}[
    blanker,
    sidebyside,
    sidebyside align=top,
    lefthand width=0.55\textwidth,
    top=4pt, bottom=4pt, left=0pt, right=0pt,
    middle=16pt,
    underlay={
        \draw[dashed, black!80, line width=0.8pt, dash pattern=on 4pt off 3pt]
            ([xshift=0.55\textwidth+8pt]frame.north west) --
            ([xshift=0.55\textwidth+8pt]frame.south west);
    }
]
\textbf{Q.51} \quad {\footnotesize\textsc{[MSQ $\cdot$ 2 Mark(s) $\cdot$ No Neg]}} \hfill {\footnotesize\textbf{GATE 2022 CS} \quad \tiny\texttt{2022\_CS\_CS\_Q46}}
\vspace{2pt}\hrule height 0.3pt\vspace{4pt}

Which of the following is/are undecidable?

\par\vspace{6pt}\noindent
\begin{tabular}{@{}p{0.98\linewidth}@{}}
$\square$\ \textbf{(A)} Given two Turing machines $M_1$ and $M_2$, decide if $L(M_1) = L(M_2)$. \\
\addlinespace[3pt]
$\square$\ \textbf{(B)} Given a Turing machine $M$, decide if $L(M)$ is regular. \\
\addlinespace[3pt]
$\square$\ \textbf{(C)} Given a Turing machine $M$, decide if $M$ accepts all strings. \\
\addlinespace[3pt]
$\square$\ \textbf{(D)} Given a Turing machine $M$, decide if $M$ takes more than $10^{73}$ steps on every string.
\end{tabular}

\tcblower

\textbf{Official Key:} \textbf{\large A,B,C} \hfill \textit{\small [Medium]}
\vspace{2pt}\hrule height 0.3pt\vspace{3pt}
\textbf{Topic:} Theory of Computation \\
\textbf{Subtopic:} \textit{Turing machines and undecidability}
\vspace{2pt}\hrule height 0.3pt\vspace{3pt}
\textbf{Guide / Hint:} \\
{\small By Rice's theorem, any non-trivial semantic property of Turing machine languages is undecidable, making equivalence $L(M_1) = L(M_2)$ (A), regularity of $L(M)$ (B), and totality $L(M) = \Sigma^*$ (C) undecidable. Deciding if $M$ takes $> 10^{73}$ steps is bounded and decidable by simulation (D is decidable).}

\end{tcolorbox}
\vspace{6pt}\hrule height 0.35pt\vspace{6pt}
\noindent
\begin{tcolorbox}[
    blanker,
    sidebyside,
    sidebyside align=top,
    lefthand width=0.55\textwidth,
    top=4pt, bottom=4pt, left=0pt, right=0pt,
    middle=16pt,
    underlay={
        \draw[dashed, black!80, line width=0.8pt, dash pattern=on 4pt off 3pt]
            ([xshift=0.55\textwidth+8pt]frame.north west) --
            ([xshift=0.55\textwidth+8pt]frame.south west);
    }
]
\textbf{Q.52} \quad {\footnotesize\textsc{[MSQ $\cdot$ 2 Mark(s) $\cdot$ No Neg]}} \hfill {\footnotesize\textbf{GATE 2022 CS} \quad \tiny\texttt{2022\_CS\_CS\_Q47}}
\vspace{2pt}\hrule height 0.3pt\vspace{4pt}

Consider the following languages:
\[ L\_1 = \{a^n w a^n \mid w \in \{a, b\}^*\} \]
\[ L\_2 = \{w x w^R \mid w, x \in \{a, b\}^*, |w|, |x| > 0\} \]
Note that $w^R$ is the reversal of the string $w$. Which of the following is/are TRUE?

\par\vspace{6pt}\noindent
\begin{tabular}{@{}p{0.98\linewidth}@{}}
$\square$\ \textbf{(A)} $L_1$ and $L_2$ are regular. \\
\addlinespace[3pt]
$\square$\ \textbf{(B)} $L_1$ and $L_2$ are context-free. \\
\addlinespace[3pt]
$\square$\ \textbf{(C)} $L_1$ is regular and $L_2$ is context-free. \\
\addlinespace[3pt]
$\square$\ \textbf{(D)} $L_1$ and $L_2$ are context-free but not regular.
\end{tabular}

\tcblower

\textbf{Official Key:} \textbf{\large A,B,C} \hfill \textit{\small [Medium]}
\vspace{2pt}\hrule height 0.3pt\vspace{3pt}
\textbf{Topic:} Theory of Computation \\
\textbf{Subtopic:} \textit{Regular and context-free languages, pumping lemma}
\vspace{2pt}\hrule height 0.3pt\vspace{3pt}
\textbf{Guide / Hint:} \\
{\small Setting $n=1$ in $L_1 = \{a^n w a^n\}$ shows that any string starting and ending with $a$ of length $\ge 2$ belongs to $L_1$, making $L_1 = a(a+b)^*a$ regular. For $L_2 = \{w x w^R \mid |w|, |x| > 0\}$, taking $|w|=1$ yields $a(a+b)^+ a + b(a+b)^+ b$, which is regular. Every regular language is context-free, so statements A, B, and C are all TRUE.}

\end{tcolorbox}
\vspace{6pt}\hrule height 0.35pt\vspace{6pt}
\noindent
\begin{tcolorbox}[
    blanker,
    sidebyside,
    sidebyside align=top,
    lefthand width=0.55\textwidth,
    top=4pt, bottom=4pt, left=0pt, right=0pt,
    middle=16pt,
    underlay={
        \draw[dashed, black!80, line width=0.8pt, dash pattern=on 4pt off 3pt]
            ([xshift=0.55\textwidth+8pt]frame.north west) --
            ([xshift=0.55\textwidth+8pt]frame.south west);
    }
]
\textbf{Q.53} \quad {\footnotesize\textsc{[MSQ $\cdot$ 2 Mark(s) $\cdot$ No Neg]}} \hfill {\footnotesize\textbf{GATE 2022 CS} \quad \tiny\texttt{2022\_CS\_CS\_Q48}}
\vspace{2pt}\hrule height 0.3pt\vspace{4pt}

Consider the following languages:
\[ L\_1 = \{w w \mid w \in \{a, b\}^*\} \]
\[ L\_2 = \{a^n b^n c^m \mid m, n \ge 0\} \]
\[ L\_3 = \{a^m b^n c^n \mid m, n \ge 0\} \]
Which of the following statements is/are FALSE?

\par\vspace{6pt}\noindent
\begin{tabular}{@{}p{0.98\linewidth}@{}}
$\square$\ \textbf{(A)} $L_1$ is not context-free but $L_2$ and $L_3$ are deterministic context-free. \\
\addlinespace[3pt]
$\square$\ \textbf{(B)} Neither $L_1$ nor $L_2$ is context-free. \\
\addlinespace[3pt]
$\square$\ \textbf{(C)} $L_2$, $L_3$ and $L_2 \cap L_3$ all are context-free. \\
\addlinespace[3pt]
$\square$\ \textbf{(D)} Neither $L_1$ nor its complement is context-free.
\end{tabular}

\tcblower

\textbf{Official Key:} \textbf{\large B,C,D} \hfill \textit{\small [Hard]}
\vspace{2pt}\hrule height 0.3pt\vspace{3pt}
\textbf{Topic:} Theory of Computation \\
\textbf{Subtopic:} \textit{Context-free grammars and push-down automata}
\vspace{2pt}\hrule height 0.3pt\vspace{3pt}
\textbf{Guide / Hint:} \\
{\small Language $L_2 = \{a^n b^n c^m\}$ is a DCFL, so stating that neither $L_1$ nor $L_2$ is CFL is false (B is FALSE). $L_2 \cap L_3 = \{a^n b^n c^n\}$ is not context-free, so claiming all are context-free is false (C is FALSE). The complement $\overline{L_1}$ is a well-known context-free language, so D is FALSE. Thus, statements B, C, and D are FALSE.}

\end{tcolorbox}
\vspace{6pt}\hrule height 0.35pt\vspace{6pt}
\noindent
\begin{tcolorbox}[
    blanker,
    sidebyside,
    sidebyside align=top,
    lefthand width=0.55\textwidth,
    top=4pt, bottom=4pt, left=0pt, right=0pt,
    middle=16pt,
    underlay={
        \draw[dashed, black!80, line width=0.8pt, dash pattern=on 4pt off 3pt]
            ([xshift=0.55\textwidth+8pt]frame.north west) --
            ([xshift=0.55\textwidth+8pt]frame.south west);
    }
]
\textbf{Q.54} \quad {\footnotesize\textsc{[MSQ $\cdot$ 1 Mark(s) $\cdot$ No Neg]}} \hfill {\footnotesize\textbf{GATE 2023 CS} \quad \tiny\texttt{2023\_CS\_CS\_Q24}}
\vspace{2pt}\hrule height 0.3pt\vspace{4pt}

Which of the following statements is/are CORRECT?

\par\vspace{6pt}\noindent
\begin{tabular}{@{}p{0.98\linewidth}@{}}
$\square$\ \textbf{(A)} The intersection of two regular languages is regular. \\
\addlinespace[3pt]
$\square$\ \textbf{(B)} The intersection of two context-free languages is context-free. \\
\addlinespace[3pt]
$\square$\ \textbf{(C)} The intersection of two recursive languages is recursive. \\
\addlinespace[3pt]
$\square$\ \textbf{(D)} The intersection of two recursively enumerable languages is recursively enumerable.
\end{tabular}

\tcblower

\textbf{Official Key:} \textbf{\large A, C, D} \hfill \textit{\small [Easy]}
\vspace{2pt}\hrule height 0.3pt\vspace{3pt}
\textbf{Topic:} Theory of Computation \\
\textbf{Subtopic:} \textit{Regular and context-free languages, pumping lemma}
\vspace{2pt}\hrule height 0.3pt\vspace{3pt}
\textbf{Guide / Hint:} \\
{\small Regular languages, recursive languages, and recursively enumerable languages are all closed under intersection (A, C, and D are CORRECT). Context-free languages are not closed under intersection; for example, $L_1 = \{a^n b^n c^m\}$ and $L_2 = \{a^m b^n c^n\}$ are both CFLs, but $L_1 \cap L_2 = \{a^n b^n c^n\}$ is not context-free.}

\end{tcolorbox}
\vspace{6pt}\hrule height 0.35pt\vspace{6pt}
\noindent
\begin{tcolorbox}[
    blanker,
    sidebyside,
    sidebyside align=top,
    lefthand width=0.55\textwidth,
    top=4pt, bottom=4pt, left=0pt, right=0pt,
    middle=16pt,
    underlay={
        \draw[dashed, black!80, line width=0.8pt, dash pattern=on 4pt off 3pt]
            ([xshift=0.55\textwidth+8pt]frame.north west) --
            ([xshift=0.55\textwidth+8pt]frame.south west);
    }
]
\textbf{Q.55} \quad {\footnotesize\textsc{[MSQ $\cdot$ 1 Mark(s) $\cdot$ No Neg]}} \hfill {\footnotesize\textbf{GATE 2024 CS1} \quad \tiny\texttt{2024\_CS1\_CS\_Q23}}
\vspace{2pt}\hrule height 0.3pt\vspace{4pt}

Let $L_1, L_2$ be two regular languages and $L_3$ a language which is not regular. Which of the following statements is/are always TRUE?

\par\vspace{6pt}\noindent
\begin{tabular}{@{}p{0.98\linewidth}@{}}
$\square$\ \textbf{(A)} $L_1 = L_2$ if and only if $L_1 \cap \overline{L_2} = \emptyset$ \\
\addlinespace[3pt]
$\square$\ \textbf{(B)} $L_1 \cup L_3$ is not regular \\
\addlinespace[3pt]
$\square$\ \textbf{(C)} $\overline{L_3}$ is not regular \\
\addlinespace[3pt]
$\square$\ \textbf{(D)} $\overline{L_1} \cup \overline{L_2}$ is regular
\end{tabular}

\tcblower

\textbf{Official Key:} \textbf{\large C;D} \hfill \textit{\small [Medium]}
\vspace{2pt}\hrule height 0.3pt\vspace{3pt}
\textbf{Topic:} Theory of Computation \\
\textbf{Subtopic:} \textit{Regular and context-free languages, pumping lemma}
\vspace{2pt}\hrule height 0.3pt\vspace{3pt}
\textbf{Guide / Hint:} \\
{\small Regular languages are closed under complementation and union; since $L_1, L_2$ are regular, $\overline{L_1}$ and $\overline{L_2}$ are regular, so $\overline{L_1} \cup \overline{L_2}$ is regular (D is TRUE). If $L_3$ is non-regular, its complement $\overline{L_3}$ must also be non-regular, because if $\overline{L_3}$ were regular, its complement $\overline{\overline{L_3}} = L_3$ would be regular, a contradiction (C is TRUE).}

\end{tcolorbox}
\vspace{6pt}\hrule height 0.35pt\vspace{6pt}
\noindent
\begin{tcolorbox}[
    blanker,
    sidebyside,
    sidebyside align=top,
    lefthand width=0.55\textwidth,
    top=4pt, bottom=4pt, left=0pt, right=0pt,
    middle=16pt,
    underlay={
        \draw[dashed, black!80, line width=0.8pt, dash pattern=on 4pt off 3pt]
            ([xshift=0.55\textwidth+8pt]frame.north west) --
            ([xshift=0.55\textwidth+8pt]frame.south west);
    }
]
\textbf{Q.56} \quad {\footnotesize\textsc{[MSQ $\cdot$ 2 Mark(s) $\cdot$ No Neg]}} \hfill {\footnotesize\textbf{GATE 2024 CS1} \quad \tiny\texttt{2024\_CS1\_CS\_Q50}}
\vspace{2pt}\hrule height 0.3pt\vspace{4pt}

Consider the 5-state DFA $M$ accepting the language $L(M) \subset (0 + 1)^*$ shown below. For any string $w \in (0 + 1)^*$ let $n_0(w)$ be the number of 0's in $w$ and $n_1(w)$ be the number of 1's in $w$.
Which of the following statements is/are FALSE?

\par\vspace{6pt}\noindent
\begin{tabular}{@{}p{0.98\linewidth}@{}}
$\square$\ \textbf{(A)} States 2 and 4 are distinguishable in $M$ \\
\addlinespace[3pt]
$\square$\ \textbf{(B)} States 3 and 4 are distinguishable in $M$ \\
\addlinespace[3pt]
$\square$\ \textbf{(C)} States 2 and 5 are distinguishable in $M$ \\
\addlinespace[3pt]
$\square$\ \textbf{(D)} Any string $w$ with $n_0(w) = n_1(w)$ is in $L(M)$
\end{tabular}

\tcblower

\textbf{Official Key:} \textbf{\large B;C} \hfill \textit{\small [Medium]}
\vspace{2pt}\hrule height 0.3pt\vspace{3pt}
\textbf{Topic:} Theory of Computation \\
\textbf{Subtopic:} \textit{Regular expressions and finite automata}
\vspace{2pt}\hrule height 0.3pt\vspace{3pt}
\textbf{Guide / Hint:} \\
{\small States $p$ and $q$ are distinguishable if there exists a string $w$ such that exactly one of $\delta(p, w)$ or $\delta(q, w)$ is an accepting state. In DFA $M$, states 3 and 4 have identical transition behavior and both are non-accepting, making them equivalent/indistinguishable, so claiming they are distinguishable is FALSE (B is FALSE). Similarly, states 2 and 5 are equivalent, so claiming they are distinguishable is FALSE (C is FALSE).}

\end{tcolorbox}
\vspace{6pt}\hrule height 0.35pt\vspace{6pt}
\noindent
\begin{tcolorbox}[
    blanker,
    sidebyside,
    sidebyside align=top,
    lefthand width=0.55\textwidth,
    top=4pt, bottom=4pt, left=0pt, right=0pt,
    middle=16pt,
    underlay={
        \draw[dashed, black!80, line width=0.8pt, dash pattern=on 4pt off 3pt]
            ([xshift=0.55\textwidth+8pt]frame.north west) --
            ([xshift=0.55\textwidth+8pt]frame.south west);
    }
]
\textbf{Q.57} \quad {\footnotesize\textsc{[MSQ $\cdot$ 2 Mark(s) $\cdot$ No Neg]}} \hfill {\footnotesize\textbf{GATE 2024 CS2} \quad \tiny\texttt{2024\_CS2\_CS\_Q52}}
\vspace{2pt}\hrule height 0.3pt\vspace{4pt}

Consider a context-free grammar $G$ with the following 3 rules:
\[ S \to aS \mid aSbS \mid c \]
Let $w \in L(G)$. Let $n_a(w), n_b(w), n_c(w)$ denote the number of times $a, b, c$ occur in $w$, respectively. Which of the following statements is/are TRUE?

\par\vspace{6pt}\noindent
\begin{tabular}{@{}p{0.48\linewidth}@{\hspace{0.04\linewidth}}p{0.48\linewidth}@{}}
$\square$\ \textbf{(A)} $n_a(w) > n_b(w)$ & $\square$\ \textbf{(B)} $n_a(w) > n_c(w) - 2$ \\
\addlinespace[3pt]
$\square$\ \textbf{(C)} $n_c(w) = n_b(w) + 1$ & $\square$\ \textbf{(D)} $n_c(w) = n_b(w) * 2$
\end{tabular}

\tcblower

\textbf{Official Key:} \textbf{\large B;C} \hfill \textit{\small [Medium]}
\vspace{2pt}\hrule height 0.3pt\vspace{3pt}
\textbf{Topic:} Theory of Computation \\
\textbf{Subtopic:} \textit{Context-free grammars and push-down automata}
\vspace{2pt}\hrule height 0.3pt\vspace{3pt}
\textbf{Guide / Hint:} \\
{\small Each application of $S \to aSbS$ consumes one $S$ and introduces two new $S$'s along with one $b$, effectively adding $+1$ to the count of $S$'s. Each termination $S \to c$ eliminates one $S$ and adds one $c$. Thus, the total number of $c$'s must equal the number of $b$'s plus 1: $n_c(w) = n_b(w) + 1$ (C is true). Also, every $b$ is accompanied by an $a$, and additional $a$'s can come from $S \to aS$, giving $n_a(w) \ge n_b(w) = n_c(w) - 1 > n_c(w) - 2$ (B is true).}

\end{tcolorbox}
\vspace{6pt}\hrule height 0.35pt\vspace{6pt}
\noindent
\begin{tcolorbox}[
    blanker,
    sidebyside,
    sidebyside align=top,
    lefthand width=0.55\textwidth,
    top=4pt, bottom=4pt, left=0pt, right=0pt,
    middle=16pt,
    underlay={
        \draw[dashed, black!80, line width=0.8pt, dash pattern=on 4pt off 3pt]
            ([xshift=0.55\textwidth+8pt]frame.north west) --
            ([xshift=0.55\textwidth+8pt]frame.south west);
    }
]
\textbf{Q.58} \quad {\footnotesize\textsc{[MSQ $\cdot$ 1 Mark(s) $\cdot$ No Neg]}} \hfill {\footnotesize\textbf{GATE 2025 CS1} \quad \tiny\texttt{2025\_CS1\_CS\_Q28}}
\vspace{2pt}\hrule height 0.3pt\vspace{4pt}

A regular language $L$ is accepted by a non-deterministic finite automaton (NFA) with $n$ states. Which of the following statement(s) is/are FALSE?

\par\vspace{6pt}\noindent
\begin{tabular}{@{}p{0.98\linewidth}@{}}
$\square$\ \textbf{(A)} $L$ may have an accepting NFA with $< n$ states. \\
\addlinespace[3pt]
$\square$\ \textbf{(B)} $L$ may have an accepting DFA with $< n$ states. \\
\addlinespace[3pt]
$\square$\ \textbf{(C)} There exists a DFA with $\le 2^n$ states that accepts $L$. \\
\addlinespace[3pt]
$\square$\ \textbf{(D)} Every DFA that accepts $L$ has $> 2^n$ states.
\end{tabular}

\tcblower

\textbf{Official Key:} \textbf{\large D} \hfill \textit{\small [Easy]}
\vspace{2pt}\hrule height 0.3pt\vspace{3pt}
\textbf{Topic:} Theory of Computation \\
\textbf{Subtopic:} \textit{Regular expressions and finite automata}
\vspace{2pt}\hrule height 0.3pt\vspace{3pt}
\textbf{Guide / Hint:} \\
{\small By the standard subset construction algorithm, every $n$-state NFA can be converted into an equivalent deterministic finite automaton (DFA) having at most $2^n$ states. Statement D asserts that every DFA must have strictly more than $2^n$ states, which is directly false.}

\end{tcolorbox}
\vspace{6pt}\hrule height 0.35pt\vspace{6pt}
\noindent
\begin{tcolorbox}[
    blanker,
    sidebyside,
    sidebyside align=top,
    lefthand width=0.55\textwidth,
    top=4pt, bottom=4pt, left=0pt, right=0pt,
    middle=16pt,
    underlay={
        \draw[dashed, black!80, line width=0.8pt, dash pattern=on 4pt off 3pt]
            ([xshift=0.55\textwidth+8pt]frame.north west) --
            ([xshift=0.55\textwidth+8pt]frame.south west);
    }
]
\textbf{Q.59} \quad {\footnotesize\textsc{[MSQ $\cdot$ 2 Mark(s) $\cdot$ No Neg]}} \hfill {\footnotesize\textbf{GATE 2025 CS1} \quad \tiny\texttt{2025\_CS1\_CS\_Q50}}
\vspace{2pt}\hrule height 0.3pt\vspace{4pt}

Consider the deterministic finite automaton (DFA) defined over the alphabet $\Sigma = \{a, b\}$. Identify which of the following language(s) is/are accepted by the given DFA.

\par\vspace{6pt}\noindent
\begin{tabular}{@{}p{0.98\linewidth}@{}}
$\square$\ \textbf{(A)} The set of all strings containing an even number of $b$'s. \\
\addlinespace[3pt]
$\square$\ \textbf{(B)} The set of all strings containing the pattern $bab$. \\
\addlinespace[3pt]
$\square$\ \textbf{(C)} The set of all strings ending with the pattern $bab$. \\
\addlinespace[3pt]
$\square$\ \textbf{(D)} The set of all strings not containing the pattern $aba$.
\end{tabular}

\tcblower

\textbf{Official Key:} \textbf{\large C} \hfill \textit{\small [Medium]}
\vspace{2pt}\hrule height 0.3pt\vspace{3pt}
\textbf{Topic:} Theory of Computation \\
\textbf{Subtopic:} \textit{Regular expressions and finite automata}
\vspace{2pt}\hrule height 0.3pt\vspace{3pt}
\textbf{Guide / Hint:} \\
{\small In the given DFA, the sequence of transitions leading to the unique accepting state $q_3$ corresponds precisely to recognizing the suffix $bab$. Any character following $q_3$ transitions back according to pattern matching prefixes of $bab$, accepting exactly strings ending with the pattern $bab$.}

\end{tcolorbox}
\vspace{6pt}\hrule height 0.35pt\vspace{6pt}
\noindent
\begin{tcolorbox}[
    blanker,
    sidebyside,
    sidebyside align=top,
    lefthand width=0.55\textwidth,
    top=4pt, bottom=4pt, left=0pt, right=0pt,
    middle=16pt,
    underlay={
        \draw[dashed, black!80, line width=0.8pt, dash pattern=on 4pt off 3pt]
            ([xshift=0.55\textwidth+8pt]frame.north west) --
            ([xshift=0.55\textwidth+8pt]frame.south west);
    }
]
\textbf{Q.60} \quad {\footnotesize\textsc{[MSQ $\cdot$ 1 Mark(s) $\cdot$ No Neg]}} \hfill {\footnotesize\textbf{GATE 2025 CS2} \quad \tiny\texttt{2025\_CS2\_CS\_Q30}}
\vspace{2pt}\hrule height 0.3pt\vspace{4pt}

Consider the two lists List I and List II given below:
\par\medskip
\begin{center}\small
\begin{tabular}{ll}
\toprule
List I & List II \\
\midrule
(i) Context free languages & (a) Closed under union \\
(ii) Recursive languages & (b) Not closed under complementation \\
(iii) Regular languages & (c) Closed under intersection \\
\bottomrule
\end{tabular}\end{center}
\par\medskip
For matching of items in List I with those in List II, which of the following option(s) is/are CORRECT?

\par\vspace{6pt}\noindent
\begin{tabular}{@{}p{0.98\linewidth}@{}}
$\square$\ \textbf{(A)} (i) – (a), (ii) – (b), and (iii) – (c) \\
\addlinespace[3pt]
$\square$\ \textbf{(B)} (i) – (b), (ii) – (a), and (iii) – (c) \\
\addlinespace[3pt]
$\square$\ \textbf{(C)} (i) – (b), (ii) – (c), and (iii) – (a) \\
\addlinespace[3pt]
$\square$\ \textbf{(D)} (i) – (a), (ii) – (c), and (iii) – (b)
\end{tabular}

\tcblower

\textbf{Official Key:} \textbf{\large B;C} \hfill \textit{\small [Medium]}
\vspace{2pt}\hrule height 0.3pt\vspace{3pt}
\textbf{Topic:} Theory of Computation \\
\textbf{Subtopic:} \textit{Regular and context-free languages: closure properties}
\vspace{2pt}\hrule height 0.3pt\vspace{3pt}
\textbf{Guide / Hint:} \\
{\small Context-free languages are not closed under complementation (i - b) or intersection, but are closed under union. Recursive languages are closed under union and complementation and intersection (ii - a or c). Regular languages are closed under all three.}

\end{tcolorbox}
\vspace{6pt}\hrule height 0.35pt\vspace{6pt}
\noindent
\begin{tcolorbox}[
    blanker,
    sidebyside,
    sidebyside align=top,
    lefthand width=0.55\textwidth,
    top=4pt, bottom=4pt, left=0pt, right=0pt,
    middle=16pt,
    underlay={
        \draw[dashed, black!80, line width=0.8pt, dash pattern=on 4pt off 3pt]
            ([xshift=0.55\textwidth+8pt]frame.north west) --
            ([xshift=0.55\textwidth+8pt]frame.south west);
    }
]
\textbf{Q.61} \quad {\footnotesize\textsc{[MSQ $\cdot$ 2 Mark(s) $\cdot$ No Neg]}} \hfill {\footnotesize\textbf{GATE 2025 CS2} \quad \tiny\texttt{2025\_CS2\_CS\_Q52}}
\vspace{2pt}\hrule height 0.3pt\vspace{4pt}

Let $\Sigma = \{a, b, c\}$. For $x \in \Sigma^*$, let $\#_\alpha(x)$ denote the number of occurrences of $\alpha$ in $x$. Which one or more of the following languages is/are regular?

\par\vspace{6pt}\noindent
\begin{tabular}{@{}p{0.98\linewidth}@{}}
$\square$\ \textbf{(A)} $\{a^m b^n \mid m, n \ge 0\}$ \\
\addlinespace[3pt]
$\square$\ \textbf{(B)} $\{a, b\}^* \cap \{a^m b^n c^{m-n} \mid m \ge n \ge 0\}$ \\
\addlinespace[3pt]
$\square$\ \textbf{(C)} $\{w \in \{a, b\}^* \mid \#_a(w) \equiv 2 \pmod 7 \text{ and } \#_b(w) \equiv 3 \pmod 9\}$ \\
\addlinespace[3pt]
$\square$\ \textbf{(D)} $\{w \in \{a, b\}^* \mid \#_a(w) \equiv 2 \pmod 7 \text{ and } \#_a(w) = \#_b(w)\}$
\end{tabular}

\tcblower

\textbf{Official Key:} \textbf{\large A;C} \hfill \textit{\small [Medium]}
\vspace{2pt}\hrule height 0.3pt\vspace{3pt}
\textbf{Topic:} Theory of Computation \\
\textbf{Subtopic:} \textit{Regular expressions and finite automata: regular languages}
\vspace{2pt}\hrule height 0.3pt\vspace{3pt}
\textbf{Guide / Hint:} \\
{\small Language A ($a^m b^n$) is generated by the regular expression $a^* b^*$. Language C counts occurrences modulo 7 and 9, which requires finite memory ($7 \times 9 = 63$ states) and is therefore regular.}

\end{tcolorbox}
\vspace{6pt}\hrule height 0.35pt\vspace{6pt}
\vspace{8pt}
\subsection*{Compiler Design}
\vspace{-4pt}
\noindent
\begin{tcolorbox}[
    blanker,
    sidebyside,
    sidebyside align=top,
    lefthand width=0.55\textwidth,
    top=4pt, bottom=4pt, left=0pt, right=0pt,
    middle=16pt,
    underlay={
        \draw[dashed, black!80, line width=0.8pt, dash pattern=on 4pt off 3pt]
            ([xshift=0.55\textwidth+8pt]frame.north west) --
            ([xshift=0.55\textwidth+8pt]frame.south west);
    }
]
\textbf{Q.62} \quad {\footnotesize\textsc{[MSQ $\cdot$ 1 Mark(s) $\cdot$ No Neg]}} \hfill {\footnotesize\textbf{GATE 2024 CS1} \quad \tiny\texttt{2024\_CS1\_CS\_Q26}}
\vspace{2pt}\hrule height 0.3pt\vspace{4pt}

Which of the following is/are Bottom-Up Parser(s)?

\par\vspace{6pt}\noindent
\begin{tabular}{@{}p{0.48\linewidth}@{\hspace{0.04\linewidth}}p{0.48\linewidth}@{}}
$\square$\ \textbf{(A)} Shift-reduce Parser & $\square$\ \textbf{(B)} Predictive Parser \\
\addlinespace[3pt]
$\square$\ \textbf{(C)} LL(1) Parser & $\square$\ \textbf{(D)} LR Parser
\end{tabular}

\tcblower

\textbf{Official Key:} \textbf{\large A;D} \hfill \textit{\small [Easy]}
\vspace{2pt}\hrule height 0.3pt\vspace{3pt}
\textbf{Topic:} Compiler Design \\
\textbf{Subtopic:} \textit{Lexical analysis, parsing, syntax-directed translation}
\vspace{2pt}\hrule height 0.3pt\vspace{3pt}
\textbf{Guide / Hint:} \\
{\small Bottom-up parsers reduce an input string back to the start symbol by finding handles. Shift-reduce parsers and LR parsers (SLR, CLR, LALR) are bottom-up parsers. In contrast, Predictive and LL(1) parsers construct the derivation tree from the root downward and are top-down parsers.}

\end{tcolorbox}
\vspace{6pt}\hrule height 0.35pt\vspace{6pt}
\noindent
\begin{tcolorbox}[
    blanker,
    sidebyside,
    sidebyside align=top,
    lefthand width=0.55\textwidth,
    top=4pt, bottom=4pt, left=0pt, right=0pt,
    middle=16pt,
    underlay={
        \draw[dashed, black!80, line width=0.8pt, dash pattern=on 4pt off 3pt]
            ([xshift=0.55\textwidth+8pt]frame.north west) --
            ([xshift=0.55\textwidth+8pt]frame.south west);
    }
]
\textbf{Q.63} \quad {\footnotesize\textsc{[MSQ $\cdot$ 1 Mark(s) $\cdot$ No Neg]}} \hfill {\footnotesize\textbf{GATE 2024 CS2} \quad \tiny\texttt{2024\_CS2\_CS\_Q29}}
\vspace{2pt}\hrule height 0.3pt\vspace{4pt}

Which of the following statements is/are FALSE?

\par\vspace{6pt}\noindent
\begin{tabular}{@{}p{0.98\linewidth}@{}}
$\square$\ \textbf{(A)} An attribute grammar is a syntax-directed definition (SDD) in which the functions in the semantic rules have no side effects \\
\addlinespace[3pt]
$\square$\ \textbf{(B)} The attributes in a L-attributed definition cannot always be evaluated in a depth-first order \\
\addlinespace[3pt]
$\square$\ \textbf{(C)} Synthesized attributes can be evaluated by a bottom-up parser as the input is parsed \\
\addlinespace[3pt]
$\square$\ \textbf{(D)} All L-attributed definitions based on LR(1) grammar can be evaluated using a bottom-up parsing strategy
\end{tabular}

\tcblower

\textbf{Official Key:} \textbf{\large B;D} \hfill \textit{\small [Hard]}
\vspace{2pt}\hrule height 0.3pt\vspace{3pt}
\textbf{Topic:} Compiler Design \\
\textbf{Subtopic:} \textit{Lexical analysis, parsing, syntax-directed translation}
\vspace{2pt}\hrule height 0.3pt\vspace{3pt}
\textbf{Guide / Hint:} \\
{\small L-attributed definitions are explicitly defined such that their attributes can always be evaluated via depth-first (left-to-right) traversal (making statement B false). Not all L-attributed definitions on LR(1) grammars can be evaluated during bottom-up parsing without rewriting or marker nonterminals (making statement D false).}

\end{tcolorbox}
\vspace{6pt}\hrule height 0.35pt\vspace{6pt}
\noindent
\begin{tcolorbox}[
    blanker,
    sidebyside,
    sidebyside align=top,
    lefthand width=0.55\textwidth,
    top=4pt, bottom=4pt, left=0pt, right=0pt,
    middle=16pt,
    underlay={
        \draw[dashed, black!80, line width=0.8pt, dash pattern=on 4pt off 3pt]
            ([xshift=0.55\textwidth+8pt]frame.north west) --
            ([xshift=0.55\textwidth+8pt]frame.south west);
    }
]
\textbf{Q.64} \quad {\footnotesize\textsc{[MSQ $\cdot$ 2 Mark(s) $\cdot$ No Neg]}} \hfill {\footnotesize\textbf{GATE 2025 CS1} \quad \tiny\texttt{2025\_CS1\_CS\_Q46}}
\vspace{2pt}\hrule height 0.3pt\vspace{4pt}

Which of the following statement(s) is/are TRUE while computing First and Follow during top down parsing by a compiler?

\par\vspace{6pt}\noindent
\begin{tabular}{@{}p{0.98\linewidth}@{}}
$\square$\ \textbf{(A)} For a production $A \to \epsilon$, $\epsilon$ will be added to $First(A)$. \\
\addlinespace[3pt]
$\square$\ \textbf{(B)} If there is any input right end marker, it will be added to $First(S)$, where $S$ is the start symbol. \\
\addlinespace[3pt]
$\square$\ \textbf{(C)} For a production $A \to \epsilon$, $\epsilon$ will be added to $Follow(A)$. \\
\addlinespace[3pt]
$\square$\ \textbf{(D)} If there is any input right end marker, it will be added to $Follow(S)$, where $S$ is the start symbol.
\end{tabular}

\tcblower

\textbf{Official Key:} \textbf{\large A;D} \hfill \textit{\small [Easy]}
\vspace{2pt}\hrule height 0.3pt\vspace{3pt}
\textbf{Topic:} Compiler Design \\
\textbf{Subtopic:} \textit{Lexical analysis, parsing, syntax-directed translation}
\vspace{2pt}\hrule height 0.3pt\vspace{3pt}
\textbf{Guide / Hint:} \\
{\small By definition of FIRST and FOLLOW: if $A \to \epsilon$, then $\epsilon \in \text{FIRST}(A)$ (A is true). $\epsilon$ can never belong to any FOLLOW set (C is false). The right end marker (\$) is placed into $\text{FOLLOW}(S)$ for the start symbol $S$ to mark the end of input (D is true, B is false).}

\end{tcolorbox}
\vspace{6pt}\hrule height 0.35pt\vspace{6pt}
\noindent
\begin{tcolorbox}[
    blanker,
    sidebyside,
    sidebyside align=top,
    lefthand width=0.55\textwidth,
    top=4pt, bottom=4pt, left=0pt, right=0pt,
    middle=16pt,
    underlay={
        \draw[dashed, black!80, line width=0.8pt, dash pattern=on 4pt off 3pt]
            ([xshift=0.55\textwidth+8pt]frame.north west) --
            ([xshift=0.55\textwidth+8pt]frame.south west);
    }
]
\textbf{Q.65} \quad {\footnotesize\textsc{[MSQ $\cdot$ 2 Mark(s) $\cdot$ No Neg]}} \hfill {\footnotesize\textbf{GATE 2025 CS2} \quad \tiny\texttt{2025\_CS2\_CS\_Q51}}
\vspace{2pt}\hrule height 0.3pt\vspace{4pt}

Consider two grammars $G_1$ and $G_2$ with production rules for dangling-else conditional statements. Which of the following statements is/are CORRECT?

\par\vspace{6pt}\noindent
\begin{tabular}{@{}p{0.98\linewidth}@{}}
$\square$\ \textbf{(A)} $G_1$ is not $LL(1)$ and $G_2$ is $LL(1)$. \\
\addlinespace[3pt]
$\square$\ \textbf{(B)} $G_1$ is $LL(1)$ and $G_2$ is not $LL(1)$. \\
\addlinespace[3pt]
$\square$\ \textbf{(C)} $G_1$ and $G_2$ are not $LL(1)$. \\
\addlinespace[3pt]
$\square$\ \textbf{(D)} $G_1$ and $G_2$ are ambiguous.
\end{tabular}

\tcblower

\textbf{Official Key:} \textbf{\large C;D} \hfill \textit{\small [Medium]}
\vspace{2pt}\hrule height 0.3pt\vspace{3pt}
\textbf{Topic:} Compiler Design \\
\textbf{Subtopic:} \textit{Parsing: LL(1) grammars and ambiguity}
\vspace{2pt}\hrule height 0.3pt\vspace{3pt}
\textbf{Guide / Hint:} \\
{\small Both $G_1$ and $G_2$ exhibit the classic dangling-else ambiguity, generating two distinct parse trees for nested if-then-else statements. Since ambiguous grammars can never be $LL(1)$, both C and D are correct.}

\end{tcolorbox}
\vspace{6pt}\hrule height 0.35pt\vspace{6pt}
\vspace{8pt}
\subsection*{Operating Systems}
\vspace{-4pt}
\noindent
\begin{tcolorbox}[
    blanker,
    sidebyside,
    sidebyside align=top,
    lefthand width=0.55\textwidth,
    top=4pt, bottom=4pt, left=0pt, right=0pt,
    middle=16pt,
    underlay={
        \draw[dashed, black!80, line width=0.8pt, dash pattern=on 4pt off 3pt]
            ([xshift=0.55\textwidth+8pt]frame.north west) --
            ([xshift=0.55\textwidth+8pt]frame.south west);
    }
]
\textbf{Q.66} \quad {\footnotesize\textsc{[MSQ $\cdot$ 1 Mark(s) $\cdot$ No Neg]}} \hfill {\footnotesize\textbf{GATE 2022 CS} \quad \tiny\texttt{2022\_CS\_CS\_Q26}}
\vspace{2pt}\hrule height 0.3pt\vspace{4pt}

Which of the following statements is/are TRUE with respect to deadlocks?

\par\vspace{6pt}\noindent
\begin{tabular}{@{}p{0.98\linewidth}@{}}
$\square$\ \textbf{(A)} Circular wait is a necessary condition for the formation of deadlock. \\
\addlinespace[3pt]
$\square$\ \textbf{(B)} In a system where each resource has more than one instance, a cycle in its wait-for graph indicates the presence of a deadlock. \\
\addlinespace[3pt]
$\square$\ \textbf{(C)} If the current allocation of resources to processes leads the system to unsafe state, then deadlock will necessarily occur. \\
\addlinespace[3pt]
$\square$\ \textbf{(D)} In the resource-allocation graph of a system, if every edge is an assignment edge, then the system is not in deadlock state.
\end{tabular}

\tcblower

\textbf{Official Key:} \textbf{\large A,D} \hfill \textit{\small [Easy]}
\vspace{2pt}\hrule height 0.3pt\vspace{3pt}
\textbf{Topic:} Operating System \\
\textbf{Subtopic:} \textit{Deadlock}
\vspace{2pt}\hrule height 0.3pt\vspace{3pt}
\textbf{Guide / Hint:} \\
{\small Circular wait is one of Coffman's four necessary conditions for deadlock (A is TRUE). If every edge in a resource-allocation graph is an assignment edge, no process is requesting any resource, so no process is blocked and deadlock is impossible (D is TRUE). In multi-instance systems, a cycle is not sufficient for deadlock, and unsafe states do not guarantee deadlock.}

\end{tcolorbox}
\vspace{6pt}\hrule height 0.35pt\vspace{6pt}
\noindent
\begin{tcolorbox}[
    blanker,
    sidebyside,
    sidebyside align=top,
    lefthand width=0.55\textwidth,
    top=4pt, bottom=4pt, left=0pt, right=0pt,
    middle=16pt,
    underlay={
        \draw[dashed, black!80, line width=0.8pt, dash pattern=on 4pt off 3pt]
            ([xshift=0.55\textwidth+8pt]frame.north west) --
            ([xshift=0.55\textwidth+8pt]frame.south west);
    }
]
\textbf{Q.67} \quad {\footnotesize\textsc{[MSQ $\cdot$ 1 Mark(s) $\cdot$ No Neg]}} \hfill {\footnotesize\textbf{GATE 2023 CS} \quad \tiny\texttt{2023\_CS\_CS\_Q22}}
\vspace{2pt}\hrule height 0.3pt\vspace{4pt}

Which one or more of the following need to be saved on a context switch from one thread (T1) of a process to another thread (T2) of the same process?

\par\vspace{6pt}\noindent
\begin{tabular}{@{}p{0.48\linewidth}@{\hspace{0.04\linewidth}}p{0.48\linewidth}@{}}
$\square$\ \textbf{(A)} Page table base register & $\square$\ \textbf{(B)} Stack pointer \\
\addlinespace[3pt]
$\square$\ \textbf{(C)} Program counter & $\square$\ \textbf{(D)} General purpose registers
\end{tabular}

\tcblower

\textbf{Official Key:} \textbf{\large B, C, D} \hfill \textit{\small [Easy]}
\vspace{2pt}\hrule height 0.3pt\vspace{3pt}
\textbf{Topic:} Operating System \\
\textbf{Subtopic:} \textit{System calls, processes, threads, inter-process communication}
\vspace{2pt}\hrule height 0.3pt\vspace{3pt}
\textbf{Guide / Hint:} \\
{\small Threads of the same process share the process address space, open files, and the page table base register (PTBR is not switched). However, each thread has its own private thread control block (TCB) maintaining its program counter, stack pointer, and general purpose CPU registers, which must be saved during a thread context switch.}

\end{tcolorbox}
\vspace{6pt}\hrule height 0.35pt\vspace{6pt}
\noindent
\begin{tcolorbox}[
    blanker,
    sidebyside,
    sidebyside align=top,
    lefthand width=0.55\textwidth,
    top=4pt, bottom=4pt, left=0pt, right=0pt,
    middle=16pt,
    underlay={
        \draw[dashed, black!80, line width=0.8pt, dash pattern=on 4pt off 3pt]
            ([xshift=0.55\textwidth+8pt]frame.north west) --
            ([xshift=0.55\textwidth+8pt]frame.south west);
    }
]
\textbf{Q.68} \quad {\footnotesize\textsc{[MSQ $\cdot$ 1 Mark(s) $\cdot$ No Neg]}} \hfill {\footnotesize\textbf{GATE 2023 CS} \quad \tiny\texttt{2023\_CS\_CS\_Q23}}
\vspace{2pt}\hrule height 0.3pt\vspace{4pt}

Which one or more of the following options guarantee that a computer system will transition from user mode to kernel mode?

\par\vspace{6pt}\noindent
\begin{tabular}{@{}p{0.48\linewidth}@{\hspace{0.04\linewidth}}p{0.48\linewidth}@{}}
$\square$\ \textbf{(A)} Function Call & $\square$\ \textbf{(B)} `malloc` Call \\
\addlinespace[3pt]
$\square$\ \textbf{(C)} Page Fault & $\square$\ \textbf{(D)} System Call
\end{tabular}

\tcblower

\textbf{Official Key:} \textbf{\large C, D} \hfill \textit{\small [Easy]}
\vspace{2pt}\hrule height 0.3pt\vspace{3pt}
\textbf{Topic:} Operating System \\
\textbf{Subtopic:} \textit{System calls, processes, threads, inter-process communication}
\vspace{2pt}\hrule height 0.3pt\vspace{3pt}
\textbf{Guide / Hint:} \\
{\small A transition to kernel mode is triggered by traps, exceptions, or hardware interrupts. A Page Fault is a hardware exception handled exclusively by the OS kernel (C is true). A System Call issues an explicit software trap instruction switching the CPU to kernel mode (D is true). Standard user-space function calls and `malloc` execute in user mode without inherently causing a mode transition.}

\end{tcolorbox}
\vspace{6pt}\hrule height 0.35pt\vspace{6pt}
\noindent
\begin{tcolorbox}[
    blanker,
    sidebyside,
    sidebyside align=top,
    lefthand width=0.55\textwidth,
    top=4pt, bottom=4pt, left=0pt, right=0pt,
    middle=16pt,
    underlay={
        \draw[dashed, black!80, line width=0.8pt, dash pattern=on 4pt off 3pt]
            ([xshift=0.55\textwidth+8pt]frame.north west) --
            ([xshift=0.55\textwidth+8pt]frame.south west);
    }
]
\textbf{Q.69} \quad {\footnotesize\textsc{[MSQ $\cdot$ 1 Mark(s) $\cdot$ No Neg]}} \hfill {\footnotesize\textbf{GATE 2023 CS} \quad \tiny\texttt{2023\_CS\_CS\_Q27}}
\vspace{2pt}\hrule height 0.3pt\vspace{4pt}

Which one or more of the following CPU scheduling algorithms can potentially cause starvation?

\par\vspace{6pt}\noindent
\begin{tabular}{@{}p{0.48\linewidth}@{\hspace{0.04\linewidth}}p{0.48\linewidth}@{}}
$\square$\ \textbf{(A)} First-in First-Out & $\square$\ \textbf{(B)} Round Robin \\
\addlinespace[3pt]
$\square$\ \textbf{(C)} Priority Scheduling & $\square$\ \textbf{(D)} Shortest Job First
\end{tabular}

\tcblower

\textbf{Official Key:} \textbf{\large A, C, D or C, D} \hfill \textit{\small [Easy]}
\vspace{2pt}\hrule height 0.3pt\vspace{3pt}
\textbf{Topic:} Operating System \\
\textbf{Subtopic:} \textit{CPU and I/O scheduling}
\vspace{2pt}\hrule height 0.3pt\vspace{3pt}
\textbf{Guide / Hint:} \\
{\small Priority Scheduling can cause starvation for low-priority processes if higher-priority processes arrive continuously (C). Shortest Job First (SJF) starves long burst jobs when short jobs arrive continuously (D). Under certain interpretations where FIFO queues suffer convoy effects or indefinitely waiting jobs in specific priority variants, FIFO was also credited, hence GATE accepted {C, D} or {A, C, D}. Round Robin provides bounded waiting and prevents starvation.}

\end{tcolorbox}
\vspace{6pt}\hrule height 0.35pt\vspace{6pt}
\noindent
\begin{tcolorbox}[
    blanker,
    sidebyside,
    sidebyside align=top,
    lefthand width=0.55\textwidth,
    top=4pt, bottom=4pt, left=0pt, right=0pt,
    middle=16pt,
    underlay={
        \draw[dashed, black!80, line width=0.8pt, dash pattern=on 4pt off 3pt]
            ([xshift=0.55\textwidth+8pt]frame.north west) --
            ([xshift=0.55\textwidth+8pt]frame.south west);
    }
]
\textbf{Q.70} \quad {\footnotesize\textsc{[MSQ $\cdot$ 1 Mark(s) $\cdot$ No Neg]}} \hfill {\footnotesize\textbf{GATE 2024 CS1} \quad \tiny\texttt{2024\_CS1\_CS\_Q24}}
\vspace{2pt}\hrule height 0.3pt\vspace{4pt}

Which of the following statements about threads is/are TRUE?

\par\vspace{6pt}\noindent
\begin{tabular}{@{}p{0.98\linewidth}@{}}
$\square$\ \textbf{(A)} Threads can only be implemented in kernel space \\
\addlinespace[3pt]
$\square$\ \textbf{(B)} Each thread has its own file descriptor table for open files \\
\addlinespace[3pt]
$\square$\ \textbf{(C)} All the threads belonging to a process share a common stack \\
\addlinespace[3pt]
$\square$\ \textbf{(D)} Threads belonging to a process are by default not protected from each other
\end{tabular}

\tcblower

\textbf{Official Key:} \textbf{\large D} \hfill \textit{\small [Easy]}
\vspace{2pt}\hrule height 0.3pt\vspace{3pt}
\textbf{Topic:} Operating System \\
\textbf{Subtopic:} \textit{System calls, processes, threads, inter-process communication}
\vspace{2pt}\hrule height 0.3pt\vspace{3pt}
\textbf{Guide / Hint:} \\
{\small All threads of a single process share the same address space (code, global data, heap) and open files. Because they run within the same virtual address space without memory protection barriers between them, one thread can corrupt another thread's memory; hence they are by default not protected from each other (D is TRUE). Each thread maintains its own private stack and registers.}

\end{tcolorbox}
\vspace{6pt}\hrule height 0.35pt\vspace{6pt}
\noindent
\begin{tcolorbox}[
    blanker,
    sidebyside,
    sidebyside align=top,
    lefthand width=0.55\textwidth,
    top=4pt, bottom=4pt, left=0pt, right=0pt,
    middle=16pt,
    underlay={
        \draw[dashed, black!80, line width=0.8pt, dash pattern=on 4pt off 3pt]
            ([xshift=0.55\textwidth+8pt]frame.north west) --
            ([xshift=0.55\textwidth+8pt]frame.south west);
    }
]
\textbf{Q.71} \quad {\footnotesize\textsc{[MSQ $\cdot$ 1 Mark(s) $\cdot$ No Neg]}} \hfill {\footnotesize\textbf{GATE 2024 CS1} \quad \tiny\texttt{2024\_CS1\_CS\_Q25}}
\vspace{2pt}\hrule height 0.3pt\vspace{4pt}

Which of the following process state transitions is/are NOT possible?

\par\vspace{6pt}\noindent
\begin{tabular}{@{}p{0.48\linewidth}@{\hspace{0.04\linewidth}}p{0.48\linewidth}@{}}
$\square$\ \textbf{(A)} Running to Ready & $\square$\ \textbf{(B)} Waiting to Running \\
\addlinespace[3pt]
$\square$\ \textbf{(C)} Ready to Waiting & $\square$\ \textbf{(D)} Running to Terminated
\end{tabular}

\tcblower

\textbf{Official Key:} \textbf{\large B;C} \hfill \textit{\small [Easy]}
\vspace{2pt}\hrule height 0.3pt\vspace{3pt}
\textbf{Topic:} Operating System \\
\textbf{Subtopic:} \textit{System calls, processes, threads, inter-process communication}
\vspace{2pt}\hrule height 0.3pt\vspace{3pt}
\textbf{Guide / Hint:} \\
{\small In the standard process state model, a process in the Waiting (Blocked) state transitions to Ready upon event/I/O completion; it cannot transition directly to Running without CPU dispatch (B is impossible). Similarly, only a Running process can initiate an I/O request and transition to Waiting; a Ready process cannot transition directly to Waiting (C is impossible).}

\end{tcolorbox}
\vspace{6pt}\hrule height 0.35pt\vspace{6pt}
\noindent
\begin{tcolorbox}[
    blanker,
    sidebyside,
    sidebyside align=top,
    lefthand width=0.55\textwidth,
    top=4pt, bottom=4pt, left=0pt, right=0pt,
    middle=16pt,
    underlay={
        \draw[dashed, black!80, line width=0.8pt, dash pattern=on 4pt off 3pt]
            ([xshift=0.55\textwidth+8pt]frame.north west) --
            ([xshift=0.55\textwidth+8pt]frame.south west);
    }
]
\textbf{Q.72} \quad {\footnotesize\textsc{[MSQ $\cdot$ 1 Mark(s) $\cdot$ No Neg]}} \hfill {\footnotesize\textbf{GATE 2024 CS2} \quad \tiny\texttt{2024\_CS2\_CS\_Q24}}
\vspace{2pt}\hrule height 0.3pt\vspace{4pt}

Which of the following tasks is/are the responsibility/responsibilities of the memory management unit (MMU) in a system with paging-based memory management?

\par\vspace{6pt}\noindent
\begin{tabular}{@{}p{0.98\linewidth}@{}}
$\square$\ \textbf{(A)} Allocate a new page table for a newly created process \\
\addlinespace[3pt]
$\square$\ \textbf{(B)} Translate a virtual address to a physical address using the page table \\
\addlinespace[3pt]
$\square$\ \textbf{(C)} Raise a trap when a virtual address is not found in the page table \\
\addlinespace[3pt]
$\square$\ \textbf{(D)} Raise a trap when a process tries to write to a page marked with read-only permission in the page table
\end{tabular}

\tcblower

\textbf{Official Key:} \textbf{\large B;C;D} \hfill \textit{\small [Medium]}
\vspace{2pt}\hrule height 0.3pt\vspace{3pt}
\textbf{Topic:} Operating System \\
\textbf{Subtopic:} \textit{Memory management and virtual memory}
\vspace{2pt}\hrule height 0.3pt\vspace{3pt}
\textbf{Guide / Hint:} \\
{\small The MMU hardware translates virtual addresses to physical addresses using page tables/TLB (B), generates a page-fault trap when a page table entry is invalid (C), and raises a protection trap on access permission violations such as writing to read-only pages (D). Page table allocation (A) is an OS kernel duty.}

\end{tcolorbox}
\vspace{6pt}\hrule height 0.35pt\vspace{6pt}
\noindent
\begin{tcolorbox}[
    blanker,
    sidebyside,
    sidebyside align=top,
    lefthand width=0.55\textwidth,
    top=4pt, bottom=4pt, left=0pt, right=0pt,
    middle=16pt,
    underlay={
        \draw[dashed, black!80, line width=0.8pt, dash pattern=on 4pt off 3pt]
            ([xshift=0.55\textwidth+8pt]frame.north west) --
            ([xshift=0.55\textwidth+8pt]frame.south west);
    }
]
\textbf{Q.73} \quad {\footnotesize\textsc{[MSQ $\cdot$ 1 Mark(s) $\cdot$ No Neg]}} \hfill {\footnotesize\textbf{GATE 2024 CS2} \quad \tiny\texttt{2024\_CS2\_CS\_Q25}}
\vspace{2pt}\hrule height 0.3pt\vspace{4pt}

Consider a process P running on a CPU. Which one or more of the following events will always trigger a context switch by the OS that results in process P moving to a non-running state (e.g., ready, blocked)?

\par\vspace{6pt}\noindent
\begin{tabular}{@{}p{0.98\linewidth}@{}}
$\square$\ \textbf{(A)} P makes a blocking system call to read a block of data from the disk \\
\addlinespace[3pt]
$\square$\ \textbf{(B)} P tries to access a page that is in the swap space, triggering a page fault \\
\addlinespace[3pt]
$\square$\ \textbf{(C)} An interrupt is raised by the disk to deliver data requested by some other process \\
\addlinespace[3pt]
$\square$\ \textbf{(D)} A timer interrupt is raised by the hardware
\end{tabular}

\tcblower

\textbf{Official Key:} \textbf{\large A;B} \hfill \textit{\small [Medium]}
\vspace{2pt}\hrule height 0.3pt\vspace{3pt}
\textbf{Topic:} Operating System \\
\textbf{Subtopic:} \textit{System calls, processes, threads, inter-process communication}
\vspace{2pt}\hrule height 0.3pt\vspace{3pt}
\textbf{Guide / Hint:} \\
{\small When process P executes a blocking I/O system call, it must wait for disk transfer and is moved to the blocked state (A). When P experiences a page fault on a swapped page, it is put into the blocked state until disk I/O completes (B). Interrupts from other devices (C) or timer ticks (D) in preemptive systems might return directly to P after ISR handling if P has highest priority.}

\end{tcolorbox}
\vspace{6pt}\hrule height 0.35pt\vspace{6pt}
\noindent
\begin{tcolorbox}[
    blanker,
    sidebyside,
    sidebyside align=top,
    lefthand width=0.55\textwidth,
    top=4pt, bottom=4pt, left=0pt, right=0pt,
    middle=16pt,
    underlay={
        \draw[dashed, black!80, line width=0.8pt, dash pattern=on 4pt off 3pt]
            ([xshift=0.55\textwidth+8pt]frame.north west) --
            ([xshift=0.55\textwidth+8pt]frame.south west);
    }
]
\textbf{Q.74} \quad {\footnotesize\textsc{[MSQ $\cdot$ 2 Mark(s) $\cdot$ No Neg]}} \hfill {\footnotesize\textbf{GATE 2024 CS2} \quad \tiny\texttt{2024\_CS2\_CS\_Q46}}
\vspace{2pt}\hrule height 0.3pt\vspace{4pt}

Consider a multi-threaded program with two threads T1 and T2. The threads share two semaphores: s1 (initialized to 1) and s2 (initialized to 0). The threads also share a global variable x (initialized to 0). The threads execute the code shown below:
\begin{verbatim}
// code of T1        // code of T2
wait(s1);            wait(s1);
x = x + 1;           x = x + 1;
print(x);            print(x);
wait(s2);            signal(s2);
signal(s1);          signal(s1);
\end{verbatim}
Which of the following outcomes is/are possible when threads T1 and T2 execute concurrently?

\par\vspace{6pt}\noindent
\begin{tabular}{@{}p{0.98\linewidth}@{}}
$\square$\ \textbf{(A)} T1 runs first and prints 1, T2 runs next and prints 2 \\
\addlinespace[3pt]
$\square$\ \textbf{(B)} T2 runs first and prints 1, T1 runs next and prints 2 \\
\addlinespace[3pt]
$\square$\ \textbf{(C)} T1 runs first and prints 1, T2 does not print anything (deadlock) \\
\addlinespace[3pt]
$\square$\ \textbf{(D)} T2 runs first and prints 1, T1 does not print anything (deadlock)
\end{tabular}

\tcblower

\textbf{Official Key:} \textbf{\large B;C} \hfill \textit{\small [Medium]}
\vspace{2pt}\hrule height 0.3pt\vspace{3pt}
\textbf{Topic:} Operating System \\
\textbf{Subtopic:} \textit{Concurrency and synchronization}
\vspace{2pt}\hrule height 0.3pt\vspace{3pt}
\textbf{Guide / Hint:} \\
{\small If T2 runs first: T2 takes s1, increments x (x=1), prints 1, signals s2, signals s1. Then T1 executes: takes s1, increments x (x=2), prints 2, completes wait(s2) without blocking (since s2 was signaled), signals s1 (outcome B). If T1 runs first: T1 takes s1, prints 1, and blocks on wait(s2) while holding s1, causing deadlock where T2 blocks on wait(s1) (outcome C).}

\end{tcolorbox}
\vspace{6pt}\hrule height 0.35pt\vspace{6pt}
\noindent
\begin{tcolorbox}[
    blanker,
    sidebyside,
    sidebyside align=top,
    lefthand width=0.55\textwidth,
    top=4pt, bottom=4pt, left=0pt, right=0pt,
    middle=16pt,
    underlay={
        \draw[dashed, black!80, line width=0.8pt, dash pattern=on 4pt off 3pt]
            ([xshift=0.55\textwidth+8pt]frame.north west) --
            ([xshift=0.55\textwidth+8pt]frame.south west);
    }
]
\textbf{Q.75} \quad {\footnotesize\textsc{[MSQ $\cdot$ 2 Mark(s) $\cdot$ No Neg]}} \hfill {\footnotesize\textbf{GATE 2025 CS2} \quad \tiny\texttt{2025\_CS2\_CS\_Q47}}
\vspace{2pt}\hrule height 0.3pt\vspace{4pt}

Consider a demand paging system with three frames, and the following page reference string: 1 2 3 4 5 4 1 6 4 2 1. Which page replacement policy produces the minimum number of page faults for this reference string?

\par\vspace{6pt}\noindent
\begin{tabular}{@{}p{0.98\linewidth}@{}}
$\square$\ \textbf{(A)} Least Recently Used page replacement policy \\
\addlinespace[3pt]
$\square$\ \textbf{(B)} Least Frequently Used page replacement policy \\
\addlinespace[3pt]
$\square$\ \textbf{(C)} Most Frequently Used page replacement policy \\
\addlinespace[3pt]
$\square$\ \textbf{(D)} Optimal page replacement policy
\end{tabular}

\tcblower

\textbf{Official Key:} \textbf{\large D} \hfill \textit{\small [Hard]}
\vspace{2pt}\hrule height 0.3pt\vspace{3pt}
\textbf{Topic:} Operating System \\
\textbf{Subtopic:} \textit{Memory management and virtual memory: page replacement algorithms}
\vspace{2pt}\hrule height 0.3pt\vspace{3pt}
\textbf{Guide / Hint:} \\
{\small Trace the page reference string through 3 frames under each algorithm. Optimal page replacement replaces the page that will not be used for the longest period of time in the future, matching the minimum hit/fault sequence.}

\end{tcolorbox}
\vspace{6pt}\hrule height 0.35pt\vspace{6pt}
\noindent
\begin{tcolorbox}[
    blanker,
    sidebyside,
    sidebyside align=top,
    lefthand width=0.55\textwidth,
    top=4pt, bottom=4pt, left=0pt, right=0pt,
    middle=16pt,
    underlay={
        \draw[dashed, black!80, line width=0.8pt, dash pattern=on 4pt off 3pt]
            ([xshift=0.55\textwidth+8pt]frame.north west) --
            ([xshift=0.55\textwidth+8pt]frame.south west);
    }
]
\textbf{Q.76} \quad {\footnotesize\textsc{[MSQ $\cdot$ 2 Mark(s) $\cdot$ No Neg]}} \hfill {\footnotesize\textbf{GATE 2025 CS2} \quad \tiny\texttt{2025\_CS2\_CS\_Q48}}
\vspace{2pt}\hrule height 0.3pt\vspace{4pt}

Let $P = \{P_1, P_2, P_3, P_4\}$ be active processes and $R = \{R_1, R_2, R_3, R_4\}$ be resources in an operating system. Given the allocation and request graph, which of the following actions makes the system deadlock free?

\par\vspace{6pt}\noindent
\begin{tabular}{@{}p{0.98\linewidth}@{}}
$\square$\ \textbf{(A)} Aborting $P_1$ makes the system deadlock free. \\
\addlinespace[3pt]
$\square$\ \textbf{(B)} Aborting $P_3$ makes the system deadlock free. \\
\addlinespace[3pt]
$\square$\ \textbf{(C)} Aborting $P_2$ makes the system deadlock free. \\
\addlinespace[3pt]
$\square$\ \textbf{(D)} Aborting $P_1$ and $P_4$ makes the system deadlock free.
\end{tabular}

\tcblower

\textbf{Official Key:} \textbf{\large C;D} \hfill \textit{\small [Hard]}
\vspace{2pt}\hrule height 0.3pt\vspace{3pt}
\textbf{Topic:} Operating System \\
\textbf{Subtopic:} \textit{Deadlock: Resource Allocation Graph and deadlock recovery}
\vspace{2pt}\hrule height 0.3pt\vspace{3pt}
\textbf{Guide / Hint:} \\
{\small In the Resource Allocation Graph, identify the directed cycles creating deadlock. Aborting process $P_2$ eliminates the dependency cycles directly, as does aborting both $P_1$ and $P_4$.}

\end{tcolorbox}
\vspace{6pt}\hrule height 0.35pt\vspace{6pt}
\vspace{8pt}
\subsection*{Databases}
\vspace{-4pt}
\noindent
\begin{tcolorbox}[
    blanker,
    sidebyside,
    sidebyside align=top,
    lefthand width=0.55\textwidth,
    top=4pt, bottom=4pt, left=0pt, right=0pt,
    middle=16pt,
    underlay={
        \draw[dashed, black!80, line width=0.8pt, dash pattern=on 4pt off 3pt]
            ([xshift=0.55\textwidth+8pt]frame.north west) --
            ([xshift=0.55\textwidth+8pt]frame.south west);
    }
]
\textbf{Q.77} \quad {\footnotesize\textsc{[MSQ $\cdot$ 1 Mark(s) $\cdot$ No Neg]}} \hfill {\footnotesize\textbf{GATE 2022 CS} \quad \tiny\texttt{2022\_CS\_CS\_Q25}}
\vspace{2pt}\hrule height 0.3pt\vspace{4pt}

Consider the following three relations in a relational database:
\[ \text{Employee}(eId, Name), \quad \text{Brand}(bId, bName), \quad \text{Own}(eId, bId) \]
Which of the following relational algebra expressions return the set of $eIds$ who own all the brands?

\par\vspace{6pt}\noindent
\begin{tabular}{@{}p{0.98\linewidth}@{}}
$\square$\ \textbf{(A)} $\Pi_{eId}(\Pi_{eId, bId}(Own) / \Pi_{bId}(Brand))$ \\
\addlinespace[3pt]
$\square$\ \textbf{(B)} $\Pi_{eId}(Own) - \Pi_{eId}((\Pi_{eId}(Own) \times \Pi_{bId}(Brand)) - Own)$ \\
\addlinespace[3pt]
$\square$\ \textbf{(C)} $\Pi_{eId}(\Pi_{eId, bId}(Own) / \Pi_{bId}(Own))$ \\
\addlinespace[3pt]
$\square$\ \textbf{(D)} $\Pi_{eId}((\Pi_{eId}(Own) \times \Pi_{bId}(Brand)) / \Pi_{bId}(Own))$
\end{tabular}

\tcblower

\textbf{Official Key:} \textbf{\large A,B} \hfill \textit{\small [Medium]}
\vspace{2pt}\hrule height 0.3pt\vspace{3pt}
\textbf{Topic:} Databases \\
\textbf{Subtopic:} \textit{Relational model: relational algebra, tuple calculus, SQL}
\vspace{2pt}\hrule height 0.3pt\vspace{3pt}
\textbf{Guide / Hint:} \\
{\small Relational division $\Pi_{eId, bId}(Own) / \Pi_{bId}(Brand)$ directly computes employees who own all brands (A is TRUE). Equivalently, taking all employees, finding those missing at least one brand via cartesian product with Brand followed by set difference with Own, and subtracting them from all employees yields the same set (B is TRUE).}

\end{tcolorbox}
\vspace{6pt}\hrule height 0.35pt\vspace{6pt}
\noindent
\begin{tcolorbox}[
    blanker,
    sidebyside,
    sidebyside align=top,
    lefthand width=0.55\textwidth,
    top=4pt, bottom=4pt, left=0pt, right=0pt,
    middle=16pt,
    underlay={
        \draw[dashed, black!80, line width=0.8pt, dash pattern=on 4pt off 3pt]
            ([xshift=0.55\textwidth+8pt]frame.north west) --
            ([xshift=0.55\textwidth+8pt]frame.south west);
    }
]
\textbf{Q.78} \quad {\footnotesize\textsc{[MSQ $\cdot$ 1 Mark(s) $\cdot$ No Neg]}} \hfill {\footnotesize\textbf{GATE 2024 CS1} \quad \tiny\texttt{2024\_CS1\_CS\_Q22}}
\vspace{2pt}\hrule height 0.3pt\vspace{4pt}

Which of the following statements about a relation $R$ in first normal form (1NF) is/are TRUE?

\par\vspace{6pt}\noindent
\begin{tabular}{@{}p{0.98\linewidth}@{}}
$\square$\ \textbf{(A)} $R$ can have a multi-attribute key \\
\addlinespace[3pt]
$\square$\ \textbf{(B)} $R$ cannot have a foreign key \\
\addlinespace[3pt]
$\square$\ \textbf{(C)} $R$ cannot have a composite attribute \\
\addlinespace[3pt]
$\square$\ \textbf{(D)} $R$ cannot have more than one candidate key
\end{tabular}

\tcblower

\textbf{Official Key:} \textbf{\large A;C} \hfill \textit{\small [Easy]}
\vspace{2pt}\hrule height 0.3pt\vspace{3pt}
\textbf{Topic:} Databases \\
\textbf{Subtopic:} \textit{Integrity constraints, normal forms}
\vspace{2pt}\hrule height 0.3pt\vspace{3pt}
\textbf{Guide / Hint:} \\
{\small A relation is in 1NF if and only if all attribute values are atomic, which explicitly forbids composite or multi-valued attributes (making (C) TRUE). 1NF places no limitation on candidate keys having multiple attributes (multi-attribute keys are fully permitted, making (A) TRUE) or the presence of foreign keys.}

\end{tcolorbox}
\vspace{6pt}\hrule height 0.35pt\vspace{6pt}
\noindent
\begin{tcolorbox}[
    blanker,
    sidebyside,
    sidebyside align=top,
    lefthand width=0.55\textwidth,
    top=4pt, bottom=4pt, left=0pt, right=0pt,
    middle=16pt,
    underlay={
        \draw[dashed, black!80, line width=0.8pt, dash pattern=on 4pt off 3pt]
            ([xshift=0.55\textwidth+8pt]frame.north west) --
            ([xshift=0.55\textwidth+8pt]frame.south west);
    }
]
\textbf{Q.79} \quad {\footnotesize\textsc{[MSQ $\cdot$ 2 Mark(s) $\cdot$ No Neg]}} \hfill {\footnotesize\textbf{GATE 2024 CS1} \quad \tiny\texttt{2024\_CS1\_CS\_Q44}}
\vspace{2pt}\hrule height 0.3pt\vspace{4pt}

The symbol $\to$ indicates functional dependency in the context of a relational database. Which of the following options is/are TRUE?

\par\vspace{6pt}\noindent
\begin{tabular}{@{}p{0.98\linewidth}@{}}
$\square$\ \textbf{(A)} $(X, Y) \to (Z, W)$ implies $X \to (Z, W)$ \\
\addlinespace[3pt]
$\square$\ \textbf{(B)} $(X, Y) \to (Z, W)$ implies $(X, Y) \to Z$ \\
\addlinespace[3pt]
$\square$\ \textbf{(C)} $((X, Y) \to Z \text{ and } W \to Y)$ implies $(X, W) \to Z$ \\
\addlinespace[3pt]
$\square$\ \textbf{(D)} $(X \to Y \text{ and } Y \to Z)$ implies $X \to Z$
\end{tabular}

\tcblower

\textbf{Official Key:} \textbf{\large B;C;D} \hfill \textit{\small [Medium]}
\vspace{2pt}\hrule height 0.3pt\vspace{3pt}
\textbf{Topic:} Databases \\
\textbf{Subtopic:} \textit{Integrity constraints, normal forms}
\vspace{2pt}\hrule height 0.3pt\vspace{3pt}
\textbf{Guide / Hint:} \\
{\small By Armstrong's decomposition axiom, $(X, Y) \to (Z, W)$ implies $(X, Y) \to Z$ (B is TRUE). By Armstrong's transitivity axiom, $X \to Y$ and $Y \to Z$ implies $X \to Z$ (D is TRUE). By pseudotransitivity, $(X, Y) \to Z$ and $W \to Y$ implies $(X, W) \to Z$ (C is TRUE). Decomposing the left-hand side determinant is invalid in general (A is FALSE).}

\end{tcolorbox}
\vspace{6pt}\hrule height 0.35pt\vspace{6pt}
\noindent
\begin{tcolorbox}[
    blanker,
    sidebyside,
    sidebyside align=top,
    lefthand width=0.55\textwidth,
    top=4pt, bottom=4pt, left=0pt, right=0pt,
    middle=16pt,
    underlay={
        \draw[dashed, black!80, line width=0.8pt, dash pattern=on 4pt off 3pt]
            ([xshift=0.55\textwidth+8pt]frame.north west) --
            ([xshift=0.55\textwidth+8pt]frame.south west);
    }
]
\textbf{Q.80} \quad {\footnotesize\textsc{[MSQ $\cdot$ 2 Mark(s) $\cdot$ No Neg]}} \hfill {\footnotesize\textbf{GATE 2024 CS1} \quad \tiny\texttt{2024\_CS1\_CS\_Q46}}
\vspace{2pt}\hrule height 0.3pt\vspace{4pt}

Consider the following read-write schedule $S$ over three transactions $T_1, T_2$, and $T_3$, where the subscripts in the schedule indicate transaction IDs:
\[ S: r\_1(z); w\_1(z); r\_2(x); r\_3(y); w\_3(y); r\_2(y); w\_2(x); w\_2(y); \]
Which of the following transaction schedules is/are conflict equivalent to $S$?

\par\vspace{6pt}\noindent
\begin{tabular}{@{}p{0.48\linewidth}@{\hspace{0.04\linewidth}}p{0.48\linewidth}@{}}
$\square$\ \textbf{(A)} $T_1 T_2 T_3$ & $\square$\ \textbf{(B)} $T_1 T_3 T_2$ \\
\addlinespace[3pt]
$\square$\ \textbf{(C)} $T_3 T_2 T_1$ & $\square$\ \textbf{(D)} $T_3 T_1 T_2$
\end{tabular}

\tcblower

\textbf{Official Key:} \textbf{\large B;C;D} \hfill \textit{\small [Medium]}
\vspace{2pt}\hrule height 0.3pt\vspace{3pt}
\textbf{Topic:} Databases \\
\textbf{Subtopic:} \textit{Transactions and concurrency control}
\vspace{2pt}\hrule height 0.3pt\vspace{3pt}
\textbf{Guide / Hint:} \\
{\small Examining data item conflicts: on item $y$, $w_3(y)$ in $T_3$ precedes both $r_2(y)$ and $w_2(y)$ in $T_2$, creating a precedence conflict edge $T_3 \to T_2$. Transaction $T_1$ only accesses $z$ and has no conflicting operations with $T_2$ or $T_3$. Any serial schedule preserving $T_3$ before $T_2$ is conflict equivalent to $S$, namely $T_1 T_3 T_2$, $T_3 T_1 T_2$, and $T_3 T_2 T_1$.}

\end{tcolorbox}
\vspace{6pt}\hrule height 0.35pt\vspace{6pt}
\noindent
\begin{tcolorbox}[
    blanker,
    sidebyside,
    sidebyside align=top,
    lefthand width=0.55\textwidth,
    top=4pt, bottom=4pt, left=0pt, right=0pt,
    middle=16pt,
    underlay={
        \draw[dashed, black!80, line width=0.8pt, dash pattern=on 4pt off 3pt]
            ([xshift=0.55\textwidth+8pt]frame.north west) --
            ([xshift=0.55\textwidth+8pt]frame.south west);
    }
]
\textbf{Q.81} \quad {\footnotesize\textsc{[MSQ $\cdot$ 1 Mark(s) $\cdot$ No Neg]}} \hfill {\footnotesize\textbf{GATE 2024 CS2} \quad \tiny\texttt{2024\_CS2\_CS\_Q26}}
\vspace{2pt}\hrule height 0.3pt\vspace{4pt}

Which of the following file organizations is/are I/O efficient for the scan operation in DBMS?

\par\vspace{6pt}\noindent
\begin{tabular}{@{}p{0.48\linewidth}@{\hspace{0.04\linewidth}}p{0.48\linewidth}@{}}
$\square$\ \textbf{(A)} Sorted & $\square$\ \textbf{(B)} Heap \\
\addlinespace[3pt]
$\square$\ \textbf{(C)} Unclustered tree index & $\square$\ \textbf{(D)} Unclustered hash index
\end{tabular}

\tcblower

\textbf{Official Key:} \textbf{\large A;B} \hfill \textit{\small [Easy]}
\vspace{2pt}\hrule height 0.3pt\vspace{3pt}
\textbf{Topic:} Databases \\
\textbf{Subtopic:} \textit{File organization, indexing (e.g., B and B+ trees)}
\vspace{2pt}\hrule height 0.3pt\vspace{3pt}
\textbf{Guide / Hint:} \\
{\small A full table scan retrieves all pages sequentially. Sorted and Heap files store data records densely packed across consecutive data pages, requiring reading each page exactly once. Unclustered indexes store data out of order, incurring multiple random I/Os per record.}

\end{tcolorbox}
\vspace{6pt}\hrule height 0.35pt\vspace{6pt}
\noindent
\begin{tcolorbox}[
    blanker,
    sidebyside,
    sidebyside align=top,
    lefthand width=0.55\textwidth,
    top=4pt, bottom=4pt, left=0pt, right=0pt,
    middle=16pt,
    underlay={
        \draw[dashed, black!80, line width=0.8pt, dash pattern=on 4pt off 3pt]
            ([xshift=0.55\textwidth+8pt]frame.north west) --
            ([xshift=0.55\textwidth+8pt]frame.south west);
    }
]
\textbf{Q.82} \quad {\footnotesize\textsc{[MSQ $\cdot$ 1 Mark(s) $\cdot$ No Neg]}} \hfill {\footnotesize\textbf{GATE 2024 CS2} \quad \tiny\texttt{2024\_CS2\_CS\_Q27}}
\vspace{2pt}\hrule height 0.3pt\vspace{4pt}

Which of the following statements about the Two Phase Locking (2PL) protocol is/are TRUE?

\par\vspace{6pt}\noindent
\begin{tabular}{@{}p{0.98\linewidth}@{}}
$\square$\ \textbf{(A)} 2PL permits only serializable schedules \\
\addlinespace[3pt]
$\square$\ \textbf{(B)} With 2PL, a transaction always locks the data item being read or written just before every operation \\
\addlinespace[3pt]
$\square$\ \textbf{(C)} With 2PL, once a lock is released on any data item inside a transaction, no more locks on any data item can be acquired \\
\addlinespace[3pt]
$\square$\ \textbf{(D)} A deadlock is possible with 2PL
\end{tabular}

\tcblower

\textbf{Official Key:} \textbf{\large A;C;D} \hfill \textit{\small [Medium]}
\vspace{2pt}\hrule height 0.3pt\vspace{3pt}
\textbf{Topic:} Databases \\
\textbf{Subtopic:} \textit{Transactions and concurrency control}
\vspace{2pt}\hrule height 0.3pt\vspace{3pt}
\textbf{Guide / Hint:} \\
{\small Two-Phase Locking (2PL) guarantees conflict serializability (A). In 2PL, once a lock is released (shrinking phase), no new locks can be acquired (C). Basic 2PL is susceptible to deadlocks due to circular wait conditions (D). Statement B describes conservative/static locking, not general 2PL.}

\end{tcolorbox}
\vspace{6pt}\hrule height 0.35pt\vspace{6pt}
\noindent
\begin{tcolorbox}[
    blanker,
    sidebyside,
    sidebyside align=top,
    lefthand width=0.55\textwidth,
    top=4pt, bottom=4pt, left=0pt, right=0pt,
    middle=16pt,
    underlay={
        \draw[dashed, black!80, line width=0.8pt, dash pattern=on 4pt off 3pt]
            ([xshift=0.55\textwidth+8pt]frame.north west) --
            ([xshift=0.55\textwidth+8pt]frame.south west);
    }
]
\textbf{Q.83} \quad {\footnotesize\textsc{[MSQ $\cdot$ 1 Mark(s) $\cdot$ No Neg]}} \hfill {\footnotesize\textbf{GATE 2025 CS1} \quad \tiny\texttt{2025\_CS1\_CS\_Q21}}
\vspace{2pt}\hrule height 0.3pt\vspace{4pt}

Consider a $B^+$ tree with 5 nodes, in which a node can store at most 3 key values. The value 23 is now inserted in the $B^+$ tree. Which of the following option(s) is/are CORRECT?

\par\vspace{6pt}\noindent
\begin{tabular}{@{}p{0.98\linewidth}@{}}
$\square$\ \textbf{(A)} None of the nodes will split. \\
\addlinespace[3pt]
$\square$\ \textbf{(B)} At least one node will split and redistribute. \\
\addlinespace[3pt]
$\square$\ \textbf{(C)} The total number of nodes will remain same. \\
\addlinespace[3pt]
$\square$\ \textbf{(D)} The height of the tree will increase.
\end{tabular}

\tcblower

\textbf{Official Key:} \textbf{\large B;D} \hfill \textit{\small [Medium]}
\vspace{2pt}\hrule height 0.3pt\vspace{3pt}
\textbf{Topic:} Databases \\
\textbf{Subtopic:} \textit{File organization, indexing (e.g., B and B+ trees)}
\vspace{2pt}\hrule height 0.3pt\vspace{3pt}
\textbf{Guide / Hint:} \\
{\small With maximum key capacity of 3 keys per node, the targeted leaf node containing keys $[20, 21, 22]$ is already full. Inserting 23 causes this leaf node to overflow and split, propagating a split upward to the already-saturated root, which creates a new root node and increases tree height.}

\end{tcolorbox}
\vspace{6pt}\hrule height 0.35pt\vspace{6pt}
\noindent
\begin{tcolorbox}[
    blanker,
    sidebyside,
    sidebyside align=top,
    lefthand width=0.55\textwidth,
    top=4pt, bottom=4pt, left=0pt, right=0pt,
    middle=16pt,
    underlay={
        \draw[dashed, black!80, line width=0.8pt, dash pattern=on 4pt off 3pt]
            ([xshift=0.55\textwidth+8pt]frame.north west) --
            ([xshift=0.55\textwidth+8pt]frame.south west);
    }
]
\textbf{Q.84} \quad {\footnotesize\textsc{[MSQ $\cdot$ 2 Mark(s) $\cdot$ No Neg]}} \hfill {\footnotesize\textbf{GATE 2025 CS1} \quad \tiny\texttt{2025\_CS1\_CS\_Q47}}
\vspace{2pt}\hrule height 0.3pt\vspace{4pt}

Consider a relational schema $team(name, city, owner)$, with functional dependencies $\{name \to city, name \to owner\}$. The relation $team$ is decomposed into two relations, $t_1(name, city)$ and $t_2(name, owner)$. Which of the following statement(s) is/are TRUE?

\par\vspace{6pt}\noindent
\begin{tabular}{@{}p{0.98\linewidth}@{}}
$\square$\ \textbf{(A)} The relation $team$ is NOT in BCNF. \\
\addlinespace[3pt]
$\square$\ \textbf{(B)} The relations $t_1$ and $t_2$ are in BCNF. \\
\addlinespace[3pt]
$\square$\ \textbf{(C)} The decomposition constitutes a lossless join. \\
\addlinespace[3pt]
$\square$\ \textbf{(D)} The relation $team$ is NOT in 3NF.
\end{tabular}

\tcblower

\textbf{Official Key:} \textbf{\large B;C} \hfill \textit{\small [Medium]}
\vspace{2pt}\hrule height 0.3pt\vspace{3pt}
\textbf{Topic:} Databases \\
\textbf{Subtopic:} \textit{Integrity constraints, normal forms}
\vspace{2pt}\hrule height 0.3pt\vspace{3pt}
\textbf{Guide / Hint:} \\
{\small In $team$, $name$ is a candidate key since $name^+ = \{name, city, owner\}$. Since every FD has a superkey on the LHS, $team$ was already in BCNF and 3NF (A and D are false). In $t_1$ and $t_2$, $name$ is the primary key, so both are in BCNF (B is true). The common attribute $R_1 \cap R_2 = \{name\}$ is a key of both $t_1$ and $t_2$, ensuring a lossless join (C is true).}

\end{tcolorbox}
\vspace{6pt}\hrule height 0.35pt\vspace{6pt}
\noindent
\begin{tcolorbox}[
    blanker,
    sidebyside,
    sidebyside align=top,
    lefthand width=0.55\textwidth,
    top=4pt, bottom=4pt, left=0pt, right=0pt,
    middle=16pt,
    underlay={
        \draw[dashed, black!80, line width=0.8pt, dash pattern=on 4pt off 3pt]
            ([xshift=0.55\textwidth+8pt]frame.north west) --
            ([xshift=0.55\textwidth+8pt]frame.south west);
    }
]
\textbf{Q.85} \quad {\footnotesize\textsc{[MSQ $\cdot$ 1 Mark(s) $\cdot$ No Neg]}} \hfill {\footnotesize\textbf{GATE 2025 CS2} \quad \tiny\texttt{2025\_CS2\_CS\_Q27}}
\vspace{2pt}\hrule height 0.3pt\vspace{4pt}

An audit of a banking transactions system has found that on an earlier occasion, two joint holders of account $A$ attempted simultaneous transfers of Rs. 10000 each from account $A$ to account $B$. Both transactions read the same value, Rs. 11000, as the initial balance in $A$ and were allowed to go through. $B$ was credited Rs. 10000 twice. $A$ was debited only once and ended up with a balance of Rs. 1000. Which of the following properties is/are certain to have been violated by the system?

\par\vspace{6pt}\noindent
\begin{tabular}{@{}p{0.48\linewidth}@{\hspace{0.04\linewidth}}p{0.48\linewidth}@{}}
$\square$\ \textbf{(A)} Atomicity & $\square$\ \textbf{(B)} Consistency \\
\addlinespace[3pt]
$\square$\ \textbf{(C)} Isolation & $\square$\ \textbf{(D)} Durability
\end{tabular}

\tcblower

\textbf{Official Key:} \textbf{\large B;C} \hfill \textit{\small [Medium]}
\vspace{2pt}\hrule height 0.3pt\vspace{3pt}
\textbf{Topic:} Databases \\
\textbf{Subtopic:} \textit{Transactions and concurrency control: ACID properties}
\vspace{2pt}\hrule height 0.3pt\vspace{3pt}
\textbf{Guide / Hint:} \\
{\small Both transactions read the same uncommitted state and overwrote each other (lost update anomaly), violating Isolation. The final balance violates the database invariant (starting 11000 debited $2 \times 10000$ cannot yield 1000), violating Consistency.}

\end{tcolorbox}
\vspace{6pt}\hrule height 0.35pt\vspace{6pt}
\noindent
\begin{tcolorbox}[
    blanker,
    sidebyside,
    sidebyside align=top,
    lefthand width=0.55\textwidth,
    top=4pt, bottom=4pt, left=0pt, right=0pt,
    middle=16pt,
    underlay={
        \draw[dashed, black!80, line width=0.8pt, dash pattern=on 4pt off 3pt]
            ([xshift=0.55\textwidth+8pt]frame.north west) --
            ([xshift=0.55\textwidth+8pt]frame.south west);
    }
]
\textbf{Q.86} \quad {\footnotesize\textsc{[MSQ $\cdot$ 2 Mark(s) $\cdot$ No Neg]}} \hfill {\footnotesize\textbf{GATE 2025 CS2} \quad \tiny\texttt{2025\_CS2\_CS\_Q46}}
\vspace{2pt}\hrule height 0.3pt\vspace{4pt}

Consider the following relational schema along with all the functional dependencies that hold on them:
$R_1(A, B, C, D)$ with FDs, and $R_2(E, F, G, H)$ with FDs.
Which of the following statements is/are CORRECT?

\par\vspace{6pt}\noindent
\begin{tabular}{@{}p{0.48\linewidth}@{\hspace{0.04\linewidth}}p{0.48\linewidth}@{}}
$\square$\ \textbf{(A)} $R_1$ is in 3NF & $\square$\ \textbf{(B)} $R_2$ is in 3NF \\
\addlinespace[3pt]
$\square$\ \textbf{(C)} $R_1$ is NOT in 3NF & $\square$\ \textbf{(D)} $R_2$ is NOT in 3NF
\end{tabular}

\tcblower

\textbf{Official Key:} \textbf{\large C;D} \hfill \textit{\small [Medium]}
\vspace{2pt}\hrule height 0.3pt\vspace{3pt}
\textbf{Topic:} Databases \\
\textbf{Subtopic:} \textit{Integrity constraints, normal forms: 3NF and BCNF}
\vspace{2pt}\hrule height 0.3pt\vspace{3pt}
\textbf{Guide / Hint:} \\
{\small A schema is in 3NF if for every non-trivial FD $X \to Y$, either $X$ is a superkey or $Y$ consists of prime attributes. Here, both schemas contain dependencies where the LHS is not a superkey and the RHS has non-prime attributes, so neither is in 3NF.}

\end{tcolorbox}
\vspace{6pt}\hrule height 0.35pt\vspace{6pt}
\noindent
\begin{tcolorbox}[
    blanker,
    sidebyside,
    sidebyside align=top,
    lefthand width=0.55\textwidth,
    top=4pt, bottom=4pt, left=0pt, right=0pt,
    middle=16pt,
    underlay={
        \draw[dashed, black!80, line width=0.8pt, dash pattern=on 4pt off 3pt]
            ([xshift=0.55\textwidth+8pt]frame.north west) --
            ([xshift=0.55\textwidth+8pt]frame.south west);
    }
]
\textbf{Q.87} \quad {\footnotesize\textsc{[MSQ $\cdot$ 2 Mark(s) $\cdot$ No Neg]}} \hfill {\footnotesize\textbf{GATE 2025 CS2} \quad \tiny\texttt{2025\_CS2\_CS\_Q54}}
\vspace{2pt}\hrule height 0.3pt\vspace{4pt}

Consider the relational schema:
- $\text{Students}(\underline{\text{rollno}}, \text{name}, \text{age}, \text{cgpa})$
- $\text{Enrolled}(\underline{\text{rollno}, \text{courseno}})$
\par\medskip
Which of the following SQL query/queries correctly find(s) the name of all students who have enrolled in course 1470?

\par\vspace{6pt}\noindent
\begin{tabular}{@{}p{0.98\linewidth}@{}}
$\square$\ \textbf{(A)} SELECT S.name \newline FROM Students S \newline WHERE EXISTS (SELECT * FROM Enrolled E WHERE E.courseno = 1470 AND E.rollno = S.rollno); \\
\addlinespace[3pt]
$\square$\ \textbf{(B)} SELECT S.name \newline FROM Students S \newline WHERE SIZEOF (SELECT * FROM Enrolled E WHERE E.courseno = 1470 AND E.rollno = S.rollno) > 0; \\
\addlinespace[3pt]
$\square$\ \textbf{(C)} SELECT S.name \newline FROM Students S \newline WHERE 0 < (SELECT COUNT(*) FROM Enrolled E WHERE E.courseno = 1470 AND E.rollno = S.rollno); \\
\addlinespace[3pt]
$\square$\ \textbf{(D)} SELECT S.name \newline FROM Students S NATURAL JOIN Enrolled E \newline WHERE E.courseno = 1470;
\end{tabular}

\tcblower

\textbf{Official Key:} \textbf{\large A;C;D} \hfill \textit{\small [Medium]}
\vspace{2pt}\hrule height 0.3pt\vspace{3pt}
\textbf{Topic:} Databases \\
\textbf{Subtopic:} \textit{Relational model: SQL queries and joins}
\vspace{2pt}\hrule height 0.3pt\vspace{3pt}
\textbf{Guide / Hint:} \\
{\small Queries A, C, and D all correctly return the names of students enrolled in course 1470 using EXISTS, COUNT, and natural join respectively. Query B is invalid because SQL does not have a `SIZEOF` operator.}

\end{tcolorbox}
\vspace{6pt}\hrule height 0.35pt\vspace{6pt}
\vspace{8pt}
\subsection*{Computer Networks}
\vspace{-4pt}
\noindent
\begin{tcolorbox}[
    blanker,
    sidebyside,
    sidebyside align=top,
    lefthand width=0.55\textwidth,
    top=4pt, bottom=4pt, left=0pt, right=0pt,
    middle=16pt,
    underlay={
        \draw[dashed, black!80, line width=0.8pt, dash pattern=on 4pt off 3pt]
            ([xshift=0.55\textwidth+8pt]frame.north west) --
            ([xshift=0.55\textwidth+8pt]frame.south west);
    }
]
\textbf{Q.88} \quad {\footnotesize\textsc{[MSQ $\cdot$ 1 Mark(s) $\cdot$ No Neg]}} \hfill {\footnotesize\textbf{GATE 2023 CS} \quad \tiny\texttt{2023\_CS\_CS\_Q25}}
\vspace{2pt}\hrule height 0.3pt\vspace{4pt}

Which of the following statements is/are INCORRECT about the OSPF (Open Shortest Path First) routing protocol used in the Internet?

\par\vspace{6pt}\noindent
\begin{tabular}{@{}p{0.98\linewidth}@{}}
$\square$\ \textbf{(A)} OSPF implements Bellman-Ford algorithm to find shortest paths. \\
\addlinespace[3pt]
$\square$\ \textbf{(B)} OSPF uses Dijkstra's shortest path algorithm to implement least-cost path routing. \\
\addlinespace[3pt]
$\square$\ \textbf{(C)} OSPF is used as an inter-domain routing protocol. \\
\addlinespace[3pt]
$\square$\ \textbf{(D)} OSPF implements hierarchical routing.
\end{tabular}

\tcblower

\textbf{Official Key:} \textbf{\large A, C} \hfill \textit{\small [Easy]}
\vspace{2pt}\hrule height 0.3pt\vspace{3pt}
\textbf{Topic:} Computer Networks \\
\textbf{Subtopic:} \textit{Routing protocols: shortest path, flooding, distance vector and link state routing}
\vspace{2pt}\hrule height 0.3pt\vspace{3pt}
\textbf{Guide / Hint:} \\
{\small OSPF is an intra-domain (Interior Gateway) routing protocol that uses Dijkstra's link-state algorithm (not Bellman-Ford, which is used by distance-vector protocols like RIP). Therefore, statement A (implements Bellman-Ford) and statement C (used as inter-domain) are INCORRECT, which are the required answers.}

\end{tcolorbox}
\vspace{6pt}\hrule height 0.35pt\vspace{6pt}
\noindent
\begin{tcolorbox}[
    blanker,
    sidebyside,
    sidebyside align=top,
    lefthand width=0.55\textwidth,
    top=4pt, bottom=4pt, left=0pt, right=0pt,
    middle=16pt,
    underlay={
        \draw[dashed, black!80, line width=0.8pt, dash pattern=on 4pt off 3pt]
            ([xshift=0.55\textwidth+8pt]frame.north west) --
            ([xshift=0.55\textwidth+8pt]frame.south west);
    }
]
\textbf{Q.89} \quad {\footnotesize\textsc{[MSQ $\cdot$ 2 Mark(s) $\cdot$ No Neg]}} \hfill {\footnotesize\textbf{GATE 2023 CS} \quad \tiny\texttt{2023\_CS\_CS\_Q50}}
\vspace{2pt}\hrule height 0.3pt\vspace{4pt}

Suppose you are asked to design a new reliable byte-stream transport protocol like TCP. This protocol, named myTCP, runs over a 100 Mbps network with Round Trip Time of 150 milliseconds and the maximum segment lifetime of 2 minutes.
Which of the following is/are valid lengths of the Sequence Number field in the myTCP header?

\par\vspace{6pt}\noindent
\begin{tabular}{@{}p{0.48\linewidth}@{\hspace{0.04\linewidth}}p{0.48\linewidth}@{}}
$\square$\ \textbf{(A)} 30 bits & $\square$\ \textbf{(B)} 32 bits \\
\addlinespace[3pt]
$\square$\ \textbf{(C)} 34 bits & $\square$\ \textbf{(D)} 36 bits
\end{tabular}

\tcblower

\textbf{Official Key:} \textbf{\large B, C, D} \hfill \textit{\small [Hard]}
\vspace{2pt}\hrule height 0.3pt\vspace{3pt}
\textbf{Topic:} Computer Networks \\
\textbf{Subtopic:} \textit{Transport layer: flow control and congestion control, UDP, TCP, sockets}
\vspace{2pt}\hrule height 0.3pt\vspace{3pt}
\textbf{Guide / Hint:} \\
{\small To prevent sequence number wraparound within the Maximum Segment Lifetime ($MSL = 2\text{ min} = 120\text{ s}$), the sequence number space $2^k$ must be at least the data bytes transmitted at line rate during MSL: $2^k \ge \frac{100 \times 10^6\text{ bps}}{8} \times 120\text{ s} = 1.5 \times 10^9\text{ bytes}$. Since $2^{30} \approx 1.07 \times 10^9 < 1.5 \times 10^9$ and $2^{31} \approx 2.15 \times 10^9 > 1.5 \times 10^9$, at least 31 bits are required. Thus, 32 bits, 34 bits, and 36 bits are all valid lengths.}

\end{tcolorbox}
\vspace{6pt}\hrule height 0.35pt\vspace{6pt}
\noindent
\begin{tcolorbox}[
    blanker,
    sidebyside,
    sidebyside align=top,
    lefthand width=0.55\textwidth,
    top=4pt, bottom=4pt, left=0pt, right=0pt,
    middle=16pt,
    underlay={
        \draw[dashed, black!80, line width=0.8pt, dash pattern=on 4pt off 3pt]
            ([xshift=0.55\textwidth+8pt]frame.north west) --
            ([xshift=0.55\textwidth+8pt]frame.south west);
    }
]
\textbf{Q.90} \quad {\footnotesize\textsc{[MSQ $\cdot$ 2 Mark(s) $\cdot$ No Neg]}} \hfill {\footnotesize\textbf{GATE 2023 CS} \quad \tiny\texttt{2023\_CS\_CS\_Q52}}
\vspace{2pt}\hrule height 0.3pt\vspace{4pt}

Suppose in a web browser, you click on the `www.gate-2023.in` URL. The browser cache is empty. The IP address for this URL is not cached in your local host, so a DNS lookup is triggered (by the local DNS server deployed on your local host) over the 3-tier DNS hierarchy in an iterative mode. No resource records are cached anywhere across all DNS servers.
Let RTT denote the round trip time between your local host and DNS servers in the DNS hierarchy. The round trip time between the local host and the web server hosting `www.gate-2023.in` is also equal to RTT. The HTML file associated with the URL is small enough to have negligible transmission time and negligible rendering time by your web browser, which references 10 equally small objects on the same web server.
Which of the following statements is/are CORRECT about the minimum elapsed time between clicking on the URL and your browser fully rendering it?

\par\vspace{6pt}\noindent
\begin{tabular}{@{}p{0.98\linewidth}@{}}
$\square$\ \textbf{(A)} 7 RTTs, in case of non-persistent HTTP with 5 parallel TCP connections. \\
\addlinespace[3pt]
$\square$\ \textbf{(B)} 5 RTTs, in case of persistent HTTP with pipelining. \\
\addlinespace[3pt]
$\square$\ \textbf{(C)} 9 RTTs, in case of non-persistent HTTP with 5 parallel TCP connections. \\
\addlinespace[3pt]
$\square$\ \textbf{(D)} 6 RTTs, in case of persistent HTTP with pipelining.
\end{tabular}

\tcblower

\textbf{Official Key:} \textbf{\large C, D} \hfill \textit{\small [Hard]}
\vspace{2pt}\hrule height 0.3pt\vspace{3pt}
\textbf{Topic:} Computer Networks \\
\textbf{Subtopic:} \textit{Application layer protocols: DNS, SMTP, HTTP, FTP, Email}
\vspace{2pt}\hrule height 0.3pt\vspace{3pt}
\textbf{Guide / Hint:} \\
{\small Resolving the URL via 3-tier iterative DNS requires $3\text{ RTT}$. Initiating TCP connection and requesting the base HTML page requires $1 + 1 = 2\text{ RTT}$, so the base HTML is received after $3 + 2 = 5\text{ RTT}$. With persistent HTTP pipelining, all 10 objects are fetched in 1 round trip over the established connection, giving $5 + 1 = 6\text{ RTT}$ (D is TRUE). With non-persistent HTTP and 5 parallel connections, fetching 10 objects takes 2 sequential batches; each batch takes $1\text{ RTT (TCP)} + 1\text{ RTT (HTTP)} = 2\text{ RTT}$, adding $2 \times 2 = 4\text{ RTT}$, giving $5 + 4 = 9\text{ RTT}$ (C is TRUE).}

\end{tcolorbox}
\vspace{6pt}\hrule height 0.35pt\vspace{6pt}
\noindent
\begin{tcolorbox}[
    blanker,
    sidebyside,
    sidebyside align=top,
    lefthand width=0.55\textwidth,
    top=4pt, bottom=4pt, left=0pt, right=0pt,
    middle=16pt,
    underlay={
        \draw[dashed, black!80, line width=0.8pt, dash pattern=on 4pt off 3pt]
            ([xshift=0.55\textwidth+8pt]frame.north west) --
            ([xshift=0.55\textwidth+8pt]frame.south west);
    }
]
\textbf{Q.91} \quad {\footnotesize\textsc{[MSQ $\cdot$ 1 Mark(s) $\cdot$ No Neg]}} \hfill {\footnotesize\textbf{GATE 2024 CS1} \quad \tiny\texttt{2024\_CS1\_CS\_Q29}}
\vspace{2pt}\hrule height 0.3pt\vspace{4pt}

TCP client P successfully establishes a connection to TCP server Q. Let $N_P$ denote the sequence number in the SYN sent from P to Q. Let $N_Q$ denote the acknowledgement number in the SYN ACK from Q to P. Which of the following statements is/are CORRECT?

\par\vspace{6pt}\noindent
\begin{tabular}{@{}p{0.98\linewidth}@{}}
$\square$\ \textbf{(A)} The sequence number $N_P$ is chosen randomly by P \\
\addlinespace[3pt]
$\square$\ \textbf{(B)} The sequence number $N_P$ is always 0 for a new connection \\
\addlinespace[3pt]
$\square$\ \textbf{(C)} The acknowledgement number $N_Q$ is equal to $N_P$ \\
\addlinespace[3pt]
$\square$\ \textbf{(D)} The acknowledgement number $N_Q$ is equal to $N_P + 1$
\end{tabular}

\tcblower

\textbf{Official Key:} \textbf{\large A;D} \hfill \textit{\small [Easy]}
\vspace{2pt}\hrule height 0.3pt\vspace{3pt}
\textbf{Topic:} Computer Networks \\
\textbf{Subtopic:} \textit{Transport layer: flow control and congestion control, UDP, TCP, sockets}
\vspace{2pt}\hrule height 0.3pt\vspace{3pt}
\textbf{Guide / Hint:} \\
{\small During TCP connection establishment, the initial sequence number $N_P$ is chosen pseudo-randomly to avoid cross-connection packet collisions and mitigate spoofing attacks (A is TRUE). Because a SYN packet consumes exactly 1 sequence number in the sequence space, server $Q$ acknowledges it with $N_Q = N_P + 1$ (D is TRUE).}

\end{tcolorbox}
\vspace{6pt}\hrule height 0.35pt\vspace{6pt}
\noindent
\begin{tcolorbox}[
    blanker,
    sidebyside,
    sidebyside align=top,
    lefthand width=0.55\textwidth,
    top=4pt, bottom=4pt, left=0pt, right=0pt,
    middle=16pt,
    underlay={
        \draw[dashed, black!80, line width=0.8pt, dash pattern=on 4pt off 3pt]
            ([xshift=0.55\textwidth+8pt]frame.north west) --
            ([xshift=0.55\textwidth+8pt]frame.south west);
    }
]
\textbf{Q.92} \quad {\footnotesize\textsc{[MSQ $\cdot$ 1 Mark(s) $\cdot$ No Neg]}} \hfill {\footnotesize\textbf{GATE 2024 CS1} \quad \tiny\texttt{2024\_CS1\_CS\_Q31}}
\vspace{2pt}\hrule height 0.3pt\vspace{4pt}

Which of the following fields is/are modified in the IP header of a packet going out of a network address translation (NAT) device from an internal network to an external network?

\par\vspace{6pt}\noindent
\begin{tabular}{@{}p{0.48\linewidth}@{\hspace{0.04\linewidth}}p{0.48\linewidth}@{}}
$\square$\ \textbf{(A)} Source IP & $\square$\ \textbf{(B)} Destination IP \\
\addlinespace[3pt]
$\square$\ \textbf{(C)} Header Checksum & $\square$\ \textbf{(D)} Total Length
\end{tabular}

\tcblower

\textbf{Official Key:} \textbf{\large A;C} \hfill \textit{\small [Easy]}
\vspace{2pt}\hrule height 0.3pt\vspace{3pt}
\textbf{Topic:} Computer Networks \\
\textbf{Subtopic:} \textit{Network Address Translation (NAT)}
\vspace{2pt}\hrule height 0.3pt\vspace{3pt}
\textbf{Guide / Hint:} \\
{\small When a packet leaves a private subnet through an outbound NAT gateway, the router rewrites the private Source IP address to its public external IP address (A is modified). Because the IP header fields have changed, the IPv4 Header Checksum must be recalculated and replaced (C is modified). Destination IP and Total Length are unmodified.}

\end{tcolorbox}
\vspace{6pt}\hrule height 0.35pt\vspace{6pt}
\noindent
\begin{tcolorbox}[
    blanker,
    sidebyside,
    sidebyside align=top,
    lefthand width=0.55\textwidth,
    top=4pt, bottom=4pt, left=0pt, right=0pt,
    middle=16pt,
    underlay={
        \draw[dashed, black!80, line width=0.8pt, dash pattern=on 4pt off 3pt]
            ([xshift=0.55\textwidth+8pt]frame.north west) --
            ([xshift=0.55\textwidth+8pt]frame.south west);
    }
]
\textbf{Q.93} \quad {\footnotesize\textsc{[MSQ $\cdot$ 1 Mark(s) $\cdot$ No Neg]}} \hfill {\footnotesize\textbf{GATE 2024 CS2} \quad \tiny\texttt{2024\_CS2\_CS\_Q23}}
\vspace{2pt}\hrule height 0.3pt\vspace{4pt}

Node X has a TCP connection open to node Y. The packets from X to Y go through an intermediate IP router R. Ethernet switch S is the first switch on the network path between X and R. Consider a packet sent from X to Y over this connection.
Which of the following statements is/are TRUE about the destination addresses?

\par\vspace{6pt}\noindent
\begin{tabular}{@{}p{0.98\linewidth}@{}}
$\square$\ \textbf{(A)} The destination IP address is the IP address of R \\
\addlinespace[3pt]
$\square$\ \textbf{(B)} The destination IP address is the IP address of Y \\
\addlinespace[3pt]
$\square$\ \textbf{(C)} The destination MAC address is the MAC address of S \\
\addlinespace[3pt]
$\square$\ \textbf{(D)} The destination MAC address is the MAC address of Y
\end{tabular}

\tcblower

\textbf{Official Key:} \textbf{\large B} \hfill \textit{\small [Easy]}
\vspace{2pt}\hrule height 0.3pt\vspace{3pt}
\textbf{Topic:} Computer Networks \\
\textbf{Subtopic:} \textit{Concept of layering: OSI and TCP/IP Protocol Stacks}
\vspace{2pt}\hrule height 0.3pt\vspace{3pt}
\textbf{Guide / Hint:} \\
{\small In an end-to-end IP network, the destination IP address remains invariant throughout transit and is always the IP address of final host Y. The destination MAC address changes hop-by-hop and is set to the MAC address of the next-hop router R (not switch S, which operates transparently at Layer 2).}

\end{tcolorbox}
\vspace{6pt}\hrule height 0.35pt\vspace{6pt}
\noindent
\begin{tcolorbox}[
    blanker,
    sidebyside,
    sidebyside align=top,
    lefthand width=0.55\textwidth,
    top=4pt, bottom=4pt, left=0pt, right=0pt,
    middle=16pt,
    underlay={
        \draw[dashed, black!80, line width=0.8pt, dash pattern=on 4pt off 3pt]
            ([xshift=0.55\textwidth+8pt]frame.north west) --
            ([xshift=0.55\textwidth+8pt]frame.south west);
    }
]
\textbf{Q.94} \quad {\footnotesize\textsc{[MSQ $\cdot$ 1 Mark(s) $\cdot$ No Neg]}} \hfill {\footnotesize\textbf{GATE 2024 CS2} \quad \tiny\texttt{2024\_CS2\_CS\_Q28}}
\vspace{2pt}\hrule height 0.3pt\vspace{4pt}

Which of the following statements about IPv4 fragmentation is/are TRUE?

\par\vspace{6pt}\noindent
\begin{tabular}{@{}p{0.98\linewidth}@{}}
$\square$\ \textbf{(A)} The fragmentation of an IP datagram is performed only at the source of the datagram \\
\addlinespace[3pt]
$\square$\ \textbf{(B)} The fragmentation of an IP datagram is performed at any IP router which finds that the size of the datagram is greater than the MTU of the outgoing link \\
\addlinespace[3pt]
$\square$\ \textbf{(C)} The reassembly of fragments is performed only at the destination of the datagram \\
\addlinespace[3pt]
$\square$\ \textbf{(D)} The reassembly of fragments is performed at all intermediate routers along the path from the source to destination
\end{tabular}

\tcblower

\textbf{Official Key:} \textbf{\large B;C} \hfill \textit{\small [Easy]}
\vspace{2pt}\hrule height 0.3pt\vspace{3pt}
\textbf{Topic:} Computer Networks \\
\textbf{Subtopic:} \textit{Fragmentation and IP addressing, IPv4, CIDR notation}
\vspace{2pt}\hrule height 0.3pt\vspace{3pt}
\textbf{Guide / Hint:} \\
{\small In IPv4, any intermediate router encountering a packet larger than the outgoing link MTU performs fragmentation (B). In contrast, fragment reassembly is performed strictly at the final destination host (C) to prevent re-fragmentation overhead and handle out-of-order routing.}

\end{tcolorbox}
\vspace{6pt}\hrule height 0.35pt\vspace{6pt}
\noindent
\begin{tcolorbox}[
    blanker,
    sidebyside,
    sidebyside align=top,
    lefthand width=0.55\textwidth,
    top=4pt, bottom=4pt, left=0pt, right=0pt,
    middle=16pt,
    underlay={
        \draw[dashed, black!80, line width=0.8pt, dash pattern=on 4pt off 3pt]
            ([xshift=0.55\textwidth+8pt]frame.north west) --
            ([xshift=0.55\textwidth+8pt]frame.south west);
    }
]
\textbf{Q.95} \quad {\footnotesize\textsc{[MSQ $\cdot$ 1 Mark(s) $\cdot$ No Neg]}} \hfill {\footnotesize\textbf{GATE 2024 CS2} \quad \tiny\texttt{2024\_CS2\_CS\_Q32}}
\vspace{2pt}\hrule height 0.3pt\vspace{4pt}

Which of the following fields of an IP header is/are always modified by any router before it forwards the IP packet?

\par\vspace{6pt}\noindent
\begin{tabular}{@{}p{0.48\linewidth}@{\hspace{0.04\linewidth}}p{0.48\linewidth}@{}}
$\square$\ \textbf{(A)} Source IP Address & $\square$\ \textbf{(B)} Protocol \\
\addlinespace[3pt]
$\square$\ \textbf{(C)} Time to Live (TTL) & $\square$\ \textbf{(D)} Header Checksum
\end{tabular}

\tcblower

\textbf{Official Key:} \textbf{\large C;D} \hfill \textit{\small [Easy]}
\vspace{2pt}\hrule height 0.3pt\vspace{3pt}
\textbf{Topic:} Computer Networks \\
\textbf{Subtopic:} \textit{Fragmentation and IP addressing, IPv4, CIDR notation}
\vspace{2pt}\hrule height 0.3pt\vspace{3pt}
\textbf{Guide / Hint:} \\
{\small Whenever an IP packet is processed by an intermediate router, its Time to Live (TTL) field is decremented by at least 1. Because the header content changes, the router must recompute and update the Header Checksum field. Source and Protocol remain unchanged.}

\end{tcolorbox}
\vspace{6pt}\hrule height 0.35pt\vspace{6pt}
\noindent
\begin{tcolorbox}[
    blanker,
    sidebyside,
    sidebyside align=top,
    lefthand width=0.55\textwidth,
    top=4pt, bottom=4pt, left=0pt, right=0pt,
    middle=16pt,
    underlay={
        \draw[dashed, black!80, line width=0.8pt, dash pattern=on 4pt off 3pt]
            ([xshift=0.55\textwidth+8pt]frame.north west) --
            ([xshift=0.55\textwidth+8pt]frame.south west);
    }
]
\textbf{Q.96} \quad {\footnotesize\textsc{[MSQ $\cdot$ 1 Mark(s) $\cdot$ No Neg]}} \hfill {\footnotesize\textbf{GATE 2025 CS1} \quad \tiny\texttt{2025\_CS1\_CS\_Q22}}
\vspace{2pt}\hrule height 0.3pt\vspace{4pt}

Consider the 3-way handshaking protocol for TCP connection establishment. Let the three packets exchanged during the connection establishment be denoted as P1, P2, and P3, in order. Which of the following option(s) is/are TRUE with respect to TCP header flags that are set in the packets?

\par\vspace{6pt}\noindent
\begin{tabular}{@{}p{0.48\linewidth}@{\hspace{0.04\linewidth}}p{0.48\linewidth}@{}}
$\square$\ \textbf{(A)} P3: SYN = 1, ACK = 1 & $\square$\ \textbf{(B)} P2: SYN = 1, ACK = 1 \\
\addlinespace[3pt]
$\square$\ \textbf{(C)} P2: SYN = 0, ACK = 1 & $\square$\ \textbf{(D)} P1: SYN = 1
\end{tabular}

\tcblower

\textbf{Official Key:} \textbf{\large B;D} \hfill \textit{\small [Easy]}
\vspace{2pt}\hrule height 0.3pt\vspace{3pt}
\textbf{Topic:} Computer Networks \\
\textbf{Subtopic:} \textit{Transport layer: flow control and congestion control, UDP, TCP, sockets}
\vspace{2pt}\hrule height 0.3pt\vspace{3pt}
\textbf{Guide / Hint:} \\
{\small In TCP 3-way handshaking: Packet 1 (Client to Server) sends a SYN segment with flag $[\text{SYN}=1, \text{ACK}=0]$; Packet 2 (Server to Client) responds with $[\text{SYN}=1, \text{ACK}=1]$; Packet 3 (Client to Server) completes setup with $[\text{SYN}=0, \text{ACK}=1]$. Hence, P1 has SYN = 1 and P2 has SYN = 1, ACK = 1.}

\end{tcolorbox}
\vspace{6pt}\hrule height 0.35pt\vspace{6pt}
\end{document}